% How the radiation field sets the molecular content in MOCASSIN-PDR,
% branch h2-model (commit 8fc1512, 2026-10-15).
%
% Build:  pdflatex main && bibtex main && pdflatex main && pdflatex main
% Only standard LaTeX packages are used (no TikZ, siunitx or enumitem), so
% the document compiles on a minimal TeX Live installation.

\documentclass[11pt,a4paper]{article}

\usepackage[utf8]{inputenc}
\usepackage[T1]{fontenc}
\usepackage{lmodern}
\usepackage{textcomp}
\usepackage{amsmath,amssymb,mathtools,bm}
\usepackage[margin=2.3cm]{geometry}
\usepackage{microtype}
\usepackage{booktabs,longtable,array}
\usepackage{caption}
\usepackage{float}
\usepackage{xcolor}
\usepackage{fancyhdr}
\usepackage[round,sort]{natbib}
\usepackage[colorlinks=true,linkcolor=blue!50!black,citecolor=green!35!black,urlcolor=blue!60!black]{hyperref}
\usepackage{url}

% ---------------------------------------------------------------- macros
% Code names are typeset in typewriter and may break after an underscore,
% a colon or a slash (the latter two are marked in the text with \allowbreak).
\newcommand{\brkus}{\textunderscore\allowbreak}
\DeclareRobustCommand{\src}[1]{\texttt{\let\_\brkus #1}}   % file:line, with \_ escaped
\DeclareRobustCommand{\rtn}[1]{\texttt{\let\_\brkus #1}}   % routine or variable name
\DeclareRobustCommand{\kw}[1]{\texttt{\let\_\brkus #1}}    % input keyword or compile flag
\newcommand{\dlog}[1]{\textsf{#1}}         % DEVLOG entry (D-nnn, H-nnn, L-nn)
\newcommand{\Htwo}{H\ensuremath{_2}}
\newcommand{\Hp}{H\ensuremath{^+}}
\newcommand{\Cp}{C\ensuremath{^+}}
\newcommand{\unit}[1]{\,\mathrm{#1}}
\newcommand{\chiD}{\chi_{\mathrm{dir}}}
\newcommand{\chiF}{\chi_{\mathrm{dif}}}
\newcommand{\Jste}{J^{\mathrm{ste}}}
\newcommand{\Jdif}{J^{\mathrm{dif}}}
\newcommand{\uD}{u_{\mathrm{D}}}
\newcommand{\RFUV}{R_{\mathrm{FUV}}}
\newcommand{\AV}{A_{V}}
\newcommand{\dd}{\mathrm{d}}
\newcommand{\note}[1]{\par\smallskip\noindent\fbox{\parbox{0.96\linewidth}{\small #1}}\par\smallskip}

\setlength{\parskip}{0.35em}
\setlength{\emergencystretch}{3em}
\setlength{\parindent}{0pt}
\renewcommand{\arraystretch}{1.15}
\newcolumntype{L}[1]{>{\raggedright\arraybackslash}p{#1}}

\pagestyle{fancy}
\fancyhf{}
\lhead{\small MOCASSIN-PDR, branch \texttt{h2-model}}
\rhead{\small Radiation field $\rightarrow$ molecular content}
\cfoot{\thepage}

\hypersetup{pdftitle={How the radiation field sets the molecular content in MOCASSIN-PDR (h2-model branch)},
            pdfauthor={Compiled from the MOCASSIN-PDR source tree}}

\title{How the radiation field sets the molecular content\\ in MOCASSIN-PDR\\[6pt]
\large The \texttt{h2-model} branch: physics, prescriptions, numerical strategy\\ and code locations}
\author{Compiled from the source tree of MOCASSIN-PDR\\
\small branch \texttt{h2-model}, commit \texttt{8fc1512} (2026-10-15),
worktree \texttt{\textasciitilde/software/mocassin\_h2}\\
\small (document prepared with Claude Code from a reading of the source; no run was made for it)}
\date{16 October 2026}

\begin{document}
\maketitle
\thispagestyle{empty}

\begin{abstract}
\noindent
This document describes, process by process, how the radiation field computed by
MOCASSIN's three-dimensional Monte Carlo (MC) photoionisation and dust radiative
transfer determines the molecular content of a model in the \texttt{h2-model} branch
of MOCASSIN-PDR, in which that transfer is coupled to the 33-species REDUCED
steady-state chemistry of 3D-PDR and to a 302-level model of the ground electronic
state of \Htwo{} whose Solomon photodissociation rate replaces the network's
single-line prescription. The photorates are evaluated on each cell's own MC field:
its never-scattered part arrives along the straight lines from the sources, its
scattered part is distributed over scattering sites by the formal solution of the
transfer equation, and the dust attenuation, already in the MC field, is not applied
again. The absorption in the \Htwo{} and CO lines, which the MC does not transport,
is applied through column densities integrated along those paths, source-ward for the
direct beam and through each site for the diffuse light; the level model shields each
of its 32\,453 lines of the Lyman, Werner, $B'$--$X$ and $D$--$X$ systems with Voigt
profiles averaged through the cell's own column, with line overlap, trapping and
formation pumping. For every step the document gives the physics, the prescription
with its reference, the numerical strategy and the file and line of the
implementation, in the order in which the code executes them: the MC estimator of the
field and its direct and diffuse parts; the dust quantities and the ray geometry; the
\Htwo{} level model; the photoreaction rates; the steady-state chemical solve; and the
feedback of the molecules on MOCASSIN's opacity, ionisation and thermal balance.
Measured numbers from the project's logs size the main approximations: on the
H\,\textsc{ii}-region fixture the scattered field carries 75--100\% of the FUV field
beyond $\AV\approx0.5$; the \Htwo{} lines and the H/\Htwo{} transition are converged
at the default packet count, while the carbon transition and CO are not; and the
damping width adopted for the \Htwo{} lines moves the transition by 0.1--0.2\,mag.
The physical assumptions are collected in Section~\ref{sec:assumptions} and the known
limitations, including several found while preparing this document, in
Section~\ref{sec:limitations}. Appendix~\ref{app:index} is an index of code locations.
\end{abstract}

\tableofcontents
\clearpage

% ======================================================================
\section{Scope, conventions and overview}
\label{sec:scope}
% ======================================================================

\subsection{What is described}

MOCASSIN-PDR couples three pieces of physics:
\begin{enumerate}
\item \textbf{MOCASSIN} \citep{Ercolano2003,Ercolano2005,Ercolano2008}: a
  three-dimensional Monte Carlo (MC) code for photoionisation and dust
  radiative transfer on a Cartesian grid of homogeneous cells. It owns the
  radiation field at every frequency, the dust (optical properties,
  temperatures, charge), the ionisation state of every element, the electron
  density and the gas temperature.
\item \textbf{3D-PDR} \citep{Bisbas2012}, vendored under
  \src{source/\allowbreak{}pdr/\allowbreak{}} (upstream commit \texttt{fabcdc6}, imported 2026-08-28;
  \src{source/\allowbreak{}pdr/\allowbreak{}UPSTREAM.md}): the REDUCED 33-species chemical network solved
  to steady state, its photoreaction prescriptions, its \Htwo{} formation rate
  and its escape-probability level-population solver for the coolants. 3D-PDR's
  own treatment of the radiation field (an external field attenuated along
  rays to the cloud boundary) and its own thermal balance are not used.
  MOCASSIN drives the vendored routines through the bridge module
  \src{source/\allowbreak{}pdr\_mod.f90}.
\item \textbf{The \Htwo{} level model} of the \texttt{h2-model} branch
  (\src{source/\allowbreak{}h2\_mod.f90}, driven by \rtn{h2LevelModel} in
  \src{pdr\_mod.f90}): the 302 rovibrational levels of $\mathrm{H_2}\,X$ in
  statistical equilibrium under collisions, quadrupole decay, ultraviolet
  pumping through the excited electronic states, Solomon dissociation and
  formation on grains. It is a re-implementation from the data and the
  conventions of Cloudy's large \Htwo{} model \citep{Shaw2005,Ferland2017,Chatzikos2023}.
  Its Solomon rate is the \Htwo{} photodissociation rate the chemistry uses.
\end{enumerate}
By \emph{molecular content} we mean the abundances of the 33 species in every
cell and, for \Htwo, the populations of its 302 levels. The document follows
the path by which the MC radiation field determines them, and the feedback of
the molecules on the field, the ionisation and the temperatures.

The text follows the source tree at \texttt{\textasciitilde/software/mocassin\_h2}
(git worktree, branch \texttt{h2-model}, HEAD
\texttt{8fc15120dea2}, 2026-10-15). Every \src{file:\allowbreak{}line} reference is to that
commit; files are under \src{source/\allowbreak{}} unless a path says otherwise. The
branch has moved on since (its tip was \texttt{78a74bd} on 2026-10-16). The only
source change on the path described here is \dlog{D-414} (commit \texttt{df9a626},
merged into \texttt{h2-model} as \texttt{901f1e8}), which makes the mean grain
temperature a mean over the size bins that hold grains
(Section~\ref{sec:Td}); it changes nothing where every grain species fills the
size grid, as in all the fixtures. \dlog{D-417} registers a temperature per
grain material for \Htwo{} formation for future development, without a code
change. DEVLOG
entries are cited as \dlog{D-nnn} (\src{DEVLOG.md}), branch entries as
\dlog{H-nnn} (\src{DEVLOG-H2.md}) and the limitations of
\src{doc/\allowbreak{}pdr-status.md} as \dlog{L-nn}. Measured numbers quoted without a
log entry are read from the run logs of the two production fixtures:
\texttt{hiipdr}, an H\,\textsc{ii} region and PDR around a 40\,kK star in gas of
$n_{\mathrm H}=100$\,cm$^{-3}$ with silicate grains on $31^3$ cells
(\src{runs/\allowbreak{}d405-hii31-d0/\allowbreak{}run.log}), and \texttt{pn150k}, a 150\,kK planetary
nebula with $n_{\mathrm H}=10^3$\,cm$^{-3}$ and graphite grains on $41^3$ cells
(\src{runs/\allowbreak{}d405-pn41-d0/\allowbreak{}run.log}). Both are uniform spheres around one
source, computed in one octant (\kw{symmetricXYZ}).

\paragraph{Build flavour.} The description is of the ray-traced build,
\begin{center}\verb|make PDR=1 PDRGEOM=3D mocassin|\end{center}
which compiles with the flags
\begin{center}\small\ttfamily
-DPDR -DRAYTHEIA -DRAYTHEIA\_MO\\
-DTHERMALBALANCE -DREDUCED -DH2FORM\_CT02 -DCHEMSTEADY -DCHEMRETRY\\
-DFORCECONVERGENCE -DNGACCEL -DNGCYCLE -DNGRELAX -DILLINOIS -DREBRACKET -DSOBOLEV
\end{center}
(\src{Makefile:\allowbreak{}82--89}): the first line selects the ray-traced
geometry, the others are fixed for every PDR build. Of the fixed flags,
\kw{-DREDUCED}, \kw{-DH2FORM\_CT02}, \kw{-DCHEMSTEADY}, \kw{-DCHEMRETRY} and
\kw{-DSOBOLEV} change what the coupled code computes (the network, the \Htwo{}
formation rate, the steady-state integration and the escape-probability
form); \kw{-DTHERMALBALANCE}, \kw{-DFORCECONVERGENCE}, \kw{-DNGACCEL},
\kw{-DNGCYCLE}, \kw{-DNGRELAX}, \kw{-DILLINOIS} and \kw{-DREBRACKET} act only in
3D-PDR's own iteration and temperature routines, which the bridge does not
call. Two other PDR flavours exist and are not
described except where noted: the 1D-slab build (\kw{-DONEDIMENSIONAL}),
whose chemistry still receives 3D-PDR's prescription
$\chi_0\,e^{-\gamma \AV}$ on twelve rays and in which the level model is a
stub that leaves the network's \Htwo{} rate in place
(\src{pdr\_mod.f90:\allowbreak{}6015--6027}); and the warm start
(\verb|make PDR=1 mocassinWarm|). The level model also skips plane-parallel
illumination and nested grids (\src{pdr\_mod.f90:\allowbreak{}5256--5259}), so the
description applies to one grid illuminated by point sources.

\subsection{Notation and units}

Table~\ref{tab:notation} lists the symbols used in more than one section.
cgs units are used throughout; frequencies are in rydbergs in the code
(1\,Ryd $=13.6057$\,eV, $\nu_{1\,\mathrm{Ryd}}=3.28984\times10^{15}$\,Hz).
$J_\nu$ is the angle-averaged specific intensity (the mean intensity), so that
$u_\nu=4\pi J_\nu/c$ is the spectral energy density.

\begin{longtable}{@{}L{0.22\linewidth}L{0.73\linewidth}@{}}
\caption{Notation and units.}\label{tab:notation}\\
\toprule
Symbol & Meaning \\
\midrule
\endfirsthead
\toprule
Symbol & Meaning \\
\midrule
\endhead
$n_{\mathrm H}$ & density of hydrogen nuclei, \rtn{grid\%Hden} [cm$^{-3}$] \\
$x_i = n_i/n_{\mathrm H}$ & fractional abundance of species $i$, \rtn{grid\%pdrAbun} [dimensionless] \\
$J_\nu$ & mean intensity [erg\,s$^{-1}$\,cm$^{-2}$\,Hz$^{-1}$\,sr$^{-1}$] \\
$k$, $\nu_k$, $\Delta\nu_k$ & index, tabulated frequency and width of a bin of the MC frequency grid
            (\rtn{nuArray}, \rtn{widFlx}) \\
$\Jste_k$ & MC estimator of the \emph{total} field in bin $k$,
            $4\pi\int_{\mathrm{bin}\,k} J_\nu\,\dd\nu$ [erg\,s$^{-1}$\,cm$^{-2}$], \rtn{grid\%Jste} \\
$\Jdif_k$ & the part of $\Jste_k$ carried by scattered or re-emitted packets, \rtn{grid\%Jdif} [same units] \\
$\uD$ & $8.94\times10^{-14}$\,erg\,cm$^{-3}$, the 6--13.6\,eV energy density of the
        \citet{Draine1978} interstellar field (\src{constants\_mod.f90:\allowbreak{}29});
        $c\,\uD=2.68\times10^{-3}$\,erg\,s$^{-1}$\,cm$^{-2}$ \\
$\chi$ & band-integrated (6--13.6\,eV) field in Draine units, a dimensionless energy density,
         $\chi = \sum_{k\in\mathrm{FUV}} \Jste_k/(c\,\uD)$; $\chi=1$ is $G_0=1.69$ in
         Habing units of the same band (Section~\ref{sec:bands}) \\
$\chi_0$ & geometrically diluted, unattenuated field of the sources, same units (Equation~\ref{eq:chi0}) \\
$\chiD,\ \chiF$ & direct (never scattered or re-emitted) and diffuse parts of $\chi$ \\
$s_{\mathrm D}$ & share of the Draine field's 6--13.6\,eV energy in the Lyman--Werner (LW) window 11.2--13.6\,eV \\
$f_{\mathrm{dir}},\ f_{\mathrm{dif}}$ & LW window factors of the direct and diffuse parts (Section~\ref{sec:window}) \\
$\eta_{lu}$ & photon occupation number at the line $l\to u$, $J_\nu c^2/(2h\nu^3)$ [dimensionless] \\
$\kappa^{\mathrm{abs}}_\nu,\ \kappa^{\mathrm{sca}}_\nu$ & dust absorption and scattering coefficients
         [cm$^{-1}$], \rtn{absOpac}, \rtn{scaOpac}; $\kappa_V$ their sum at 5500\,\AA \\
$\AV$ & visual extinction [mag]; $1.0857 = 2.5\log_{10}e$ mag per unit optical depth (\rtn{magPerTau}) \\
$\RFUV$ & ratio of the band-mean FUV dust extinction coefficient to $\kappa_V$, per cell (Section~\ref{sec:dust}) \\
$N_X$ & column density of species $X$ [cm$^{-2}$]; $N_l$ the column of \Htwo{} in level $l$ \\
$b$ & Doppler parameter, $b^2 = 2kT/m_{\mathrm{H_2}} + v_{\mathrm{turb}}^2$ [cm\,s$^{-1}$] \\
$\tau_0,\ a$ & line-centre optical depth and Voigt damping parameter of an \Htwo{} line \\
$S(\tau_0,a)$, $W(\tau_0,a)$ & single-line self-shielding function (profile-weighted transmission of a continuum in a Voigt line) and its integral over $\tau_0$ (curve of growth) \\
$k_{\mathrm{Sol}}$ & Solomon photodissociation rate per \Htwo{} molecule [s$^{-1}$], \rtn{grid\%h2Sol} \\
$R_{\mathrm{f}}$ & \Htwo{} grain formation coefficient [cm$^3$\,s$^{-1}$]; formation rate per volume $R_{\mathrm f}\,n_{\mathrm H}\,n(\mathrm H)$ [cm$^{-3}$\,s$^{-1}$] \\
$\zeta$ & cosmic-ray ionisation rate, \kw{pdrZeta} [s$^{-1}$] \\
$T_e$, $T_d$ & gas (electron) temperature and mean grain temperature [K]; $T_{\mathrm{CMB}}=2.725$\,K \\
\bottomrule
\end{longtable}

\subsection{Who owns what}
\label{sec:ownership}

The coupling is designed so that every quantity has a single writer
(\src{CLAUDE.md}, ``one reality per quantity''). Table~\ref{tab:owner}
lists the quantities that lie on the path from the field to the molecules.

\begin{longtable}{@{}L{0.26\linewidth}L{0.16\linewidth}L{0.25\linewidth}L{0.24\linewidth}@{}}
\caption{Ownership of the quantities on the path from the radiation field to the molecular content.}\label{tab:owner}\\
\toprule
Quantity & Owner & Written in & Read by \\
\midrule
\endfirsthead
\toprule
Quantity & Owner & Written in & Read by \\
\midrule
\endhead
$J_\nu$ in every bin (total and diffuse) & MOCASSIN MC & \rtn{pathSegment}, \src{photon\_mod.f90:\allowbreak{}1900--1924, 2236--2250} & everything below \\
dust opacity, albedo, $T_d$ & MOCASSIN & \src{dust\_mod.f90}, \rtn{getDustT} (\src{update\_mod.f90}) & $\RFUV$, \Htwo{} formation, level model \\
$T_e$ & MOCASSIN \rtn{thermBalance} & \src{update\_mod.f90:\allowbreak{}1734} & chemistry, level model \\
ionisation stages of H, He, C, O, Mg (and all other elements), $n_e$ & MOCASSIN \rtn{ionBalance} & \src{update\_mod.f90:\allowbreak{}2527} & imposed on the chemistry \\
atom photoionisation and recombination rates of H, He, C, O, Mg; \Htwo{}$\to$\Htwo$^+$ & MOCASSIN & \src{update\_mod.f90:\allowbreak{}2840--2878} & chemistry (replace network channels) \\
$\AV$ and the source-ward columns $N(\mathrm{H_2}), N(\mathrm C), N(\mathrm{CO})$ & bridge geometry & \rtn{computeAVradial}, \src{pdr\_mod.f90:\allowbreak{}1859} & chemistry \\
scattering sites and their two-leg columns & bridge geometry & \rtn{traceDiffuseSites}, \src{pdr\_mod.f90:\allowbreak{}2136} & chemistry, level model \\
\Htwo{} level populations, $k_{\mathrm{Sol}}$, \Htwo{} heating & level model & \rtn{h2LevelModel}, \src{pdr\_mod.f90:\allowbreak{}5201} & chemistry, thermal balance \\
abundances of the 33 species & chemistry & \rtn{pdrStep}, \src{pdr\_mod.f90:\allowbreak{}3180} & MOCASSIN, level model, columns \\
fraction of each element not in molecules (\rtn{elemAvail}) & chemistry & \rtn{setElemAvail}, \src{pdr\_mod.f90:\allowbreak{}5120} & \rtn{ionBalance} \\
$n(\mathrm{H_2})$ for the gas opacity (\rtn{H2den}) & chemistry & \src{pdr\_mod.f90:\allowbreak{}4573} & \rtn{addOpacity} \\
cooling by CO & chemistry (3D-PDR level populations, tabulated in $T$) & \rtn{pdrStep} & \rtn{thermBalance} \\
cooling by \Cp, C, O & MOCASSIN's own atoms, with the chemistry's neutral colliders and escape probabilities (\dlog{D-363}) & \src{emission\_mod.f90:\allowbreak{}1861--1887} & \rtn{thermBalance}, line fluxes \\
heating terms MOCASSIN lacks (\rtn{pdrHeatXtra}) & chemistry and level model & \src{pdr\_mod.f90:\allowbreak{}4499--4546} & \rtn{thermBalance} \\
\bottomrule
\end{longtable}

\subsection{Overview of the chain}

Figure~\ref{fig:chain} shows the chain inside one MC iteration. The field
reaches the molecules through five channels:
(i) directly, through the photodissociation and photoionisation rates, which
read the local MC field (Sections~\ref{sec:field} and \ref{sec:photorates});
(ii) through the gas temperature, which sets the rate coefficients, the
collisional excitation of \Htwo{} and, through the thermal speed and the
sticking of H atoms, the \Htwo{} formation rate;
(iii) through the dust temperature, which sets the efficiency of \Htwo{}
formation on grains (Section~\ref{sec:Td});
(iv) through MOCASSIN's ionisation balance, which fixes the \Cp/C ratio and
the electron density the chemistry works with; and
(v) through the molecules themselves, whose columns, integrated along the
paths the photons took, shield the \Htwo{} and CO lines at the next
iteration (Section~\ref{sec:geometry}).
The molecules act back on the transfer only through the \Htwo{} photoionisation
continuum and, indirectly, through the ionisation and energy balances
(Section~\ref{sec:feedback}); no molecular line opacity enters the MC transfer.

\begin{figure}[tbp]
\centering
\small
\fbox{\parbox{0.86\linewidth}{\centering
\textbf{1. MC transfer} (\rtn{energyPacketDriver}): $\Jste_k$, $\Jdif_k$ in every cell and bin;
dust temperatures; packets flagged as never scattered (direct) or scattered/re-emitted (diffuse)}}\\[2pt]
$\Big\downarrow$\\[2pt]
\fbox{\parbox{0.86\linewidth}{\centering
\textbf{2. Dust and geometry} (\rtn{dustPerCell}, \rtn{computeAVdust}):
$\RFUV$, $\AV/N_{\mathrm H}$, grain area per H;
exact rays to the sources: $\AV$ and $N(\mathrm{H_2})$, $N(\mathrm C)$, $N(\mathrm{CO})$ of iteration $n-1$'s abundances;
snapshot of the scattering emissivity for the diffuse walk; $\chi$, $\chi_0$}}\\[2pt]
$\Big\downarrow$\\[2pt]
\fbox{\parbox{0.86\linewidth}{\centering
\textbf{3. \Htwo{} level model} (\rtn{h2LevelModel}, one sweep): per line
$\eta = \eta_{\mathrm{dir}}\langle S\rangle_{\mathrm{dir}}\mathcal O_{\mathrm{dir}}+\eta_{\mathrm{dif}}\langle S\rangle_{\mathrm{dif}}$;
Voigt self-shielding with level columns, line overlap, trapping;
302-level statistical equilibrium at $T_e^{(n-1)}$ $\Rightarrow$ $k_{\mathrm{Sol}}$, \Htwo{} heating}}\\[2pt]
$\Big\downarrow$\\[2pt]
\fbox{\parbox{0.86\linewidth}{\centering
\textbf{4. MOCASSIN cell update} (\rtn{updateCell}): ionisation and $T_e$ of every cell,
with the chemistry's channels of iteration $n-1$;
stores the photoionisation and recombination rates of H, He, C, O, Mg and the \Htwo{} photoionisation rate}}\\[2pt]
$\Big\downarrow$\\[2pt]
\fbox{\parbox{0.86\linewidth}{\centering
\textbf{5. Chemistry} (\rtn{pdrStep}): rays $=$ direct beam $+$ scattering sites, all at $\AV=0$;
generic photorates $\alpha\chi$; CO with \citet{vanDishoeckBlack1988} shielding and window factor;
\Htwo{} photodissociation $=k_{\mathrm{Sol}}$; MOCASSIN's ionisation imposed;
CVODE to steady state; CO cooling table; heating terms}}\\[2pt]
$\Big\downarrow$\\[2pt]
\fbox{\parbox{0.86\linewidth}{\centering
\textbf{6. Feedback} (\rtn{pdrToMocassin}, \rtn{rescaleIonToAvail}): fraction of each element in molecules
$\rightarrow$ ion stages; $n(\mathrm{H_2})\rightarrow$ gas opacity; coolants and heating $\rightarrow$ next $T_e$;
new abundances $\rightarrow$ next iteration's shielding columns; second ionisation solve (kept at the end)}}
\caption{The chain from the radiation field to the molecular content inside one MC iteration $n$ of the
\texttt{h2-model} branch (3D build), in the order of execution. Inputs marked ``iteration $n-1$'' or
with the superscript $(n-1)$ are one iteration old (Section~\ref{sec:lags}); the lags vanish at a fixed
point of the iteration. The steps and their line numbers are listed in Table~\ref{tab:steps}.}
\label{fig:chain}
\end{figure}

% ======================================================================
\section{The iteration in which the chain runs}
\label{sec:iteration}
% ======================================================================

All of the coupling lives inside MOCASSIN's MC iteration, the recursive
subroutine \rtn{iterateMC} (\src{iteration\_mod.f90:\allowbreak{}56--3345}), which
advances the iteration counter and calls itself until the stop test of
Section~\ref{sec:convergence} passes or the maximum number of iterations is
reached. The PDR part switches on at iteration \kw{pdrStartIter} (default 1,
\src{common\_mod.f90:\allowbreak{}893}); a larger value lets MOCASSIN first converge an
ionisation structure, a gas and a dust temperature without molecules. From
then on every iteration executes the steps of Table~\ref{tab:steps}, in that
order.

\begin{longtable}{@{}L{0.04\linewidth}L{0.25\linewidth}L{0.25\linewidth}L{0.38\linewidth}@{}}
\caption{One MC iteration with the PDR coupling active (3D build). Lines are in \src{iteration\_mod.f90} unless stated.}\label{tab:steps}\\
\toprule
\# & Step & Location & What it does \\
\midrule
\endfirsthead
\toprule
\# & Step & Location & What it does \\
\midrule
\endhead
1 & MC transfer & \rtn{energyPacketDriver}, :611--613 (sources), :655--660 (optional diffuse source) & propagates the energy packets of every source; an absorbed packet is re-emitted on the spot by the local gas or dust and followed until it escapes; accumulates $\Jste$ and $\Jdif$ (Section~\ref{sec:field}) \\
2 & estimator units & :924, :940 & multiplies by $10^{-9}$ (packet luminosity in $10^{36}$\,erg\,s$^{-1}$, cell volume in $10^{45}$\,cm$^3$, path lengths in cm) and divides by 8 for a \kw{symmetricXYZ} octant \\
3 & dust per cell & \rtn{calibratePDRdust} :986 (first time), \rtn{dustPerCell} :993 & $\RFUV$, $\AV/N_{\mathrm H}$, grain area per H, model-wide fallbacks (Section~\ref{sec:dust}) \\
4 & columns and $\AV$ & \rtn{computeAVdust} :995 & \rtn{snapDiffuseSources} (snapshot for the diffuse walk), \rtn{computeAVradial} ($\AV$ and source-ward columns) (Section~\ref{sec:geometry}) \\
5 & field summary & \rtn{mocassinToPDR} :1023 & $\chi$ (local) and $\chi_0$ (geometric, unattenuated) per cell \\
6 & seed (first PDR iteration only) & :1027--1029, :1037 & \rtn{seedMolecules}, then steps 4--5 again on the seeded abundances, then \rtn{pdrToMocassin} (Section~\ref{sec:seed}) \\
7 & \Htwo{} level model & \rtn{h2LevelModel(grid,1)} :1043 & one sweep over all cells on this iteration's field (Section~\ref{sec:h2model}) \\
8 & MOCASSIN cell update & \rtn{updateCell} :1261 & ionisation and thermal balance of every cell; stores the atomic photoionisation and recombination rates the chemistry will use \\
9 & chemistry & \rtn{pdrStep} :1986, cells dealt round-robin to the ranks (:1980--1986), then \verb|MPI_Allreduce| of abundances, heating, cooling and rates & photorates, steady-state solve, heating and cooling for the thermal balance (Sections~\ref{sec:photorates}--\ref{sec:chemistry}) \\
10 & handback & \rtn{pdrToMocassin} :2154 & fraction of each element not in molecules (\rtn{elemAvail}) \\
11 & second ionisation solve & \rtn{ionSecondSolve} called at :2168 (routine :3361--3434); its control pass, before the chemistry, at :1274 & \rtn{ionBalance} solved again on the channels and molecular content this iteration's chemistry left, then put back (\dlog{D-404}); the result is set aside for step 15 \\
12 & damping, hydrogen budget & :2424 (only if \kw{iterDamp}$>0$; default 0), \rtn{rescaleIonToAvail} :2474 & ion stages rescaled to \rtn{elemAvail} so that $x(\mathrm{H^0})+x(\mathrm{H^+})+2x(\mathrm{H_2})=1$ \\
13 & stop test & :3243 & integrated-quantity test (\kw{convTotals}, Section~\ref{sec:convergence}) \\
14 & next iteration & recursive call of \rtn{iterateMC} & the opacity of the next transfer is built from this iteration's ionisation, dust and $n(\mathrm{H_2})$ \\
15 & finish & \rtn{ionKeepSecond} :3249/:3282 (routine :3455--3521), \rtn{h2LevelModel(grid,\allowbreak{}12)} :3254/:3283, outputs & the last second solve replaces the first in the written model; level model swept to convergence on the final model; \src{h2Cells.out}, \src{lineFluxH2.out}, \src{h2Pop.bin} \\
\bottomrule
\end{longtable}

\subsection{What each step reads, and from when}
\label{sec:lags}

The coupling is a fixed-point iteration in which some inputs are one
iteration old. Within iteration $n$:

\begin{itemize}
\item the \textbf{radiation field} used by steps 3--9 is iteration $n$'s,
  transported through the opacity built from iteration $n-1$'s ionisation,
  dust and $n(\mathrm{H_2})$;
\item the \textbf{shielding columns} (steps 4 and 7) are integrated over the
  abundances left by iteration $n-1$'s chemistry. The snapshot is taken
  before any cell of iteration $n$ is solved, so the columns do not depend
  on the rank count or on the order in which the cells are solved;
\item the \textbf{level model} (step 7) runs before \rtn{updateCell}: it sees
  iteration $n-1$'s $T_e$, collider abundances and $n(\mathrm{H_2})$, and
  iteration $n$'s field; the other cells' level populations in its columns
  are those of its previous sweep (a Jacobi iteration);
\item \textbf{MOCASSIN's ionisation balance} (step 8) takes the chemistry's
  channels to and from molecules, and normalises the ion stages to the
  molecular content, of iteration $n-1$ (\dlog{L-47}); step 11 measures what
  the delay costs and, at the end of the run, removes it from the written
  ions;
\item the \textbf{chemistry} (step 9) sees iteration $n$'s $T_e$, $n_e$,
  ionisation and atomic rates (step 8), iteration $n$'s field, the step-7
  Solomon rates and the columns of step 4;
\item the \textbf{coolant line escape} in step 9 integrates the optical depth
  over the level populations of the other cells as stored on the rank; these
  are iteration $n-1$'s except for cells solved earlier in the same sweep on
  the same rank, so this one input does depend weakly on the order of the
  sweep and the rank count;
\item the \textbf{heating terms} that MOCASSIN takes from the chemistry and
  the level model (\rtn{pdrHeatXtra}) and its \textbf{CO cooling table} enter
  step 8 as step 9 of iteration $n-1$ left them, and the heating terms are
  held fixed during the temperature search (Section~\ref{sec:thermal}). The
  cosmic-ray, turbulent and chemical terms were evaluated at $T_e^{(n-1)}$; the
  \Htwo{} terms come from the level model of iteration $n-1$, which ran at
  $T_e^{(n-2)}$.
\end{itemize}

Inside one iteration the chemistry is solved once per cell: every input of
the solve is frozen during the iteration, so repeating it would reproduce the
same answer (\dlog{D-354}, which retired the inner coupling loop after
finding that a column rescale between passes was the only thing that changed
and that it oscillated). The damping across iterations (\kw{iterDamp}) is off
by default (\dlog{D-406}): with it, the written model was a blend of solves
that trailed the undamped one by three to six iterations and, where a cell
has two solutions, wrote states no solve gave (\dlog{L-49}).

At a fixed point of the iteration every lagged input equals its current
value, so the lags determine the path to the solution, not the solution. Two
things qualify this. First, the MC field is redrawn at every iteration, so
the fixed point is reached only in a statistical sense: per-cell quantities
fluctuate about it with the MC noise of Section~\ref{sec:mcnoise}, and
because the photorates are non-linear in the field and the shielding is
non-linear in the columns, the noise can bias cell means (\dlog{L-26},
\dlog{D-408}). Second, the coupled system has cells with two stable states
(the carbon ionisation turn, \dlog{D-405}), in which the lagged inputs decide
which state a cell is in at a given iteration.

% ======================================================================
\section{The Monte Carlo radiation field}
\label{sec:field}
% ======================================================================

\subsection{Transport and the path-length estimator}

MOCASSIN transports the radiation of each source $s$ of luminosity $L_s$ as
$N_s$ energy packets of equal luminosity $\dot\varepsilon_s=L_s/N_s$
(\rtn{deltaE}; \citealt{Ercolano2003}). A packet's frequency bin is drawn
from the source's cumulative emitted-energy distribution
(\rtn{inSpectrumProbDen}; an energy, not a photon-number, cumulative, as
packets of equal energy require). The packet travels in a straight line until
it interacts, the optical depth to the interaction, $-\ln(1-\mathcal R)$ with
$\mathcal R$ uniform in $[0,1)$, being accumulated with the total opacity of
the cells it crosses: the gas continua (bound--free absorption by every ion
present, including C\,\textsc{i} and the other neutral metals in the FUV, the
\Htwo{} photoionisation continuum above 15.4\,eV, free--free) plus dust
absorption and scattering. At an interaction the packet is
\begin{itemize}
\item scattered by dust with probability $\kappa^{\mathrm{sca}}_\nu/\kappa_\nu$, with $\kappa_\nu$ the
  total opacity (\src{photon\_mod.f90:\allowbreak{}2025--2035}): coherently, keeping its
  frequency, into a new direction drawn from the \citet{HenyeyGreenstein1941}
  phase function with the bin's asymmetry parameter $g_\nu$ (\rtn{hg},
  \src{:\allowbreak{}2154--2165}), or isotropically if isotropic scattering is selected;
\item otherwise absorbed, by the dust with probability
  $\kappa^{\mathrm{abs}}_\nu/(\kappa_\nu-\kappa^{\mathrm{sca}}_\nu)$ and by the gas otherwise
  (\dlog{D-233}; \src{:\allowbreak{}2066--2103}), and re-emitted on the spot as a new
  packet of the same energy at a frequency drawn from the local dust (thermal)
  or gas (hydrogen and helium diffuse continua and lines; \dlog{L-51})
  emissivity.
\end{itemize}
Packets carry a weight $w$ (\rtn{enPacket\%weight}) that is unity unless the
optional composite path-length biasing of \citet{Baes2016} is switched on
(\kw{pathStretch}, off by default, \src{common\_mod.f90:\allowbreak{}1099--1103}).

The mean intensity is measured with the path-length estimator of
\citet[eq.~12]{Lucy1999}. A packet of luminosity $\dot\varepsilon_s$ that travels a
length $\ell$ inside a cell of volume $V$ adds $\dot\varepsilon_s\ell/c$ to the
time-averaged radiant energy in the cell, and $u=4\pi J/c$, so the integral of
the mean intensity over bin $k$ is
\begin{equation}
  \Jste_k \equiv 4\pi\int_{\mathrm{bin}\,k} J_\nu\,\dd\nu
          = \frac{1}{V}\sum_{\mathrm{segments}\in k} \dot\varepsilon_s\, w\,\ell
  \qquad [\mathrm{erg\,s^{-1}\,cm^{-2}}],
  \label{eq:estimator}
\end{equation}
the sum running over every path segment of every packet (stellar or not, of
every source) in the cell and in the bin, whether the segment ends at the
cell's face (\src{photon\_mod.f90:\allowbreak{}2236--2250}) or at an interaction
inside it (\src{:\allowbreak{}1900--1924}). $\Jste_k$ is a sum over the bin, \emph{not}
a quantity per unit frequency; the band energy density is
$u_{\mathrm{band}}=\sum_k \Jste_k/c$. The accumulation is in MOCASSIN's internal
units (packet luminosity in $10^{36}$\,erg\,s$^{-1}$, cell volume in
$10^{45}$\,cm$^3$, lengths in cm); the factor $10^{-9}$ that converts to cgs,
and the division by 8 for a \kw{symmetricXYZ} octant (into which every packet
of a single central source is launched), are applied after the transport
(\src{iteration\_mod.f90:\allowbreak{}924}, \src{:\allowbreak{}940}). The default frequency grid
has \kw{nbins} $=600$ bins (\src{set\_input\_mod.f90:\allowbreak{}77}). The estimator is
unbiased and its relative variance in a cell and bin falls as the inverse of
the number of segments recorded there, so it degrades with depth as the
packets thin out (Section~\ref{sec:mcnoise}).

\subsection{Direct and diffuse parts}
\label{sec:split}

A packet leaves a source with \rtn{lgStellar} $=$ \texttt{.true.}; it is
re-initialised as non-stellar whenever it is scattered
(\src{photon\_mod.f90:\allowbreak{}2154}) or re-emitted. In a PDR run every packet adds
to $\Jste$, and the non-stellar ones are mirrored into $\Jdif$
(\rtn{lgJdifLive} is forced true for any PDR run,
\src{set\_input\_mod.f90:\allowbreak{}775--778}; accumulation at
\src{photon\_mod.f90:\allowbreak{}1920--1922} and \src{:\allowbreak{}2247--2249}). Hence
\begin{align}
  \Jste_k &= \text{total field},\nonumber\\
  \Jdif_k &= \text{scattered} + \text{re-emitted},\\
  \Jste_k - \Jdif_k &= \text{the never-scattered source photons}.\nonumber
\end{align}
The last is the direct beam of the sources, averaged over the cell and
attenuated along its path by dust absorption \emph{and} scattering and by the
gas continua; it arrives at the cell along the straight line from each source.
In the 6--13.6\,eV band the diffuse part is dust-scattered starlight of any
order of scattering plus the gas emission that falls in the band (the
two-photon continuum of hydrogen and the line packets not destroyed by dust,
\dlog{D-376}); thermal dust emission contributes nothing there. Two properties
of the split matter for what follows. First, the diffuse part dominates the
PDR: on \texttt{hiipdr} the diffuse fraction $\chiF/\chi$ rises from 0.35 at
$\AV<0.2$ to 0.75 at $\AV=0.5$--0.8, 0.93 at 1.2--1.6 and above 0.99 beyond
$\AV=2.1$ (\dlog{D-344}), so the deep PDR sees almost no direct starlight.
Second, because scattering keeps photons in the cloud, the total field decays
with depth much more slowly than pure extinction: on the \texttt{V1} slab the
transported FUV field fell by $e^{-1.86}$ per magnitude of $\AV$, against
$e^{-3.02}$ for absorption and scattering removing the photons alike
(\src{pdr\_mod.f90:\allowbreak{}1836--1844}, \dlog{D-086}, \dlog{D-087}). (With the debug
switch \rtn{lgDebug} the diffuse packets go into $\Jdif$ alone; every
consumer adds them back in that case.)

\subsection{Bands and the Draine normalisation}
\label{sec:bands}

\begin{itemize}
\item \textbf{FUV band}: $\nu\in[0.44066,\,1.0]$\,Ryd, i.e.\ 6.0--13.6\,eV
  (\src{constants\_mod.f90:\allowbreak{}20--21}). The bins inside it,
  \rtn{FUVlowP}\ldots\rtn{FUVhighP}, run from the first bin with $\nu_k>0.44066$\,Ryd
  to the last with $\nu_k\le1$\,Ryd, located once at start-up
  (\src{continuum\_mod.f90:\allowbreak{}92--94}). On the fixtures the band is 37
  (\texttt{hiipdr}) and 39 (\texttt{pn150k}) of the 600 bins, starting at
  $\nu_k=6.17$ and 6.02\,eV; the bins are $\Delta\nu/\nu\approx2$\% wide. (The
  comment at \src{constants\_mod.f90:\allowbreak{}18--20} calls the band 912--2400\,\AA;
  the value 0.44066\,Ryd is 2068\,\AA, i.e.\ 6.0\,eV, which is the band the
  code uses.)
\item \textbf{Lyman--Werner (LW) window}: $\nu_k\in[0.8232,\,1.0)$\,Ryd,
  11.2--13.6\,eV (\src{constants\_mod.f90:\allowbreak{}25}), about 9 bins: the range in
  which \Htwo{} absorbs in its Lyman and Werner bands and is photodissociated,
  in which CO is photodissociated, and above whose lower edge (11.26\,eV)
  C\,\textsc{i} is photoionised.
\item \textbf{$V$ band}: the bin containing 5500\,\AA{} (\rtn{VbandP},
  \src{constants\_mod.f90:\allowbreak{}32}), used to define $\AV$.
\item \textbf{Draine unit}: $\uD = 8.94\times10^{-14}$\,erg\,cm$^{-3}$
  (\src{constants\_mod.f90:\allowbreak{}29}), the 6--13.6\,eV energy density
  $(4\pi/c)\int E\,F_{\mathrm D}(E)\,\dd E$ (with $E$ converted to erg in the
  integrand) of the interstellar field of \citet{Draine1978}, whose photon
  intensity $F_{\mathrm D}$ is Equation~(\ref{eq:draine}); integrating the polynomial gives
  $8.940\times10^{-14}$, and the code's own quadrature on its bins
  $8.938\times10^{-14}$ (\src{pdr\_mod.f90:\allowbreak{}2280--2286}). In photon units the
  same field carries $4\pi\int F_{\mathrm D}\,\dd E=1.94\times10^{8}$\,photons\,cm$^{-2}$\,s$^{-1}$
  over 6--13.6\,eV and $2.06\times10^{7}$ in the LW window.
\item \textbf{Relation to Habing units.} The \citet{Habing1968} field is
  usually normalised by its 6--13.6\,eV flux,
  $H_0=1.6\times10^{-3}$\,erg\,s$^{-1}$\,cm$^{-2}$ \citep[as in][]{Bisbas2015}, an
  energy density $H_0/c=5.3\times10^{-14}$\,erg\,cm$^{-3}$. In that band
  $\chi=1$ is $G_0=c\,\uD/H_0=1.68$ (1.69 with the value
  $5.29\times10^{-14}$\,erg\,cm$^{-3}$ quoted at
  \src{constants\_mod.f90:\allowbreak{}26--28}); 3D-PDR's heating routine converts with
  1.68 (\src{pdr/\allowbreak{}heatingfunctions.F90:\allowbreak{}57}), in terms this coupling does
  not use (Section~\ref{sec:thermal}). The factor 1.71 compares the two fields
  at 1000\,\AA{} instead: \citet{Habing1968} gives $\lambda u_\lambda=4\times10^{-14}$\,erg\,cm$^{-3}$
  there, Equation~(\ref{eq:draine}) gives $6.84\times10^{-14}$.
  \citet{DraineBertoldi1996} normalise their field parameter (also written
  $\chi$) to Habing's at 1000\,\AA, so that the Draine field has $\chi_{\mathrm{DB}}=1.71$;
  the PDR code comparison adopts $\chi=1.71\,G_0$ \citep[their Table~5]{Rollig2007},
  and TORUS-3DPDR divides its MC field by $1.71H_0$ \citep[their Eq.~12]{Bisbas2015},
  a Draine unit 2.1\% larger than $c\,\uD$. A band of 912--2400\,\AA, as in
  TORUS-3DPDR, holds 13\% more of a Draine-shaped field than 6--13.6\,eV. As in
  the code comparison, $\chi$ here is a $4\pi$ average (an energy density),
  whatever the angular distribution of the field.
\end{itemize}

\subsection{Band-integrated field strengths}

In \rtn{mocassinToPDR} (\src{pdr\_mod.f90:\allowbreak{}2538--2677}):
\begin{align}
  \chi &= \frac{1}{c\,\uD}\sum_{k=\mathrm{FUVlowP}}^{\mathrm{FUVhighP}} \Jste_k ,
  &&\text{(local, attenuated; \rtn{chiFUV}, :2594--2600)} \label{eq:chi}\\
  \chi_0 &= \sum_s \frac{L_{\mathrm{FUV},s}}{4\pi r_s^2\, c\,\uD},
  &&\text{(geometric, unattenuated; \rtn{chiFUV0}, :2604--2616)} \label{eq:chi0}
\end{align}
where $L_{\mathrm{FUV},s}$ is the luminosity of source $s$ in the FUV bins,
taken from the same cumulative distribution the packets are drawn from
(:2558--2562), and $r_s$ the distance of the cell centre from the source.
$\chi_0$ is the energy density of the undimmed beams, $u=F/c$ for a beam of
flux $F$; it is used only by the 1D build and by diagnostics. $\chi$ is passed
to 3D-PDR's heating routine, where only the photoelectric terms read it, and
those are not used (Section~\ref{sec:thermal}); it also enters the seed
(Section~\ref{sec:seed}). The photorates use the split of
Section~\ref{sec:split}, evaluated per cell in \rtn{pdrStep}
(\src{pdr\_mod.f90:\allowbreak{}3401--3417}):
\begin{equation}
  \chiD = \frac{1}{c\,\uD}\sum_{k\in\mathrm{FUV}} \bigl(\Jste_k-\Jdif_k\bigr),
  \qquad
  \chiF = \frac{1}{c\,\uD}\sum_{k\in\mathrm{FUV}} \Jdif_k ,
  \label{eq:chisplit}
\end{equation}
so that $\chiD+\chiF=\chi$; $\chiD$ is clipped at zero against the rounding of
the two single-precision sums.

\subsection{Lyman--Werner window factors}
\label{sec:window}

3D-PDR's photorate coefficients $\alpha$ are rates per unit Draine field,
computed for the Draine spectral shape, and are applied to the 6--13.6\,eV band
strength. For the four reactions driven only in the LW window (\Htwo, HD and
CO photodissociation, C\,\textsc{i} photoionisation) a field whose shape differs
from Draine's in the band drives them at a different rate for the same $\chi$.
The bridge therefore multiplies each such rate by the ratio of the window's
share of the band energy in the MC field to the same share in the Draine field,
separately for the direct and the diffuse part (\dlog{D-349}, \dlog{D-350}):
\begin{equation}
  f_{\mathrm{dir}} = \frac{1}{s_{\mathrm D}}
     \frac{\sum_{k\in\mathrm{LW}}(\Jste_k-\Jdif_k)}{\sum_{k\in\mathrm{FUV}}(\Jste_k-\Jdif_k)},
  \qquad
  f_{\mathrm{dif}} = \frac{1}{s_{\mathrm D}}
     \frac{\sum_{k\in\mathrm{LW}}\Jdif_k}{\sum_{k\in\mathrm{FUV}}\Jdif_k}
  \label{eq:window}
\end{equation}
(\src{pdr\_mod.f90:\allowbreak{}3418--3420}). Here $s_{\mathrm D}$ is the LW share of the
Draine field's 6--13.6\,eV energy, computed once by \rtn{draineWindowShare}
(\src{pdr\_mod.f90:\allowbreak{}2288--2314}) from the photon intensity of
\citet[eq.~11, given there as valid over 5--13.6\,eV]{Draine1978}
\begin{equation}
  F_{\mathrm D}(E) = 1.658\times10^{6}E - 2.152\times10^{5}E^2 + 6.919\times10^{3}E^3
  \quad [\mathrm{photons\,cm^{-2}\,s^{-1}\,sr^{-1}\,eV^{-1}}],\quad E\ \text{in eV},
  \label{eq:draine}
\end{equation}
weighted by $E\,\Delta E_k$ on MOCASSIN's own bins, with the same bin
membership as the MC sums. The continuous band gives $s_{\mathrm D}=0.148$; on
the bins it is 0.155 (\texttt{hiipdr}) and 0.164 (\texttt{pn150k}), printed as
\texttt{[D350]} in the run logs. Using the binned value makes $f=1$ exactly for
an MC field of Draine shape, whatever the bins.
The correction rests on three approximations. (i) It rescales the energy in
the window but not the spectral shape inside it, which matters because the
window is emptied from its blue end first (dust extinction rises towards
912\,\AA{} and C\,\textsc{i} absorbs above 11.26\,eV). (ii) It is a ratio of energy
shares, whereas the rates count photons; the two differ by the ratio of the
mean photon energies in the window, which lie between 11.2 and 13.6\,eV, so
by at most $\sim$20\%, and much less where the shape in the window is close to
Draine's. (iii) One factor serves all the sites of the diffuse field
(Section~\ref{sec:sites}). The size of what it corrects is large: the MC
field's window share falls from 0.299 at $\AV=0.07$ to 0.152, 0.094, 0.041 and
0.008 at $\AV=0.74$, 1.13, 1.70 and 2.50 on \texttt{hiipdr}
(\dlog{D-349}), so $f$ runs from 1.9 near the star to 0.05 deep in the PDR.
The factors are handed to the vendored code as the per-ray array
\rtn{ray\_lw\_factor} (\src{pdr/\allowbreak{}calc\_reac\_rates.F90:\allowbreak{}64--74}; a local
modification listed in \src{pdr/\allowbreak{}UPSTREAM.md}). In the \texttt{h2-model} branch
only the CO rate actually uses them: the \Htwo{} rate is replaced by the level
model, C photoionisation by MOCASSIN's own rate, and the network has no HD
(Section~\ref{sec:photorates}). The other photoreactions keep the Draine shape
(Section~\ref{sec:generic}).

\subsection{Specific intensity and occupation number at each \texorpdfstring{\Htwo}{H2} line}
\label{sec:eta}

The level model needs the field at each of the 32\,453 pumping lines rather
than a band integral. In \rtn{solveCell} (\src{pdr\_mod.f90:\allowbreak{}5670--5695}) each
bin's direct and diffuse mean intensities are
\begin{equation}
  J_{\nu,k}^{\mathrm{dir}} = \frac{\Jste_k-\Jdif_k}{4\pi\,\Delta\nu_k},\qquad
  J_{\nu,k}^{\mathrm{dif}} = \frac{\Jdif_k}{4\pi\,\Delta\nu_k},\qquad
  \Delta\nu_k = \mathrm{widFlx}_k \times 3.28984\times10^{15}\unit{Hz},
\end{equation}
in erg\,s$^{-1}$\,cm$^{-2}$\,Hz$^{-1}$\,sr$^{-1}$, and every line $lu$ of wavenumber
$\tilde\nu_{lu}$ takes the bin that contains it, $\nu_k \le \nu_{lu} < \nu_{k+1}$
(\rtn{tBin}, \src{:\allowbreak{}5310--5320}, the rule MOCASSIN uses to place its own line
packets). The photon occupation numbers are
\begin{equation}
  \eta^{\mathrm{dir}}_{lu} = J^{\mathrm{dir}}_{\nu,k(lu)}\,\frac{c^2}{2h\nu_{lu}^3},\qquad
  \eta^{\mathrm{dif}}_{lu} = J^{\mathrm{dif}}_{\nu,k(lu)}\,\frac{c^2}{2h\nu_{lu}^3},
  \qquad \nu_{lu} = c\,\tilde\nu_{lu}.
  \label{eq:eta}
\end{equation}
Lines outside the frequency grid get $\eta=0$. Lines from vibrationally
excited levels that lie below 11.2\,eV read their own bins as well, inside or
outside the 6--13.6\,eV band. The field at a line is the bin average of the
continuum, 2\% wide: the MC field contains no \Htwo{} (or H\,\textsc{i}
Lyman-series) line opacity, and the line shielding is applied afterwards
through column densities (Section~\ref{sec:h2model}). This separation of
continuum and line attenuation is exact for a line whose absorption does not
change the continuum it sits on. It is not exact for the continuum itself:
photons that \Htwo{} lines absorb remain in the MC field, where they go on
heating the dust, ejecting photoelectrons, photoionising C\,\textsc{i} and being
scattered. The fraction of the band so affected is at most the window's share
and falls to the line cores once the lines saturate; its effect has not been
measured.

% ======================================================================
\section{Dust quantities seen by the chemistry and the level model}
\label{sec:dust}
% ======================================================================

MOCASSIN computes the dust absorption and scattering coefficients
$\kappa^{\mathrm{abs}}_\nu$, $\kappa^{\mathrm{sca}}_\nu$ [cm$^{-1}$] and the
asymmetry parameter $g_\nu$ of every cell from the grain species and the size
distribution of the input \citep{Ercolano2005}: per-grain cross sections from
Mie theory on optical constants \citep{BohrenHuffman1983} or from tabulated
efficiencies, weighted by each species' number abundance and the size
distribution and multiplied by the grain number density $N_d$ that follows
from the dust-to-gas mass ratio (\src{doc/\allowbreak{}dust.md}). It also solves the
radiative-equilibrium temperature of every species and size in every cell from
the local MC field. The bridge derives from these, per cell, the quantities
below. None is an input parameter: they follow the model's own grains
(\src{CLAUDE.md}, ``no hardcoded values''). Table~\ref{tab:dustfix} gives
their values on the two fixtures.

\begin{table}[H]
\centering
\caption{Dust quantities measured on the two production fixtures (run logs,
lines \texttt{[D073]}, \texttt{[D109]}, \texttt{[D195]}; values at the
calibration cell, uniform across each model).}\label{tab:dustfix}
\begin{tabular}{@{}lll@{}}
\toprule
Quantity & \texttt{hiipdr} (silicate) & \texttt{pn150k} (graphite) \\
\midrule
$\AV/N_{\mathrm H}$ [mag\,cm$^2$] & $5.18\times10^{-22}$ & $3.58\times10^{-22}$ \\
$N_{\mathrm H}/\AV$ [cm$^{-2}$\,mag$^{-1}$] & $1.93\times10^{21}$ & $2.79\times10^{21}$ \\
$\RFUV$ & 4.31 & 2.98 \\
FUV albedo, energy-weighted over 6--13.6\,eV & 0.45 & 0.37 \\
FUV asymmetry parameter $g$, scattering-weighted & 0.61 & 0.29 \\
grain area per H, silicate / carbonaceous [cm$^2$] & $1.02\times10^{-21}$ / 0 & 0 / $5.45\times10^{-22}$ \\
\bottomrule
\end{tabular}
\end{table}

\subsection{FUV-to-visual extinction ratio and \texorpdfstring{$\AV$}{AV} per H}

\rtn{dustPerCell} (\src{pdr\_mod.f90:\allowbreak{}1363--1436}) computes
\begin{gather}
  \kappa_V = \kappa^{\mathrm{abs}}_V + \kappa^{\mathrm{sca}}_V,\qquad
  \kappa_{\mathrm{FUV}} = \frac{\sum_{k\in\mathrm{FUV}} w_k\,(\kappa^{\mathrm{abs}}_k+\kappa^{\mathrm{sca}}_k)}{\sum_{k\in\mathrm{FUV}} w_k},
  \nonumber\\
  \RFUV = \frac{\kappa_{\mathrm{FUV}}}{\kappa_V},\qquad
  \frac{\AV}{N_{\mathrm H}} = \frac{1.0857\,\kappa_V}{n_{\mathrm H}}\ [\mathrm{mag\,cm^2}],
  \label{eq:rfuv}
\end{gather}
stored as \rtn{grid\%dustRfuv} and \rtn{grid\%dustAVfac}. The weight $w_k$ is the
fraction of source~1's emitted energy in bin $k$ (the increment of its
cumulative distribution), so that $\kappa_{\mathrm{FUV}}$ is the mean
extinction coefficient of the band-integrated field as the source emits it; a
mean weighted by bin width would over-weight the 912\,\AA{} end of the band,
where the extinction is steepest. The extinction includes scattering, as in
the definition of $\AV$, so $e^{-\RFUV\tau_V}$ is the transmission of the
source's FUV beam with every scattered photon counted as lost.

Two approximations limit what $\RFUV$ represents. It is weighted by the
spectrum of source~1 for every source, and by the spectrum as emitted, not as
it arrives: with depth the field is depleted first at the blue end of the band
(\dlog{D-349}), so the band-mean extinction of the local field is lower than
$\RFUV\kappa_V$ and $\RFUV$ overstates it increasingly with depth. In the 3D
build $\RFUV$ does not enter any photorate. It enters the combination of the
$\AV$ of several sources (Equation~\ref{eq:avradial}, where it cancels for a
single source), the grey extinction of the diffuse-site weights
(Section~\ref{sec:sites}), the fallback field $\chi\,e^{-\RFUV\AV/1.0857}$ of a
cell the MC field did not reach (\src{pdr\_mod.f90:\allowbreak{}3326}), and the seed.

\rtn{calibratePDRdust} (\src{pdr\_mod.f90:\allowbreak{}1439--1662}) runs once, at the
first PDR iteration. From the cell whose $\AV/N_{\mathrm H}$ is closest to the
density-weighted mean it derives the model-wide fallbacks for cells without
dust, \rtn{pdrAVfac} and \rtn{pdrRfuv}, and 3D-PDR's \rtn{UV\_fac}
$=\RFUV/1.0857$, which the build does not read (it is used only under
\kw{GRAINRECOMB}, which is not compiled; \src{pdr\_mod.f90:\allowbreak{}3308--3310}).
It reports the energy-weighted FUV albedo and the scattering-weighted
asymmetry parameter (Table~\ref{tab:dustfix}), and (\dlog{D-108}) overwrites the
extinction curve that 3D-PDR's \rtn{XLAMBDA} splines (\rtn{X\_GRID}, the
\citet{SavageMathis1979} table) with this model's $\kappa(\lambda)/\kappa_V$
at the bin nearest each tabulated wavelength.

\paragraph{3D-PDR's grain attenuation factor.} \rtn{SCATTER}
(\src{pdr/\allowbreak{}shield.F90:\allowbreak{}449--497}) is the fit of
\citet[their eqs.~1a--1d]{WagenblastHartquist1989} to the attenuation of an
FUV field incident on one side of a plane-parallel slab of absorbing and
scattering grains, which they derived with the spherical-harmonics method of
\citet{Flannery1980} for grains of albedo 0.3 and asymmetry parameter 0.8:
\begin{equation}
  f_g(\tau_\lambda) = \begin{cases}
     e^{-0.7514\,\tau_\lambda}, & \tau_\lambda<1,\\[2pt]
     \sum_{i=1}^{5} a_i\,e^{-k_i\tau_\lambda}, & \tau_\lambda\ge1,
  \end{cases}
  \qquad \tau_\lambda = \frac{\AV}{1.086}\,X(\lambda),
\end{equation}
with $(a_i)=(2.006,-1.438,0.7364,-0.5076,-0.0592)$,
$(k_i)=(0.8490,1.013,1.282,2.005,5.832)$ and $X(\lambda)=\tau_\lambda/\tau_V$
the extinction curve splined by \rtn{XLAMBDA}; the two branches agree at
$\tau_\lambda=1$ to $2\times10^{-4}$. The field thus falls as
$e^{-0.75\tau_\lambda}$ to $e^{-0.85\tau_\lambda}$ rather than $e^{-\tau_\lambda}$,
because forward scattering keeps photons in the beam. The MC transfer does
this exactly, with the model's own albedo and phase function, which differ
from 0.3 and 0.8 (Table~\ref{tab:dustfix}). In the 3D build the photorates are
evaluated at $\AV=0$ (Section~\ref{sec:rays}), where $f_g=1$, so \rtn{SCATTER}
matters only for the 1D build, the seed (Section~\ref{sec:seed}) and
diagnostics.

\subsection{Grain area per H nucleus for \texorpdfstring{\Htwo}{H2} formation}
\label{sec:grainarea}

The formation rate of \citet{CazauxTielens2002,CazauxTielens2004}
(Section~\ref{sec:formation}) is linear in the geometric grain cross section
per H nucleus of each material, for which 3D-PDR ships diffuse-ISM constants
($8.473\times10^{-22}$ and $7.908\times10^{-22}$\,cm$^2$ per H for silicate and
graphite, \src{pdr/\allowbreak{}h2\_form.F90:\allowbreak{}37--39}). \rtn{measureH2grainArea}
(\src{pdr\_mod.f90:\allowbreak{}1266--1338}) computes instead, once, the mean
geometric cross section per grain of each dust component, split by material,
\begin{equation}
  \bar\sigma_{\mathrm{mat}} = \sum_{s\in\mathrm{mat}} X_s \sum_{a=a_{\min,s}}^{a_{\max,s}} \pi a^2\, w(a)
  \quad[\mathrm{cm^2}],
\end{equation}
with $X_s$ the number abundance of species $s$ within the component and
$w(a)=n(a)\Delta a/\sum n(a)\Delta a$ the number fraction of grains in size
bin $a$, normalised over the whole size grid
(\src{ph\_mod.f90:\allowbreak{}999--1012}) and summed over the species' own size range.
This is the same weighting from which the opacity is built, so the area and
the extinction describe the same grains. Species are classified by their
label (\rtn{isCarbonGrain}, \rtn{isSilicateGrain},
\src{pdr\_mod.f90:\allowbreak{}1228--1247, 1341--1360}): carbonaceous if it contains \texttt{graph},
\texttt{grph}, \texttt{amc}, \texttt{carb}, \texttt{pah}, \texttt{sic},
\texttt{gr1} or \texttt{gr2}, silicate if it contains \texttt{sil}, \texttt{oliv},
\texttt{pyr}, \texttt{enst}, \texttt{forst}, \texttt{quartz} or \texttt{oss}; an
unrecognised label goes to the silicate channel with a warning.
\rtn{dustPerCell} then stores per cell
$\sigma^{\mathrm{sil}}_{\mathrm H}=(N_d/n_{\mathrm H})\,\bar\sigma_{\mathrm{sil}}$ and
$\sigma^{\mathrm{gra}}_{\mathrm H}=(N_d/n_{\mathrm H})\,\bar\sigma_{\mathrm{gra}}$
[cm$^2$ per H] (\rtn{grid\%dustSigSil}, \rtn{grid\%dustSigGra}); a dust-free cell
gets zero, not the upstream constant (\dlog{D-195}). On the fixtures the area
is 1.2 times the upstream silicate constant on \texttt{hiipdr} and 0.69 times
the graphite one on \texttt{pn150k}, each with no grains of the other material
(Table~\ref{tab:dustfix}); the formation rate follows linearly. The geometric
area is the right weight for the collision rate of H atoms with grains; the
efficiency per unit area is applied at one mean grain temperature
(Section~\ref{sec:Td}).

\subsection{Grain albedo in the Lyman--Werner band}

Cosmic-ray-induced photoreactions are divided by $1-\omega$
(Section~\ref{sec:cr}), with $\omega$ the grain albedo at the wavelengths of
the \Htwo{} Lyman--Werner emission that cosmic rays excite: a photon is then
scattered on average $1/(1-\omega)$ times before a grain absorbs it, and
can photodissociate or photoionise at each step. Instead of 3D-PDR's
constant $\omega=0.42$ (\src{pdr\_mod.f90:\allowbreak{}354}), each cell uses
(\rtn{grainAlbedoLW}, \src{update\_mod.f90:\allowbreak{}4658--4673}; handed over at
\src{pdr\_mod.f90:\allowbreak{}3691}, \dlog{D-280})
\begin{equation}
  \omega = \frac{\sum_{k\in\mathrm{LW}} \kappa^{\mathrm{sca}}_k\,\Delta\nu_k}
                {\sum_{k\in\mathrm{LW}} (\kappa^{\mathrm{abs}}_k+\kappa^{\mathrm{sca}}_k)\,\Delta\nu_k},
\end{equation}
an average flat in frequency over the window. MOCASSIN's own cosmic-ray-induced
ionisation of the linked neutrals uses the same $\omega$
(\src{update\_mod.f90:\allowbreak{}2916--2930}).

\subsection{Dust temperature}
\label{sec:Td}

MOCASSIN solves the radiative-equilibrium temperature $T_d(s,a)$ of every
species $s$ and size $a$ in every cell from the local MC field, by equating
the absorbed power (plus the heating by resonance lines and gas collisions,
less the photoelectric losses) to $4\pi\int C_{\mathrm{abs}}B_\nu(T_d)\,\dd\nu$; stochastic
heating is off by default (keyword \kw{quantumHeatGrain}; \dlog{L-8}). The
chemistry and the level model use one mean temperature per cell,
\begin{equation}
  T_d \equiv \rtn{Tdust(0,0,cell)} = \sum_s X_s \sum_{a=a_{\min,s}}^{a_{\max,s}} w(a)\, T_d(s,a)
\end{equation}
(\src{update\_mod.f90:\allowbreak{}4546--4555}), floored at $T_{\mathrm{CMB}}=2.725$\,K
because MOCASSIN transports no cosmic microwave background
(\src{pdr\_mod.f90:\allowbreak{}3338}; \dlog{L-30}). Three properties of this mean
matter for the molecules:
\begin{itemize}
\item it is weighted by grain \emph{number}, which favours the smallest
  grains, whereas the formation rate is a sum over grains weighted by
  \emph{area};
\item the size weights are those of the whole size grid, not renormalised to
  a species' size range, so a grid wider than the species' range enters with
  zero-temperature bins: on \texttt{test\_Problem}'s silicates (0.07--5\,\textmu m
  on a $10^{-4}$--5\,\textmu m grid) the mean came out at $(2$--$8)\times10^{-6}$\,K
  and no \Htwo{} formed anywhere (\dlog{H-035}). After \texttt{8fc1512} the user
  decided this: \dlog{D-414} renormalises each species' mean to its own bins and
  weights the species by the grains they have, which is exact (byte-identical
  outputs) where a species fills the grid;
\item one efficiency $\epsilon(T_d)$ is applied to the area of all sizes, where
  the grains have a spread of temperatures; since $\epsilon$ is non-linear,
  $\epsilon(\langle T_d\rangle)\ne\langle\epsilon(T_d)\rangle$.
\end{itemize}
The temperature is the second route by which the field acts on \Htwo{}, and
a steep one. Evaluating the efficiency of Section~\ref{sec:formation} with
3D-PDR's surface energies: on silicate $\epsilon\approx1$ for
$6\lesssim T_d\lesssim15$\,K, 0.75 at 25\,K, 0.59 at 30\,K, 0.37 at 40\,K and 0.25
at 50\,K; on graphite $\epsilon\approx1$ for $10\lesssim T_d\lesssim40$\,K, 0.91 at
50\,K and 0.43 at 100\,K. Below about 5.5\,K (silicate) and 9\,K (graphite)
$\epsilon$ collapses, because newly formed \Htwo{} no longer desorbs (the
term $F_1$ of Section~\ref{sec:formation} exceeds unity); the floor at
$T_{\mathrm{CMB}}$ does not prevent this. On \texttt{hiipdr}, whose silicate
grains sit at 26--37\,K in the PDR (\dlog{L-33}), $\dd\ln\epsilon/\dd\ln T_d\approx-1.3$
to $-1.7$: a 10\% error in $T_d$ is a 13--17\% error in the formation rate,
and the formation rate sets the depth of the H/\Htwo{} transition together
with the shielding. The grains of \texttt{pn150k} are 0.92--0.93 of Cloudy's
temperature through the PDR (\dlog{L-19}), where graphite's efficiency is flat.

% ======================================================================
\section{Geometry: rays, extinction and shielding columns}
\label{sec:geometry}
% ======================================================================

The MC field already contains the attenuation by dust (absorption and
scattering) and by the gas continua. What it does not contain is the
absorption in the \Htwo{} and CO lines, which the MC does not transport. That
absorption is applied through column densities integrated along the paths the
photons took: the straight line from each source for the direct beam, and the
path source $\to$ scattering point $\to$ cell for the scattered light. This
section describes how those paths and columns are built;
Sections~\ref{sec:h2model} and \ref{sec:photorates} apply them.

\subsection{Exact ray--cell traversal}
\label{sec:raycells}

Every ray in the bridge is traced by \rtn{rayCells}
(\src{pdr\_mod.f90:\allowbreak{}1710--1833}, \dlog{D-397}). MOCASSIN's cells are centred on
the axis values, so a cell's faces are the midpoints between neighbouring axis
values and the grid ends at the first and last values (the boundary cells are
half-cells). From a start point inside the grid and a unit vector $\hat u$,
the routine computes for each axis the distance along $\hat u$ to the next
face and repeatedly steps through the nearest one (a three-dimensional digital
differential analyser in the manner of \citealt{AmanatidesWoo1987}), returning
the list of cells crossed and the exact length $\ell$ of the ray in each, at a
cost proportional to the number of cells crossed. The first entry is the cell
the ray starts in; inactive cells (outside the nebula or in a cavity) are
returned with index zero and contribute nothing. The ray ends after a given
length or at the edge of the grid. For a \kw{symmetricXYZ} octant the lower
faces are planes of symmetry: when \rtn{lgMirror} is set the ray is reflected
there ($u_q\to-u_q$) and continues, a segment that re-enters the cell it just
left being merged with it; otherwise the ray ends at the edge. This replaced
fixed-step marches that sampled each cell an integer number of times and
missed the part of the neighbouring cell nearest the target, which had made
the source-ward $N(\mathrm{H_2})$ through \texttt{hiipdr}'s transition 0.46--0.75
of its exact value; with the exact paths that transition moved 0.25\,pc
towards the star (\dlog{D-397}, \dlog{H-025}, \dlog{L-45}).

\subsection{Source-ward extinction and molecular columns}
\label{sec:avradial}

\rtn{computeAVradial} (\src{pdr\_mod.f90:\allowbreak{}1859--2015}), called from
\rtn{computeAVdust} (\src{:\allowbreak{}2410--2523}) for point sources, traces for every
active cell and every source $s$ a ray from the cell centre towards the
source (no mirror) and accumulates
\begin{equation}
  \tau_{V,s} = \sum_{\mathrm{cells}} (\kappa^{\mathrm{abs}}_V+\kappa^{\mathrm{sca}}_V)\,\ell,
  \qquad
  N_{X,s} = \sum_{\mathrm{cells}\neq\mathrm{target}} n_{\mathrm H}\, x_X\,\ell,
  \quad X\in\{\mathrm{H_2},\,\mathrm C,\,\mathrm{CO}\},
\end{equation}
with the abundances left by the previous chemistry. $\tau_{V,s}$ runs to the
cell centre, the target's half-chord included. The target cell is excluded
from $N_{X,s}$ (\dlog{D-088}): including its near half made its \Htwo{}
shield itself, a positive feedback (more \Htwo, less dissociation, more \Htwo)
that, with MOCASSIN's large cells, drove the surface cell to complete
self-shielding. The column is therefore the one to the cell's near face. It
misses the cell's own contribution, an error that shrinks with the cell size
instead of feeding back: on \texttt{hiipdr}'s $31^3$ grid a cell is 0.648\,pc,
0.104\,mag of $\AV$ (\dlog{H-026}) and $2\times10^{20}$ H nuclei\,cm$^{-2}$, so
half a cell can hold $5\times10^{19}$\,cm$^{-2}$ of \Htwo{} and, at a CO
abundance of $10^{-4}$, $10^{16}$\,cm$^{-2}$ of CO, columns at which both
self-shield strongly (the CO factor of \citealt{vanDishoeckBlack1988} is
$\approx0.05$ at $N(\mathrm{CO})=10^{16}$\,cm$^{-2}$). The level model averages the transmission through the cell's own
column instead (Section~\ref{sec:ownavg}); the CO photorate and the
chemistry's own \Htwo{} diagnostic do not: their column lags the column to
the cell centre by half a cell (\dlog{D-088}), so they are under-shielded in
the cells of their own transitions on coarse grids, which places those
transitions deeper by up to about half a cell. The length of the source-ward ray inside
the target cell, from its near face to its centre, is kept separately as
$\ell_{\mathrm{own}}$ for the scattering sites (\dlog{D-398}).

The sources are combined with the geometric weights $f_s = r_s^{-2}$
(dilution only, \src{:\allowbreak{}1986}):
\begin{align}
  \AV &= -\frac{1.0857}{\RFUV}\,\ln\!\left[\frac{\sum_s f_s\, e^{-\RFUV\tau_{V,s}}}{\sum_s f_s}\right],
  &
  N_X &= \frac{\sum_s f_s N_{X,s}}{\sum_s f_s},
  &
  \ell_{\mathrm{own}} &= \frac{\sum_s f_s\,\ell_{\mathrm{own},s}}{\sum_s f_s}
  \label{eq:avradial}
\end{align}
(\src{:\allowbreak{}1996--2009}), stored as \rtn{AVeff}, \rtn{pdrColH2}, \rtn{pdrColC},
\rtn{pdrColCO} and \rtn{difSrc(:,6)}. For one source these are the plain
column and $\AV$, and all the fixtures have one source. For several, the $\AV$
is the one that reproduces the summed transmitted FUV of sources of
\emph{equal} luminosity, and the columns are means weighted by geometry alone
(neither the luminosities nor the attenuation of the sources enter $f_s$),
applied to the summed direct field through shielding functions that are not
linear in the column; the error of this is not measured. An outward ray along
the same direction (with mirrors) gives the dust column to the edge,
\rtn{AVback}.

$\AV$ is a coordinate and a diagnostic in the 3D build: the photorates do not
use it (Section~\ref{sec:rays}). It is a geometric dust column, not an
inversion of the MC field, which would under-read it by a factor 1.8 because
the field carries scattered light (\dlog{D-086}, \dlog{D-087}); TORUS-3DPDR also
derives $\AV$ from the column density along rays \citep[their Eq.~2]{Bisbas2015}.
It is measured from the source and so includes the dust of the ionised gas:
\texttt{hiipdr}'s ionisation front sits at $\AV=0.48$ (\dlog{D-260}). A 1D PDR
model, whose $\AV$ starts at its illuminated face, has to be compared at
$\AV$ measured from the front.

Under plane-parallel illumination \rtn{computeAVdust} integrates $\AV$ and the
columns along the beam direction itself (\src{:\allowbreak{}2432--2521}); there the
columns run to the cell centre, the cell's near half included
(\src{:\allowbreak{}2492--2507}), unlike the point-source branch. The level model does
not run in that geometry.

\subsection{Snapshot for the diffuse walk}

\rtn{snapDiffuseSources} (\src{pdr\_mod.f90:\allowbreak{}2082--2121}) is called by
\rtn{computeAVdust} before any cell is solved and stores per cell
(\rtn{difSrc}):
\begin{equation}
  n(\mathrm{H_2}),\ n(\mathrm C),\ n(\mathrm{CO})\ [\mathrm{cm^{-3}}];\qquad
  \kappa_{\mathrm{ext}} = \RFUV\,\kappa_V\ [\mathrm{cm^{-1}}];\qquad
  \epsilon = \sum_{k\in\mathrm{FUV}} \kappa^{\mathrm{sca}}_k\,\Jste_k\ [\mathrm{erg\,s^{-1}\,cm^{-3}}],
\end{equation}
the number densities that shield, a grey FUV extinction coefficient, and the
FUV power scattered per unit volume in the band (\src{:\allowbreak{}2099--2117}). For
isotropic scattering $\epsilon/4\pi$ is the band-integrated scattering
emissivity. It is computed from the full local MC field, direct and diffuse,
so it contains every order of scattering.

\subsection{Scattering sites of the diffuse field}
\label{sec:sites}

The diffuse FUV field is produced by scattering inside the cloud. A photon
scattered at a point $P$ and arriving at the cell $C$ has crossed the column
$N(\star\to P)+N(P\to C)$, not the column from the star to $C$
(\dlog{D-347}, \dlog{D-348}). \rtn{traceDiffuseSites}
(\src{pdr\_mod.f90:\allowbreak{}2136--2276}) represents the diffuse field arriving at the
target cell by a set of \emph{scattering sites}, built as the discretised
formal solution of the transfer equation for the intensity arriving at $C$
from each of a set of directions $\hat k$,
\begin{equation}
  4\pi I_{\hat k} = \int_0^\infty \epsilon(s)\, e^{-\tau(s)}\,\dd s ,
  \qquad \tau(s)=\int_0^s \kappa_{\mathrm{ext}}\,\dd s',
  \label{eq:formal}
\end{equation}
with $s$ the distance from $C$ along $\hat k$ (the direction the photons come
from):
\begin{enumerate}
\item \textbf{Directions.} The twelve HEALPix $N_{\mathrm{side}}=1$ pixel
  centres \citep{Gorski2005}, which divide the sphere into twelve equal solid
  angles of $\pi/3$\,sr, built directly in RING order: rays 0--3 at
  $\mu=+2/3$, $\phi=(k+\frac12)\pi/2$; rays 4--7 at $\mu=0$, $\phi=k\pi/2$;
  rays 8--11 at $\mu=-2/3$, $\phi=(k+\frac12)\pi/2$ (\src{:\allowbreak{}2152--2165}),
  with $\mu$ measured along the grid's $z$ axis. The directions are fixed in
  the grid frame.
\item \textbf{Walk.} Each ray is traced with \rtn{rayCells} from the target's
  centre, with mirror reflection at the octant's lower faces, to the edge of
  the model. Every cell crossed that scatters ($\epsilon>0$, $\kappa_{\mathrm{ext}}>0$)
  is a site, the target cell itself included (with the half-chord the ray
  crosses); the walk is stored (\rtn{wkCell},
  \rtn{wkLen}, \rtn{siteRay}, \rtn{siteIst}) for the level model.
\item \textbf{Weight.} With $\tau$ the grey FUV optical depth between the
  target's centre and the entry of the site cell and $\Delta\tau=\kappa_{\mathrm{ext}}\ell$
  across it,
  \begin{equation}
    w_{\mathrm{site}} = \epsilon\;\frac{1-e^{-\Delta\tau}}{\kappa_{\mathrm{ext}}}\;e^{-\tau}
    \qquad\bigl(\tfrac{1-e^{-\Delta\tau}}{\kappa_{\mathrm{ext}}}\to\ell(1-\tfrac12\Delta\tau)\ \text{for}\ \Delta\tau\le10^{-4}\bigr)
    \label{eq:siteweight}
  \end{equation}
  (\src{:\allowbreak{}2207--2217}): Equation~(\ref{eq:formal}) integrated over the
  site's segment with $\epsilon$ and $\kappa_{\mathrm{ext}}$ constant in the cell.
  Since the directions subtend equal solid angles,
  $4\pi J_{\mathrm{formal}}=\frac1{12}\sum_{\mathrm{sites}}w_{\mathrm{site}}$.
\item \textbf{Two-leg column.} For each species $X$ (\src{:\allowbreak{}2219--2239}),
  \begin{equation}
    N_{X,\mathrm{site}} = N_X^{\star}(P) + n_X(P)\,\ell_{\mathrm{own}}(P)
       + \sum_{\mathrm{cells\ between}} n_X\,\ell + \tfrac12\,n_X(P)\,\ell_P ,
    \label{eq:twoleg}
  \end{equation}
  the source-ward column of the site's cell to its near face, plus the site
  cell's own gas from its near face to its centre (\dlog{D-398}), plus the
  cells crossed between the target and the site, plus half of the site cell
  along the walk. The target cell is excluded from both legs, so a site in the
  target cell carries the direct beam's column.
\end{enumerate}

Only the \emph{relative} weights are used: the magnitude of the diffuse field
is the MC's $\chiF$ (Equation~\ref{eq:chisplit}), which the sites distribute
over directions and columns, each site becoming a ray of the chemistry with
field $\chiF\,w_{\mathrm{site}}/\sum w$ (Section~\ref{sec:rays}). The
distribution is the formal solution only under the assumptions built into
Equation~(\ref{eq:siteweight}): isotropic scattering, whereas the grains of
the fixtures have $g=0.61$ and 0.29 (Table~\ref{tab:dustfix}), so the real
diffuse field arrives preferentially from the source side; a grey extinction
equal to the source-weighted band mean, whereas the window in which \Htwo{} and
CO dissociate has a larger extinction and the local field at depth a smaller
one (Section~\ref{sec:dust}); an emissivity made by dust scattering alone,
whereas $\chiF$ also holds gas emission; uniform emissivity across a cell; and
the twelve-direction quadrature. As a check of these assumptions, the formal
solution's own magnitude, $4\pi J_{\mathrm{formal}}/(c\,\uD)$, recovers 52--130\%
of the MC's $\chiF$, falling with depth, for reasons not yet identified (gas
two-photon emission, the phase function and the quadrature were named as
candidates; the source-weighted grey extinction, which overstates the
extinction at depth, is another; \dlog{L-4}). The field arrives markedly
anisotropically: the brightest of the twelve directions carries 1.9--3.1 times
the mean (\dlog{D-348}). What the assumptions affect is the distribution of
the diffuse field over the sites' columns, not its magnitude, and that
distribution acts on the photorates only through the line shielding: at fixed
abundances the two-leg columns took back 5--32\% of the increase of the \Htwo{}
photodissociation rate that the diffuse field's magnitude produced
(\dlog{D-348}).

\subsection{Angular quadrature}
\label{sec:healpix}

The twelve directions of the sites and of the line rays are the coarsest
HEALPix tessellation. Their moments are not isotropic: the mean fourth power
of a direction's component is 0.132 along the ring axis ($z$) against 0.218
across it, where an isotropic set gives 0.2 (\dlog{D-410}). The result is a
pattern fixed in the grid frame: of the cells that the octant's symmetry makes
equivalent, the one nearest the plane $z=0$ holds about 11\% more \Htwo{}
before and across the H/\Htwo{} transition, on both fixtures, independently of
the packet count (\dlog{L-53}). Rotating the set moves the excess to the new
ring axis, and 48 directions ($N_{\mathrm{side}}=2$, fourth moments 0.188 and
0.201) remove it at 1.48 times the run time, moving every integrated quantity
by less than 1\%; the twelve were kept because no line map carries the pattern
(\dlog{D-410}). 3D-PDR and TORUS-3DPDR also use HEALPix rays, twelve
($\ell=0$) in TORUS-3DPDR's production runs \citep{Bisbas2012,Bisbas2015}.

\subsection{Line-escape rays}
\label{sec:lineRays}

\rtn{traceLineRays} (\src{pdr\_mod.f90:\allowbreak{}2327--2376}, \dlog{D-377}) traces the
same twelve directions from the target cell to the edge of the model, with
mirrors at the lower faces, skipping inactive cells, and returns for each ray
the cells crossed and the path length in each; the first entry is the cell
itself with the half of it the ray crosses. These rays give the escape
probabilities of the coolant lines (Section~\ref{sec:thermal}) and of the
\Htwo{} lines (Section~\ref{sec:escape}), each crossed cell contributing with
its own level populations. Before \dlog{D-377} the 3D flavour's line escape
was a slab with four of the twelve rays shut, which capped the escape
probability of an optically thin line at 2/3 (\dlog{L-27}).

\subsection{Comparison with 3D-PDR and TORUS-3DPDR}

In 3D-PDR \citep{Bisbas2012} the FUV field is a boundary condition: at every
point the field incident on the cloud surface is attenuated along HEALPix rays
traced to the boundary, $\chi(p)=N_\ell^{-1}\sum_{q}\chi_0(q)\,e^{-\tau_{\mathrm{UV}}\AV(q)}$
with $\tau_{\mathrm{UV}}=3.02$ in the benchmark comparisons
\citep[their Eq.~1]{Bisbas2015}, and the shielding columns are integrated
along the same rays. For a source inside the cloud these rays are the wrong
paths: a ray pointing away from the source would deliver unshielded radiation
from outside a nebula that has no outside source (\dlog{D-342}). TORUS-3DPDR
\citep{Bisbas2015} takes the FUV field of every cell from its MC path-length
estimator (912--2400\,\AA, stored as two vectors for the positive and negative
directions), derives $\AV$ and the column densities along HEALPix rays from
the density field, and leaves the dust's contribution to the diffuse field to
later work (their Sect.~3.1). The present coupling differs from both in that the dust
attenuation and scattering are the MC's, the field is split into its direct
and scattered parts, and the line-shielding columns follow the photons'
paths: source-ward for the direct beam and through the scattering sites for
the diffuse field.

% ======================================================================
\section{The \texorpdfstring{\Htwo}{H2} level model}
\label{sec:h2model}
% ======================================================================

The level model computes, in every cell, the fractional populations of the
302 rovibrational levels of the electronic ground state $X\,^1\Sigma_g^+$ of
\Htwo{} under the cell's own radiation field, gas temperature and colliders,
and from them the Solomon photodissociation rate per molecule that the
chemistry uses (Section~\ref{sec:photorates}) and three \Htwo{} heating terms
(Section~\ref{sec:thermal}). The physics is in \src{h2\_mod.f90} (1570 lines);
the driver that supplies the field, the columns, the colliders and the
formation distribution is \rtn{h2LevelModel} with its internal routine
\rtn{solveCell} (\src{pdr\_mod.f90:\allowbreak{}5201--6014}). It is a re-implementation,
from the data and conventions of Cloudy's large \Htwo{} model
\citep{Shaw2005}, of the line-by-line treatment of ultraviolet pumping and
self-shielding that the Meudon PDR code \citep{LePetit2006} and the models of
\citet{DraineBertoldi1996} also use, in place of a shielding function fitted
to such models. It was validated level by level against Cloudy on one-zone and
slab models (\dlog{H-003}--\dlog{H-006}). Those comparisons isolate three
deliberate departures from Cloudy, each discussed below: the closure of the
level balance (Section~\ref{sec:balance}), the damping width of the lines
(Section~\ref{sec:shield}) and the treatment of line overlap
(Section~\ref{sec:overlap}).

\paragraph{Physical picture.} A molecule in level $l$ of $X$ absorbs a
far-ultraviolet photon in a Lyman ($B\,^1\Sigma_u^+$--$X$) or Werner
($C\,^1\Pi_u$--$X$) line, or in one of the weaker $B'\,^1\Sigma_u^+$--$X$ and
$D\,^1\Pi_u$--$X$ lines; out of $v=0$ the Lyman and Werner lines lie between 912
and 1110\,\AA{} (11.2--13.6\,eV), and the $B'$ and $D$ systems reach below
912\,\AA. Within about a nanosecond the excited level $u$ decays, either into
the vibrational continuum of $X$, which dissociates the molecule into two H(1s)
atoms (the Solomon process), or into a bound level of $X$, mostly
vibrationally excited, from which a slow electric-quadrupole cascade
($A=3\times10^{-11}$--$6\times10^{-6}$\,s$^{-1}$, median $2\times10^{-8}$\,s$^{-1}$ in
\src{transprob\_X.dat}) competes with collisions. In an unshielded field of the
Draine shape the model dissociates one molecule per 8.1 absorptions, at
$6.36\times10^{-11}$\,s$^{-1}$ per unit Draine field (\dlog{H-007}); the mean
dissociation probability per absorption is 0.12 in the computations of
\citet{Sternberg2014} and about 0.15 when averaged through complete
photodissociation fronts \citep{DraineBertoldi1996}. The strong lines out of a
level become optically thick at level columns of order
$10^{14}$\,cm$^{-2}$, saturate, and above $\sim10^{17}$--$10^{18}$\,cm$^{-2}$ absorb
mainly in their Lorentzian wings, where above $N(\mathrm{H_2})\sim10^{20}$\,cm$^{-2}$
the lines of different levels overlap \citep{DraineBertoldi1996}. The
dissociation rate in a cell therefore depends on the field at every line, on
the column of every level along the paths of the photons, on the line widths,
and on the populations of the absorbing levels, which the pumping itself
changes. The model treats all of these line by line.

\subsection{Data}
\label{sec:h2data}

All data are Cloudy c25.00's (\src{dataPDR/\allowbreak{}h2/\allowbreak{}}, copied
unmodified; provenance in \src{dataPDR/\allowbreak{}h2/\allowbreak{}README.md}). The
reading rules are Cloudy's: a line starting with \texttt{\#} is a comment, and the
data of the excited-state files end at the first empty line. The sources in
Table~\ref{tab:h2data} are those that the files' own headers give, completed
where the header is not the primary source.

\begin{longtable}{@{}L{0.33\linewidth}L{0.28\linewidth}L{0.33\linewidth}@{}}
\caption{Data of the \Htwo{} level model. Temperature ranges are those of the
tables; coverage is given as the upper levels present.}\label{tab:h2data}\\
\toprule
Content & File(s) & Source \\
\midrule
\endfirsthead
\toprule
Content & File(s) & Source \\
\midrule
\endhead
energies of the $X$ rovibrational levels (302 levels, $v\le14$, $J\le31$; sorted by energy) & \src{energy\_X.dat} & \citet[their Table~1]{Komasa2011}, compiled by K.~Kaplan (header) \\
radiative transition probabilities within $X$: electric quadrupole, plus magnetic-dipole components of 30 $v=1\to0$, $\Delta J=0$ transitions listed separately and summed & \src{transprob\_X.dat} & \citet{Wolniewicz1998}; magnetic dipole \citet{PachuckiKomasa2011} \\
H\,+\,\Htwo{}: 1432 transitions, $v\le3$, $J\le17$, 728 of them ortho$\leftrightarrow$para (reactive); 100--5000\,K & \src{coll\_rates\_H\_15.dat} & \citet{Lique2015} \\
He\,+\,\Htwo{}: 3251 transitions, $v\le14$, $J\le30$, no spin change; 1--$10^4$\,K & \src{coll\_rates\_He\_ORNL.dat} & ORNL data set; the header cites ``Lee, T.\ G.\ et al.\ 2007, ApJ, in preparation''; no published source for this table was identified \\
ortho- and para-\Htwo\,+\,\Htwo{}: 240 and 241 transitions, $v=0$, $J\le31$; 1--$10^4$\,K & \src{coll\_rates\_H2ortho\_ORNL.dat}, \src{coll\_rates\_H2para\_ORNL.dat} & \citet{Wan2018} \\
\Hp\,+\,\Htwo{}: 44 transitions, $v=0$, $J\le9$, 24 of them ortho$\leftrightarrow$para (proton exchange); 3--6000\,K & \src{coll\_rates\_Hp.dat} & \citet{Gerlich1990} \\
energies of the rovibrational levels of $B$, $C^+$, $C^-$, $B'$, $D^+$, $D^-$ (1592 levels) & \src{energy\_\{B,C\_plus,\ldots\}.dat} & derived from \citet{Abgrall2000} (headers) \\
transition probabilities of the 32\,453 lines to $X$ & \src{transprob\_\{\ldots\}.dat} & \citet{Abgrall2000}, from the band-system calculations of \citet{Abgrall1993} ($B$--$X$, $C$--$X$) and \citet{Abgrall1994} ($B'$--$X$, $D$--$X$) \\
spontaneous radiative dissociation rate $D_u$ [s$^{-1}$] and mean kinetic energy of the atoms [eV] per excited level & \src{dissprob\_\{\ldots\}.dat} & \citet{Abgrall2000} \\
dissociation energy of $X$, $D=36\,118.11$\,cm$^{-1}$ (4.478\,eV; the only value of the file the model uses) & \src{energy\_dissoc.dat} & \citet{Sharp1970} per the header (which dates it 1971), with values revised in Cloudy in 2002 \\
\bottomrule
\end{longtable}

Reading: \rtn{h2ReadX} (\src{h2\_mod.f90:\allowbreak{}206--334}) and \rtn{h2ReadElec}
(\src{:\allowbreak{}366--516}). Statistical weights are $g=(2J+1)(2I+1)$, with
nuclear-spin factor $2I+1=3$ for ortho and 1 for para. In $X$ ortho levels have
odd $J$; in the excited states $B$, $B'$ and the $\Pi^+$ components $C^+$, $D^+$
ortho levels have even $J$, and in the $\Pi^-$ components $C^-$, $D^-$ odd $J$,
following Cloudy \citep{Shaw2005}: this is the assignment under which every
allowed transition ($\Delta J=\pm1$ to the $\Sigma_u^+$ states and the $\Pi^+$
components, $\Delta J=0$ to the $\Pi^-$ components) conserves the nuclear spin,
and the reader counts the lines that would not (none, \dlog{H-004}). The $\Pi$
states have no $J=0$ level. The $D^-$ state stops at $v=2$ because Cloudy's file
has an empty line before the levels it marks ``above the dissociation limit'';
3138 $D^-$ lines are dropped with them, and 392 lines whose lower level lies
beyond the 302 levels of $X$ (\dlog{H-004}). The dissociation rates $D_u$ are
those of spontaneous emission into the vibrational continuum of $X$
\citep{Abgrall2000}. Apart from 45 high-$J$ levels of $D^-$ ($v\le2$), all the
excited levels lie below the H(1s)\,+\,H($n=2$) limit (118\,376\,cm$^{-1}$ in
\src{energy\_dissoc.dat}), where this radiative channel is the only way they
dissociate; predissociation, which acts above that limit, is not in the data.
A transition listed
twice adds its components, as Cloudy does; reading the $X$ file without adding
the magnetic-dipole components had doubled the $v=1$ populations at low density
(\dlog{H-003}). The lines are stored sorted by upper level (counting sort) and
indexed by wavenumber (heapsort) for the overlap calculation.

\subsection{Statistical equilibrium and the formation closure}
\label{sec:balance}

Let $n_i$ be the fraction of the cell's \Htwo{} molecules in level $i$ of $X$
($\sum_i n_i=1$). The populations solve (\rtn{balance},
\src{h2\_mod.f90:\allowbreak{}1456--1477})
\begin{equation}
  \sum_{j\neq i} n_j R_{ji} - n_i\Bigl(\sum_{j\neq i} R_{ij} + k_i\Bigr)
  + \phi_i \sum_j n_j k_j + s_i = 0,\qquad i=1,\ldots,302,
  \label{eq:balance}
\end{equation}
where $R_{ij}$ [s$^{-1}$] is the total rate per molecule from level $i$ to
level $j$ (collisions, spontaneous decay within $X$ and the fluorescent
transfer through the excited states; Sections~\ref{sec:coll} and
\ref{sec:pump}), $k_i$ [s$^{-1}$] the Solomon dissociation rate out of level
$i$, $\phi_i$ the fraction of newly formed molecules that enter level $i$
($\sum_i\phi_i=1$; Section~\ref{sec:formdist}) and $s_i$ [s$^{-1}$] an optional
further source per molecule, not used by \rtn{solveCell}. Fluorescent returns
to the absorbing level ($R_{ii}$) drop out. Summed over $i$, the transfer terms
cancel and so do the dissociation and formation terms, so the system
conserves $\sum_in_i$: the term $\phi_i\sum_jn_jk_j$ is the closure \emph{every
molecule destroyed by the Solomon process is replaced by one formed}. The
level model determines how the molecules are distributed over the levels;
how many there are is the chemistry's (Section~\ref{sec:chemistry}).

\paragraph{Numerics.} The equation of level 1 is replaced by $\sum_in_i=1$
and the $302\times302$ system is solved directly by Gaussian elimination with
partial pivoting (\rtn{gaussSolve}, \src{:\allowbreak{}1536--1568}; about
$n^3/3\approx9\times10^6$ operations per solution). \rtn{h2Solve}
(\src{:\allowbreak{}1516--1533}) assembles $R=C+A+W$ (collisions, radiative decay,
fluorescent transfer) and the $k_i$, and calls it. The populations of the
validation runs were never negative and summed to unity to the printed
precision (\dlog{H-003}). Cloudy instead solves its lowest 70 levels by matrix
inversion and the others one by one (\dlog{H-010}).

\paragraph{Validity of the closure.} In steady state the chemistry's formation
rate per molecule, $R_{\mathrm f}\,n_{\mathrm H}\,n(\mathrm H)/n(\mathrm{H_2})$, equals
the sum of all destruction rates per molecule, of which $k_{\mathrm{Sol}}$ is
one; the closure sets the formation flux into the levels equal to
$k_{\mathrm{Sol}}$ alone. It is therefore exact where photodissociation
dominates the destruction (93--100\% of the \Htwo{} destruction out to 0.43\,pc
in \texttt{pn150k}, \dlog{H-009}), and underestimates the formation flux by
the factor $k_{\mathrm{Sol}}/(k_{\mathrm{Sol}}+k_{\mathrm{other}})$ elsewhere. The
network's three \Htwo\,+\,CRP channels alone give $k_{\mathrm{other}}=0.99\,\zeta$,
$1.0\times10^{-16}$\,s$^{-1}$ at the default \kw{pdrZeta}, so the closure fails
where the Solomon rate is shielded below a few $10^{-16}$\,s$^{-1}$, a reduction by
more than $10^5$ from its unshielded value at $\chi\sim1$. There the formation
pumping of the $v\ge1$ levels, the infrared emission it produces and the part
of the formation energy that reaches the gas by collisional de-excitation are
underestimated; the Solomon rate itself is affected only through the
populations of the absorbing levels, and the kinetic part of the formation
energy follows the chemistry's own formation rate
(Section~\ref{sec:thermal}; Section~\ref{sec:limitations}). Cloudy closes its
level system differently: it applies its network's Solomon rate a second time
as a uniform sink on the levels below 4100\,cm$^{-1}$, solves without a sum
rule and rescales all populations by one factor (1.06--1.11 in one-zone models
with $n_{\mathrm H}\le10^4$\,cm$^{-3}$ and fields up to 100 Habing). Put through that
scheme, this model's rates reproduce Cloudy's populations to 0.002--0.018\,dex
in 15 of 16 such models, against a median difference of 0.07\,dex with its own
closure (\dlog{H-004}); Cloudy's closure raises the pumped line intensities by
about 1.2 (\dlog{H-009}).

\subsection{Collisional and radiative rates within \texorpdfstring{$X$}{X}}
\label{sec:coll}

\rtn{h2CollRad} (\src{h2\_mod.f90:\allowbreak{}1369--1415}) adds, for every pair of
levels $u>l$ (energies $E$ in cm$^{-1}$),
\begin{equation}
  R_{ul} \mathrel{+}= \sum_c n_c\,k^c_{ul}(T) + A_{ul},\qquad
  R_{lu} \mathrel{+}= \sum_c n_c\,k^c_{ul}(T)\,\frac{g_u}{g_l}\,e^{-hc(E_u-E_l)/kT},
  \label{eq:collrad}
\end{equation}
with $n_c$ [cm$^{-3}$] the density of collider
$c\in\{\mathrm H,\,\mathrm{He},\,o\text{-H}_2,\,p\text{-H}_2,\,\mathrm H^+\}$, $k^c_{ul}$
[cm$^3$\,s$^{-1}$] its de-excitation rate coefficient, $A_{ul}$ [s$^{-1}$] the
radiative transition probability, and $T=\max(T_e,1\,\mathrm K)$
(\src{pdr\_mod.f90:\allowbreak{}5667}):
\begin{itemize}
\item downward coefficients from the tables of Table~\ref{tab:h2data},
  interpolated linearly in $T$ and held at the end values outside the
  tabulated range (\rtn{tabRate}, \src{h2\_mod.f90:\allowbreak{}642--660}; Cloudy's
  \texttt{InterpCollRate} convention). Below 100\,K the H\,+\,\Htwo{} rates are
  therefore those of 100\,K, which overestimates the reactive (ortho--para)
  channels, whose barrier makes them fall steeply with temperature; above
  5000\,K (H), 6000\,K (\Hp) and $10^4$\,K (He, \Htwo) the rates are those of the
  table's end;
\item for a pair of levels of the same nuclear spin that a collider's table
  does not contain, the ``g-bar'' estimate of \citet[their eq.~1 and
  Table~2]{Shaw2005},
  \begin{equation}
    \log_{10}\bigl(k^c_{ul}/\mathrm{cm^3\,s^{-1}}\bigr) = y_0 + a_g\,\bigl[\max(\Delta E,\,100)\bigr]^{b_g},
    \qquad \Delta E = E_u-E_l\ [\mathrm{cm^{-1}}],
    \label{eq:gbar}
  \end{equation}
  a fit, per collider, of the rate coefficients then available as a function
  of the energy gap (coefficients in Table~\ref{tab:gbar}; Cloudy's
  \texttt{GbarRateCoeff}), which \citet{Shaw2005} found to lie within 1\,dex of
  60--70\% of the rates it was fitted to. A pair that changes the nuclear spin
  and is not in a table has no collisional rate;
\item upward rates by detailed balance at the gas temperature, so that
  collisions alone relax the populations to a Boltzmann distribution at $T$;
\item spontaneous radiative decay $A_{ul}$ (quadrupole plus magnetic dipole)
  from \src{transprob\_X.dat}.
\end{itemize}

\begin{table}[H]
\centering
\caption{g-bar coefficients of Equation~(\ref{eq:gbar}) (\src{h2\_mod.f90:\allowbreak{}96--101}),
identical to Table~2 of \citet{Shaw2005}.}\label{tab:gbar}
\begin{tabular}{@{}lrrr@{}}
\toprule
Collider & $y_0$ & $a_g$ & $b_g$ \\
\midrule
H & $-9.9265$ & $-0.1048$ & 0.456 \\
He & $-8.281$ & $-0.1303$ & 0.4931 \\
$o$-\Htwo & $-10.0357$ & $-0.0243$ & 0.67 \\
$p$-\Htwo & $-8.6213$ & $-0.1004$ & 0.5291 \\
\Hp & $-9.2719$ & $-0.0001$ & 1.0391 \\
\bottomrule
\end{tabular}
\end{table}

\paragraph{Coverage.} Of the 22\,675 pairs of levels of the same nuclear spin,
the tables contain 703 for H, 3250 for He, 240 and 239 for $o$- and $p$-\Htwo{}
and 20 for \Hp{} (counted from the data files); the remaining 86--99.9\% take
Equation~(\ref{eq:gbar}). Ortho and para are coupled by collisions only through
the 728 reactive H\,+\,\Htwo{} transitions ($v\le3$) and the 24 proton-exchange
\Hp\,+\,\Htwo{} transitions ($v=0$, $J\le9$), the latter being by far the faster
\citep{Shaw2005}. The tabulated pairs include the low-$J$ levels of $v=0$ that
carry most of the column; the g-bar rates matter where collisions compete
with radiative decay in levels without data, i.e.\ high-$v$ or high-$J$ levels
at high density, and \citet{Shaw2005} found the populations of the tabulated
low levels nearly insensitive to them.

\paragraph{Colliders.} The collider densities (\src{pdr\_mod.f90:\allowbreak{}5887--5888})
are taken from the chemistry's abundances of the previous iteration,
$n(\mathrm H)=n_{\mathrm H}x(\mathrm H)$, $n(\mathrm{He})$ and $n(\mathrm H^+)$, and
from the cell's own solution for the spin split,
$n(o\text{-H}_2)=f_o\,n(\mathrm{H_2})$, $n(p\text{-H}_2)=(1-f_o)\,n(\mathrm{H_2})$, with
$f_o=\sum_{i\in\mathrm{ortho}} n_i$ iterated to self-consistency
(Section~\ref{sec:solvecell}). Electrons are not a collider, as in Cloudy;
\citet{AlemanGruenwald2011} include them in their models of the \Htwo{}
emission of the ionised region of planetary nebulae. Where electrons are a
large fraction of the colliders, in the ionised gas and at the ionisation
front, the collisional excitation of \Htwo{} is therefore underestimated, in
gas whose \Htwo{} abundance is in any case uncertain (\dlog{L-43}). The rate
coefficients and Boltzmann factors are cached for the last temperature
(\rtn{crKd}, \rtn{crBf}), so the iterations over the colliders' densities at
fixed $T$ rebuild only $R$.

\subsection{Ultraviolet pumping, Solomon dissociation and trapping}
\label{sec:pump}

A molecule in $X$ level $l$ is excited in the line $l\to u$ at the rate
\begin{equation}
  P_{lu} = B_{lu}\,\bar J_{lu} = A_{ul}\,\frac{g_u}{g_l}\,\eta_{lu}\quad[\mathrm{s^{-1}}],
  \qquad
  \eta_{lu} = \frac{c^2}{2h\nu_{lu}^3}\,J_\nu(\nu_{lu}),
  \label{eq:pump}
\end{equation}
where $\bar J_{lu}$ is the mean intensity averaged over the line profile and
$\eta_{lu}$ the photon occupation number at the line after shielding
(Sections~\ref{sec:eta} and \ref{sec:shield}); the second form follows from
the Einstein relations for a field constant across the line. In the
unshielded Draine field $\eta\approx2\times10^{-19}$ near 12\,eV.

The excited levels have radiative lifetimes of 0.5--9\,ns in these data
(total decay rates $\Gamma_u=1.1\times10^8$--$1.9\times10^9$\,s$^{-1}$, median
$\approx10^9$\,s$^{-1}$ for $B$ and $C$). Collisions cannot compete with that below
$\sim10^{14}$\,cm$^{-3}$ \citep{Shaw2005}, and the excited populations, of order
$\sum_lP_{lu}/\Gamma_u\sim10^{-18}\chi$ per molecule, are never a reservoir: every
absorption is followed at once by a decay, into the continuum of $X$
(dissociation) or into a bound level of $X$. The excited states are therefore
not solved for. With the effective decay rates
\begin{equation}
  A^{\mathrm{eff}}_t = A_t\,(1+\eta_t)\,\beta_t,
  \qquad
  \Gamma^{\mathrm{eff}}_u = D_u + \sum_{t\in u} A^{\mathrm{eff}}_t ,
  \label{eq:aeff}
\end{equation}
where $t$ runs over the lines from $u$ to bound levels of $X$, $1+\eta_t$
accounts for stimulated emission (negligible, kept for completeness), $\beta_t$
is the escape probability of a photon emitted in line $t$
(Section~\ref{sec:escape}) and $D_u$ the spontaneous radiative dissociation rate
of $u$, \rtn{h2PumpRates} (\src{h2\_mod.f90:\allowbreak{}540--572}) returns
\begin{align}
  k_l &= \sum_u P_{lu}\,\frac{D_u}{\Gamma^{\mathrm{eff}}_u}
       && \text{(Solomon dissociation out of } l\text{)},\\
  W_{ll'} &= \sum_u P_{lu}\,\frac{A^{\mathrm{eff}}_{ul'}}{\Gamma^{\mathrm{eff}}_u}
       && \text{(fluorescent transfer } l\to l'\text{)},\\
  H^{\mathrm{dis}}_l &= \sum_u P_{lu}\,\frac{D_u}{\Gamma^{\mathrm{eff}}_u}\,E^{\mathrm{kin}}_u
       && \text{(kinetic energy given to the atoms [eV\,s}^{-1}\text{])},
\end{align}
and the total absorption rate $p_l=\sum_u P_{lu}$, with $E^{\mathrm{kin}}_u$ the
mean kinetic energy released per dissociation of $u$ (\src{dissprob\_\ldots.dat}).
The first step of the fluorescent cascade is $W$; the quadrupole cascade that
follows is carried by $A$ in Equation~(\ref{eq:balance}). The Solomon rate per
molecule handed to the chemistry is
\begin{equation}
  k_{\mathrm{Sol}} = \sum_l n_l\,k_l
  \qquad(\text{\src{pdr\_mod.f90:\allowbreak{}5918}}).
  \label{eq:ksol}
\end{equation}

\paragraph{Trapping.} A photon emitted in a line $t=u\to l'$ that is optically
thick is absorbed again close by, by a molecule in $l'$, which it re-excites
to the same level $u$. In the escape-probability (``on-the-spot'')
approximation the net effect is that the decay rate of the line is reduced to
$A_t\beta_t$, and the level's other decays and its dissociation take its share
in $\Gamma^{\mathrm{eff}}_u$. Only lines that end in populated low levels are
thick, so trapping closes the decays to $v=0$ and raises the dissociation
probability $D_u/\Gamma^{\mathrm{eff}}_u$ in thick gas. It is needed: without it
the rate into $v=0$, $J=0$--5 from the excited states was up to 2200 times
Cloudy's in the shielded zones of the slab of \dlog{H-005} and the Solomon rate
10--20\% low; with it the Solomon rate comes within 3--8\% of Cloudy's
(\dlog{H-005}). The total absorption rate $p_l$ counts only the external field.
Photons that dust or an overlapping \Htwo{} line absorb before they escape are
not counted as lost (Cloudy adds such a destruction probability to the escape
probability, \citealt{Shaw2005}); this slightly overestimates the dissociation
probability of levels whose decays are trapped.

\subsection{Self-shielding of the pumping lines}
\label{sec:shield}

\paragraph{Line parameters.} For line $lu$ of wavelength $\lambda=1/\tilde\nu_{lu}$
[cm] (\rtn{h2LineParams}, \src{h2\_mod.f90:\allowbreak{}1328--1336}),
\begin{equation}
  \sigma_{lu} = \frac{g_u}{g_l}\,\frac{A_{ul}\,\lambda^3}{8\pi^{3/2}\,b}\ [\mathrm{cm^2}],\qquad
  \tau_0 = \sigma_{lu}N_l,\qquad
  a = \frac{\Gamma_u\,\lambda}{4\pi\,b},\qquad
  \Gamma_u = D_u + \sum_{t\in u}A_t ,
  \label{eq:lineparams}
\end{equation}
where $\sigma_{lu}$ is the line-centre cross section of a pure Doppler profile
of Doppler width $\Delta\nu_D=b/\lambda$, $(\pi e^2/m_ec)f_{lu}/(\sqrt\pi\Delta\nu_D)$
written with $A_{ul}$; $N_l$ [cm$^{-2}$] is the column of molecules in level $l$;
$b=\sqrt{2kT/m_{\mathrm{H_2}}+v_{\mathrm{turb}}^2}$ [cm\,s$^{-1}$] the Doppler parameter of
the target cell (\src{pdr\_mod.f90:\allowbreak{}5668}; $v_{\mathrm{turb}}$ is \kw{pdrVturb},
default 1\,km\,s$^{-1}$); and $a$ the Voigt damping parameter, the ratio of the
Lorentzian half width $\Gamma_u/4\pi$ to $\Delta\nu_D$. $\Gamma_u$ is the total
decay rate of the upper level, bound--bound plus bound--free (spontaneous
radiative dissociation), i.e.\ the inverse radiative lifetime; since the lower
level of every line is a metastable level of $X$, $\Gamma_u$ is the natural width
of the line. With these data $a\le0.03$ for $b\ge1$\,km\,s$^{-1}$ (median 0.006 at
100\,K and the default $v_{\mathrm{turb}}$).

\paragraph{The damping width.} $\Gamma_u$ exceeds a line's own $A_{ul}$ by a
factor 20 in the median over the thick lines of the \dlog{H-005} slab (16--84\%:
7.8--157), and the shielding in the damping wings scales as $a^{1/2}$
(Equation~\ref{eq:Sasym}). Cloudy damps each \Htwo{} line with its own $A_{ul}$
(\texttt{dampXvel}$\,=A_{ul}\lambda/4\pi$ in \texttt{mole\_h2\_create.cpp} of c25.00,
and in its saved damping constants, \dlog{H-005}), whereas its hydrogenic lines
use the lifetime of the upper level (\texttt{iso\_create.cpp}). The natural width
is the physical choice. Cloudy's lowers the transmission in the damping regime
by up to $\sqrt{20}\approx4.5$ (pumped lines behind more than $10^{16}$\,cm$^{-2}$ of
\Htwo{} fall to 0.18--0.41 of their intensity, \dlog{H-008}) and moves the
H/\Htwo{} crossing 0.21\,mag shallower on \texttt{hiipdr} and 0.10\,mag on
\texttt{pn150k} (\dlog{H-026}; Section~\ref{sec:limitations}). The same $b$ (the target cell's) is used for the
whole column in front of the cell and along the escape rays; this matters in
the Doppler-saturated regime, where $\tau_0\propto b^{-1}$, and not in the damping
regime, where the equivalent width in frequency does not depend on $b$.

\paragraph{Transmission and curve of growth.} The rate at which a molecule
behind the column $N_l$ absorbs in the line, relative to the unshielded rate
for a continuum flat across the line, and its integral over $\tau_0$ are
\begin{align}
  S(\tau_0,a) &= \frac{1}{\sqrt\pi}\int_{-\infty}^{\infty} H(a,x)\,
       \exp\!\left[-\tau_0\frac{H(a,x)}{H(a,0)}\right]\dd x,
  \label{eq:SW}\\
  W(\tau_0,a) &= \int_0^{\tau_0} S\,\dd\tau
      = \frac{2H(a,0)}{\sqrt\pi}\int_0^{\infty}\left\{1-\exp\!\left[-\tau_0\frac{H(a,x)}{H(a,0)}\right]\right\}\dd x ,
  \label{eq:W}
\end{align}
with $x=(\nu-\nu_{lu})/\Delta\nu_D$ and $H(a,x)=(a/\pi)\int e^{-y^2}[(x-y)^2+a^2]^{-1}\dd y$
the Voigt function ($\int H\,\dd x=\sqrt\pi$), evaluated with the rational
approximation of \citet{Humlicek1982} (relative accuracy $10^{-4}$; \rtn{h2Voigt},
\src{h2\_mod.f90:\allowbreak{}664--686}). $S$ is the single-line self-shielding
function: the profile-weighted transmission, $S(0,a)=1$. $W$ is the equivalent
width in units of $\sqrt\pi\Delta\nu_D/H(a,0)$, normalised so that $W\to\tau_0$ for
$\tau_0\ll1$, and $\dd W/\dd\tau_0=S$. The three regimes are
\begin{equation}
  S \simeq \begin{cases}
    1-\tau_0/\sqrt2 & (\tau_0\ll1),\\[2pt]
    \bigl[\sqrt\pi\,\tau_0\,(\ln\tau_0)^{1/2}\bigr]^{-1} & \text{(saturated Doppler core)},\\[2pt]
    \pi^{-1/4}\,\bigl[H(a,0)\,a/\tau_0\bigr]^{1/2} & \text{(damping wings)},
  \end{cases}
  \label{eq:Sasym}
\end{equation}
the first two for $a\to0$, the third when the absorption is dominated by the
Lorentzian wings $H\simeq a/(\sqrt\pi x^2)$.

\paragraph{Profile normalisation.} $\tau_0=\sigma_{lu}N_l$ is the line-centre
optical depth of the Doppler core, and $S$ normalises the profile to its own
centre, $\tau(x)=\tau_0H(a,x)/H(a,0)$; with this convention $\tau_0$ agrees with
Cloudy's line optical depths to 0.04\% (\dlog{H-005}). The area-normalised Voigt
profile gives $\tau(x)=\tau_0H(a,x)$, smaller by the factor
$H(a,0)\simeq1-2a/\sqrt\pi$ at every frequency. For these data the difference is
at most 3\% in $\tau$ (0.8\% weighted by the unshielded pumping, and half that in
$S$ in the damping regime), negligible against the other uncertainties.

\paragraph{Numerical strategy.} \rtn{shieldQuad} (\src{:\allowbreak{}692--746})
integrates one side of the profile with the trapezoidal rule on 1200 uniform
points to $x=6$ and logarithmically (60 points per decade) beyond, to
$x_{\max}=\max(60,\,300\sqrt{\tau_0 a})$; beyond $x_{\max}$ the Lorentzian wing
$H\simeq a/(\sqrt\pi x^2)$ is optically thin ($\tau\le10^{-5}$ at $x_{\max}=300\sqrt{\tau_0a}$)
and is added analytically. It also records the offsets $x_q$ at which $1/6$,
$1/2$ and $5/6$ of one side's absorption, $\int H e^{-\tau}\dd x$, has accumulated,
for the overlap quadrature. Against an independent quadrature (the Faddeeva
function of SciPy) $S$ and $W$ agree to 0.04\% and 0.03\% over
$\tau_0=10^{-2}$--$10^{11}$, $a=10^{-6}$--$10^{-1}$ (\dlog{H-005}). \rtn{shieldTable}
(\src{:\allowbreak{}748--760}) tabulates $\log S$, $\log W$ and $\log x_q$ once on a
$161\times29$ grid, $\log_{10}\tau_0\in[-4,16]$ (step 0.125) and
$\log_{10}a\in[-7,0]$ (step 0.25). Look-up (\rtn{tabShieldJ}, \src{:\allowbreak{}833--851})
is bilinear in $(\log\tau_0,\log a)$ on the logarithms; below the table $S=1$ and
$W=\tau_0$; above it the damping-wing asymptote, $S\propto\tau_0^{-1/2}$,
$W\propto\tau_0^{1/2}$. The table's $S$ is within 1.2\% of the exact single-line
integral (\dlog{H-006}). The damping parameter's column and fraction in the table
are found once per line (\rtn{h2ShieldAIndex}, \src{:\allowbreak{}822--830}).

\paragraph{Comparison with shielding functions.} A shielding-function
treatment writes the Solomon rate as $k_{\mathrm{Sol}}\simeq k_{\mathrm{Sol}}(0)\,
f_{\mathrm{shield}}(N_{\mathrm{H_2}})\,e^{-\tau_{d,1000}}$ \citep[their
eq.~40]{DraineBertoldi1996}, with $\tau_{d,1000}$ the dust optical depth at
1000\,\AA{} and $f_{\mathrm{shield}}$ a fit to multi-line models for given
populations, temperature and field shape. The fit of \citet[their
eq.~37]{DraineBertoldi1996},
\begin{equation}
  \begin{split}
  f_{\mathrm{shield}} &= \frac{0.965}{(1+x/b_5)^2}+\frac{0.035}{(1+x)^{1/2}}\,
     \exp\!\left[-8.5\times10^{-4}(1+x)^{1/2}\right],\\
  x&=\frac{N(\mathrm{H_2})}{5\times10^{14}\,\mathrm{cm^{-2}}},\qquad b_5=\frac{b}{\mathrm{km\,s^{-1}}},
  \end{split}
  \label{eq:db96}
\end{equation}
describes the saturation of the lines in its first two terms and their
overlap in the exponential. 3D-PDR's own rate uses the single-line function of
\citet{Federman1979} for one representative line
(Section~\ref{sec:h2rate}). Here the equivalent of $f_{\mathrm{shield}}$
follows from the level columns, the line widths, the populations and the
field at each line, and the dust attenuation is in the MC field. On the Cloudy
slab of \dlog{H-005} (100\,K, $n_{\mathrm H}=10^4$\,cm$^{-3}$, 100 Habing,
$b=0.91$\,km\,s$^{-1}$) the model without overlap followed Equation~(\ref{eq:db96})
without its overlap factor to 0.79--1.00 at $N(\mathrm{H_2})=10^{18}$--$10^{21}$\,cm$^{-2}$,
and with overlap it lies at 0.70--0.95 of the full fit there and 20--50\% above
it at $10^{14}$--$10^{17}$\,cm$^{-2}$ (\dlog{H-005}, \dlog{H-006}). The analytic
theory of the H\,\textsc{i}-to-\Htwo{} transition of \citet{Sternberg2014}, in
which the total atomic column of a slab in a beamed field is
$N_{\mathrm{H\,I}}=\sigma_g^{-1}\ln(\alpha G/2+1)$ (their eq.~40), with $\sigma_g$ the
far-ultraviolet dust absorption cross section per H nucleus, $\alpha$ the ratio
of the free-space dissociation rate to the formation rate $R_{\mathrm f}n$ and $G$ a
cloud-averaged self-shielding factor, rests on the same kind of line-by-line
computation (an ``effective dissociation bandwidth'' computed with the Meudon
code).

\subsection{Averaging through the cell's own column}
\label{sec:ownavg}

The pumping rate that the level balance of a cell needs is its mean over the
cell, whose populations are uniform. If the cell adds $\Delta\tau=\sigma_{lu}\Delta N_l$
to the column $\tau_0$ in front of it, the mean of $S$ over a uniform layer
traversed along one chord is the mean slope of the curve of growth across it
(\rtn{h2Shield}, \rtn{h2ShieldJ}, \src{h2\_mod.f90:\allowbreak{}768--784, 854--866};
\rtn{h2LineShieldPrep/Eval}, \src{:\allowbreak{}876--907}):
\begin{equation}
  \langle S\rangle = \frac{W(\tau_0+\Delta\tau,a)-W(\tau_0,a)}{\Delta\tau}
  \quad\text{if }\Delta\tau > 10^{-3}\max(\tau_0,10^{-3}),
  \qquad\text{else } S(\tau_0+\tfrac12\Delta\tau,a),
  \label{eq:cellavg}
\end{equation}
clamped to $[0,1]$; the second form avoids the cancellation in the difference
of $W$ for thin layers, with an error of order $\Delta\tau^2 S''/24$. A cell whose own
column takes its lines from thin to saturated absorbs the band in its front
part, and $\langle S\rangle$ is then much smaller than $S(\tau_0)$. The cell's own
column of level $l$ is $\Delta N_l = n_l\,n(\mathrm{H_2})\,\ell$ with two path
lengths (\src{pdr\_mod.f90:\allowbreak{}5609--5615, 5697--5706, 5804--5819}):
\begin{align}
  \ell_{\mathrm{dir}} &= \frac{\Delta x\,\Delta y\,\Delta z}{\Delta y\Delta z\,|u_x|+\Delta x\Delta z\,|u_y|+\Delta x\Delta y\,|u_z|},\\
  \ell_{\mathrm{iso}} &= \frac{4V}{A} = \frac{2\,\Delta x\,\Delta y\,\Delta z}{\Delta x\Delta y+\Delta y\Delta z+\Delta x\Delta z},
\end{align}
where $\ell_{\mathrm{dir}}$ is the mean chord of the cell along the unit vector
$\hat u$ towards the source (the volume over the area projected on a plane
normal to $\hat u$), averaged over sources with weights $r_s^{-2}$ and used for
the direct beam, and $\ell_{\mathrm{iso}}$ is Cauchy's mean chord of a convex body
for isotropically distributed lines, used for the diffuse light. The widths
$\Delta x,\Delta y,\Delta z$ are centred differences of the axis values
(\rtn{cWid}), one-sided at the first and last axis value; there the full
spacing is used, although the outermost cell of a grid that ends at its last
axis value is half as wide. One representative chord stands for the
distribution of chord lengths through the cell. The direct-path column uses
the populations of the current pass of the cell's iteration, the diffuse one
those of the previous sweep.

This is how the level model includes the self-shielding of the cell's own
\Htwo, which the chemistry's columns leave out (Section~\ref{sec:avradial}).
Averaging the transmission through the cell is a different operation from
adding the cell's near half to the shielding column, which made the
chemistry's surface cell shield itself completely and settle on a
self-shielded branch (\dlog{D-088}). The size of the effect shows how thin the
pumping layer is: when the diffuse part was first computed without the cell's
own column, the Solomon rate where \Htwo{} turns on was 4--12 times the
chemistry's, and averaging through the mean chord lowered the 1--0\,S(1)
luminosity of the test model by a factor 4.6 (\dlog{H-007}). The
average is exact for a uniform cell; cells across the H/\Htwo{} transition are
not, and on \texttt{hiipdr} the 1--0\,S(1) luminosity changed by $+35\%$ from
$21^3$ to $31^3$ and by $+3\%$ from $31^3$ to $41^3$ cells while the \Htwo{} profile
of the transition was still moving (\dlog{H-007}).

\subsection{Line overlap}
\label{sec:overlap}

The \Htwo{} lines shade one another. Each line's transmission is its own $S$
times the mean transmission of all the \emph{other} lines over the frequencies
where it absorbs (\rtn{lineOverlap}, \src{h2\_mod.f90:\allowbreak{}1018--1289},
\dlog{H-006}):
\begin{equation}
  \mathcal{O}_i = \frac16\sum_{k=1}^{3}\sum_{s=\pm1}
     \exp\!\Bigl[-\tau_{\mathrm{oth}}\bigl(\tilde\nu_i + s\,x_{q,k}\,\Delta\tilde\nu_{D,i}\bigr)\Bigr],
  \qquad
  \tau_{\mathrm{oth}}(\tilde\nu) = \sum_{j\neq i,\ \tau_{0,j}\ge10^{-2}} \tau_{0,j}\,\frac{H(a_j,x_j)}{H(a_j,0)},
  \label{eq:overlap}
\end{equation}
with $x_j=(\tilde\nu-\tilde\nu_j)/\Delta\tilde\nu_{D,j}$ and
$\Delta\tilde\nu_{D,j}=\tilde\nu_j b/c$ [cm$^{-1}$]. $\mathcal{O}_i$ approximates
$\int\phi_ie^{-\tau_i}e^{-\tau_{\mathrm{oth}}}\dd\nu\big/\int\phi_ie^{-\tau_i}\dd\nu$, the mean of
$e^{-\tau_{\mathrm{oth}}}$ weighted by line $i$'s own absorption, by a midpoint rule in
the cumulative absorption: each side of the line is cut into three parts that
carry equal shares of its absorption, and each part is represented by the
offset $x_{q,k}$ that halves its share ($1/6$, $1/2$ and $5/6$ of the side). A thin
line thus samples its core and a damped line its wings, hundreds of Doppler
widths out. Three points per side were chosen by test: against an exact
reference the factor on the Solomon rate is 0.86--1.03 with one point per side,
0.99--1.02 with two and 0.98--1.04 with three, and the line-by-line rms falls from
0.03--0.45 to 0.03--0.21 (\dlog{H-006}). Numerically:
\begin{itemize}
\item only lines with $\tau_0\ge10^{-2}$ shade; a shading line's reach is
  $\Delta\tilde\nu_D\max\bigl(12,\sqrt{\tau_0 a/(\sqrt\pi H(a,0)\,\varepsilon)}\bigr)$
  with $\varepsilon=10^{-4}$, the offset where its wing optical depth falls
  to $10^{-4}$;
\item the full Voigt function is used within 12 Doppler widths and the
  Lorentzian wing $\tau_0 a/(\sqrt\pi H(a,0)x^2)$ beyond;
\item lines whose reach exceeds 100\,cm$^{-1}$ (``long'') shade every line
  through their Lorentzian wings, summed as $\sum_j C_j/(\tilde\nu-\tilde\nu_j)^2$
  with four interleaved partial sums; other lines (``short'') are added where
  they reach, within 100\,cm$^{-1}$; lines within 40 Doppler widths of line $i$
  (``near'') are evaluated term by term;
\item the lines are processed in wavenumber order so that the windows only
  move forward; a thin line ($\tau_0<1$) takes the far contribution once at
  its centre; lines that no shading line reaches keep $\mathcal{O}=1$.
\end{itemize}
In \rtn{solveCell} the factor is computed once per cell and path, at the columns
in front of the cell plus half the cell's own (Section~\ref{sec:solvecell}), and
multiplies the cell-averaged $\langle S\rangle$ of Equation~(\ref{eq:cellavg}): a product of
means stands for the mean of the product, adequate where the other lines'
shading changes little across the cell. The result was checked against the
same integral on a 0.05\,cm$^{-1}$ grid across 84\,000--119\,500\,cm$^{-1}$ (every line
with $\tau_0>10^{-3}$ a Voigt profile). The exact factor on the Solomon rate of the
\dlog{H-005} slab's ground-state levels is 0.96, 0.93, 0.86, 0.65, 0.49 and 0.32 at
$N(\mathrm{H_2})=10^{17}$, $10^{18}$, $10^{19}$, $10^{20}$, $3\times10^{20}$ and
$10^{21}$\,cm$^{-2}$; the model reproduces it to 0.2--1.6\% up to
$3\times10^{20}$\,cm$^{-2}$ and to 4.5\% at $10^{21}$\,cm$^{-2}$, and the exact factor is
within 5\% of the overlap term of Equation~(\ref{eq:db96}) (\dlog{H-006}).
Overlap lowered the total pumping through that slab by 33\%.

Not included: a photon absorbed in another line's wing is counted as lost
rather than as a pump of that line (cross-pumping), which underestimates the
pumping of the weaker lines where the band is saturated; and the
H\,\textsc{i} Lyman-series lines, which the MC field does not contain either
(Section~\ref{sec:eta}), do not shade the band. Their damping wings remove
continuum around Ly$\beta$, Ly$\gamma$ and Ly$\delta$ (1026, 973 and 950\,\AA) over
equivalent widths of about 4, 1.4 and 0.7\,\AA{} at
$N(\mathrm{H\,I})=10^{21}$\,cm$^{-2}$, growing as $N(\mathrm{H\,I})^{1/2}$ (an
estimate of this document from the damping-wing curve of growth, i.e.\ a few
per cent of the 912--1110\,\AA{} band); pumping and dissociation behind such
columns are overestimated (Section~\ref{sec:limitations}). The Meudon code
includes them in its exact transfer \citep{LePetit2006}.

\subsection{Escape probabilities}
\label{sec:escape}

The probability that a photon emitted in line $t$ ($u\to l$) at the cell
centre leaves the cloud is (\rtn{h2LineEscape},
\src{h2\_mod.f90:\allowbreak{}1344--1361}; \src{pdr\_mod.f90:\allowbreak{}5821--5835})
\begin{equation}
  \beta_t = \frac{1}{12}\sum_{k=1}^{12} S\bigl(\sigma_t N_{l,k},\,a_t\bigr),
  \label{eq:beta}
\end{equation}
with $N_{l,k}$ the column of the line's lower level from the cell centre to
the edge of the model along line ray $k$ (Section~\ref{sec:lineRays}), the
cell's own half included, from the populations of the previous sweep. With
complete redistribution the photon is emitted with the absorption profile,
so the probability that it crosses a column in one direction,
$\int\phi\,e^{-\tau}\dd x/\int\phi\,\dd x$, is the function $S$ of
Equation~(\ref{eq:SW}); the mean over the twelve HEALPix directions is the mean
over the sphere. The medium is static (no velocity field, so no Sobolev-type
escape), the escape probability is that of the cell centre rather than a mean
over the cell, and each line is treated alone (no overlap in the trapping).
The Meudon code makes the same on-the-spot approximation for the ultraviolet
photons of the \Htwo{} cascades \citep{LePetit2006}.

\subsection{Formation on grains: the nascent distribution}
\label{sec:formdist}

Molecules formed on grains are placed in the $(v,J)$ distribution of
\citet{TakahashiUehara2001}, as Cloudy applies it \citep{Shaw2005}
(\rtn{h2FormDist}, \src{h2\_mod.f90:\allowbreak{}589--639}). With $D$ the dissociation
energy of $X$ from its $(0,0)$ level (36\,118.11\,cm$^{-1}$, 4.478\,eV; their
$D_e=4.48$\,eV), $E_{\mathrm{off}}=0.273$\,eV the energy of $(0,0)$ above the bottom of the
potential well (their zero-point energy $E_{v0}=0.27$\,eV; 0.273\,eV is Cloudy's
value), $E(v,J)$ the level energies above $(0,0)$, and per surface the
parameters of Table~\ref{tab:tu01}:
\begin{gather}
  E_{vm} = X_{\mathrm{vib}}D + E_{\mathrm{off}},\qquad
  E_v = E(v,0)+E_{\mathrm{off}},\\
  E_{\mathrm{dust}}(v) = D\,X_{\mathrm{dust}}\Bigl[1-\frac{E_v-E_{vm}}{D+E_{\mathrm{off}}-E_{vm}}\,\frac{1-X_{\mathrm{dust}}}{2}\Bigr],\qquad
  E_{\mathrm{H_2}}(v) = D+E_{\mathrm{off}}-E_{\mathrm{dust}}(v),\\
  F(v) = e^{-x_v^2},\quad
  x_v = \begin{cases}\dfrac{E_v-E_{vm}}{g_1E_{vm}} & E_v\le E_{vm}\\[6pt]
                     \dfrac{E_v-E_{vm}}{g_2\,(E_{\mathrm{H_2}}(v)-E_{vm})} & E_v>E_{vm}\end{cases},\\
  E_{\mathrm{rot}}(v) = \bigl(E_{\mathrm{H_2}}(v)-E_v\bigr)\frac{X_{\mathrm{rot}}}{X_{\mathrm{rot}}+X_{\mathrm{trans}}},\qquad
  \phi(v,J)\propto g_{\mathrm{ns}}\,\frac{F(v)}{\sum_vF}\,(2J+1)\,\tfrac12\Bigl[e^{-y^2}+e^{-\Delta E/E_{\mathrm{rot}}}\Bigr],
\end{gather}
with $\Delta E=E(v,J)-E(v,0)$, $y=(\Delta E-E_{\mathrm{rot}})/(E_{\mathrm{rot}}/2)$,
$g_{\mathrm{ns}}=3$ (ortho) or 1 (para), $\phi=0$ where $E_{\mathrm{rot}}\le0$ (high $v$;
J.~Takahashi's prescription as quoted in Cloudy's source), and $\phi$
normalised to one. These are equations~1--7 of \citet{TakahashiUehara2001}:
the energy left on the grain $E_{\mathrm{dust}}(v)$ (their eq.~1), the energy that
stays in the molecule (eq.~2) and its share in rotation (eq.~3), and a
vibrational distribution peaked at $E_{vm}$ with different widths below and
above it (eqs.~5--7). The rotational distribution is the mean of their two
models, a Gaussian about $E_{\mathrm{rot}}$ of width $E_{\mathrm{rot}}/2$ (model A,
eq.~12) and a thermal distribution at $E_{\mathrm{rot}}$ (model B, eq.~13); the mean
is applied before one overall normalisation, as in Cloudy. Carbonaceous,
silicate and icy surfaces put the new molecules mainly near $v=2$, 6 and 7
respectively \citep{Shaw2005}. The fraction formed above 4100\,cm$^{-1}$ agrees
with Cloudy's network to four digits (0.9771, \dlog{H-004}).

The energy budget of a formation follows by conservation: $E_{\mathrm{grain}}=\langle
E_{\mathrm{dust}}\rangle_\phi$ is left on the grain,
$E_{\mathrm{kin}}=D-E_{\mathrm{grain}}-\langle E(v,J)\rangle_\phi$ is the translational energy
of the departing molecule, which heats the gas directly, and the internal
energy $\langle E(v,J)\rangle_\phi$ is either radiated in \Htwo{} lines or returned to
the gas by collisional de-excitation, i.e.\ through
$\mathcal{H}_{\mathrm{coll}}$ (Section~\ref{sec:solvecell}). $E_{\mathrm{kin}}$ is 0.834\,eV for
silicate and 1.142\,eV for carbonaceous surfaces, against 3D-PDR's 1.5\,eV
(\dlog{H-013}). $E_{\mathrm{grain}}$ is a diagnostic only (\src{h2Cells.out}); it is not
added to MOCASSIN's grain energy balance.

\begin{table}[H]
\centering
\caption{\citet[their Table~2]{TakahashiUehara2001} parameters as used
(\src{h2\_mod.f90:\allowbreak{}132--139}); the ice set is not used by \rtn{solveCell}.}\label{tab:tu01}
\begin{tabular}{@{}lcccccc@{}}
\toprule
Surface & $X_{\mathrm{vib}}$ & $X_{\mathrm{dust}}$ & $X_{\mathrm{rot}}$ & $X_{\mathrm{trans}}$ & $g_1$ & $g_2$ \\
\midrule
ice & 0.70 & 0.04 & 0.14 & 0.12 & 0.3 & 0.6 \\
silicate & 0.60 & 0.10 & 0.15 & 0.15 & 0.4 & 0.6 \\
carbonaceous & 0.20 & 0.40 & 0.15 & 0.25 & 0.9 & 0.4 \\
\bottomrule
\end{tabular}
\end{table}

In each cell the silicate and carbonaceous distributions are mixed in the
proportion of the chemistry's own formation rate on each material:
$w_{\mathrm{sil}}$ and $w_{\mathrm{gra}}$ are \rtn{H2\_FORMATION\_RATE}
(Section~\ref{sec:formation}) evaluated with only that material's grain area
(\src{pdr\_mod.f90:\allowbreak{}5839--5876}). A cell with no grain formation (no dust,
or $T_d=0$) takes a Boltzmann distribution at $T_e$, i.e.\ no formation pumping,
and is counted. Formation through H$^-$ and \Htwo$^+$, which populates specific
levels, is not distributed over the levels: every molecule the closure forms
enters in $\phi$.

\subsection{The per-cell algorithm}
\label{sec:solvecell}

\rtn{solveCell} (\src{pdr\_mod.f90:\allowbreak{}5665--5954}) does, for one cell:
\begin{enumerate}
\item \textbf{Field} (:5667--5695): $b$ from $T_e$; $J_\nu^{\mathrm{dir}}$,
  $J_\nu^{\mathrm{dif}}$ per bin and $\eta^{\mathrm{dir}}$,
  $\eta^{\mathrm{dif}}$ per line (Equation~\ref{eq:eta}).
\item \textbf{Cell geometry} (:5697--5706): cell widths, $\ell_{\mathrm{iso}}$,
  and the cell's own column for the diffuse paths from its populations of the
  last sweep, $\Delta N^{\mathrm{iso}}_l = n_l n(\mathrm{H_2})\ell_{\mathrm{iso}}$.
\item \textbf{Diffuse shielding} (:5712--5802), if any line has
  $\eta^{\mathrm{dif}}>0$: \rtn{traceDiffuseSites} gives the sites and their
  weights $w$ (Section~\ref{sec:sites}); walking the stored rays again, each
  site gets a column per level,
  \begin{equation}
    N_{l,\mathrm{site}} = N^\star_l(P) + n_l(P)\,n(\mathrm{H_2},P)\,\ell_{\mathrm{own}}(P)
        + \sum_{\mathrm{between}} n_l\,n(\mathrm{H_2})\,\ell + \tfrac12\,n_l(P)\,n(\mathrm{H_2},P)\,\ell_P
  \end{equation}
  (the level-resolved form of Equation~\ref{eq:twoleg}, with the source-ward
  level columns $N^\star_l$ of \rtn{starColumns}). Evaluating $S$ for every
  line at every site would cost $32\,453\times N_{\mathrm{site}}$ look-ups,
  and a cell has hundreds of sites, so the sites are binned:
  \begin{itemize}
  \item four groups by total column $\sum_l N_{l,\mathrm{site}}$, with edges
    $10^{17}$, $10^{19}$ and $3\times10^{20}$\,cm$^{-2}$;
  \item within a group and for each level, a histogram of the weights in
    $\log_{10}N_{l,\mathrm{site}}$, 0.1\,dex bins from $10^9$\,cm$^{-2}$
    (150 bins, the end bins catching the columns beyond), plus a ``thin'' weight for sites where
    $\sigma_{\max,l}N_l<10^{-4}$, $\sigma_{\max,l}$ being the largest
    cross section of the lines out of $l$;
  \item the overlap factor (Equation~\ref{eq:overlap}) once per group with
    more than $10^{-4}$ of the total weight, at the group's weight-mean
    level columns plus $\frac12\Delta N^{\mathrm{iso}}$ (none for the
    $<10^{17}$\,cm$^{-2}$ group);
  \item per line, the diffuse shielding factor
    \begin{equation}
      \langle S\rangle^{\mathrm{dif}}_i = \frac{1}{\sum w}\sum_{g}\mathcal{O}_{i,g}
      \Bigl[w^{\mathrm{thin}}_{l,g}\,\langle S\rangle(0) + \sum_{\mathrm{bins}} w_{l,g,\mathrm{bin}}\,
      \langle S\rangle\bigl(\sigma_i N_{\mathrm{bin}}\bigr)\Bigr],
    \end{equation}
    each $\langle S\rangle$ averaged through the cell's own column
    $\sigma_i\Delta N^{\mathrm{iso}}_l$ by Equation~(\ref{eq:cellavg}).
  \end{itemize}
  The binning replaces each site's column by its bin centre (an error of at
  most $\pm0.05$\,dex in $N_l$) and each group's overlap by that at the group's
  mean columns. If there is diffuse light but no site (no dust on any of the
  twelve rays) the diffuse part rides the direct path (:5799--5802).
\item \textbf{Direct path} (:5804--5819): source-ward level columns
  $N^\star_l$ and the chord $\ell_{\mathrm{dir}}$.
\item \textbf{Escape} (:5821--5835): $\beta_t$ over the twelve line rays
  (Equation~\ref{eq:beta}).
\item \textbf{Formation} (:5839--5876): mixture weights, $\phi$, $E_{\mathrm{kin}}$,
  $E_{\mathrm{grain}}$.
\item \textbf{Solve} (:5879--5914). \rtn{h2LineShieldPrep} stores, for the
  direct column, each line's $\tau_0$, $W(\tau_0)$ and damping index. Then up
  to 12 passes: colliders from the current ortho fraction; the cell's own
  direct column $\Delta N^{\mathrm{dir}}_l = n_l n(\mathrm{H_2})\ell_{\mathrm{dir}}$
  from the current populations; on the first pass only, the direct overlap factor at
  $N^\star+\frac12\Delta N^{\mathrm{dir}}$; the direct shielding
  $\langle S\rangle^{\mathrm{dir}}_i\mathcal{O}^{\mathrm{dir}}_i$
  (\rtn{h2LineShieldEval}); the shielded occupation numbers
  \begin{equation}
    \eta_i = \eta^{\mathrm{dir}}_i\,\langle S\rangle^{\mathrm{dir}}_i\mathcal{O}^{\mathrm{dir}}_i
           + \eta^{\mathrm{dif}}_i\,\langle S\rangle^{\mathrm{dif}}_i
  \end{equation}
  (or $(\eta^{\mathrm{dir}}_i+\eta^{\mathrm{dif}}_i)$ times the direct factor
  when there is no site); and the level balance (\rtn{h2Solve}). The loop stops
  when the ortho fraction changes by less than $10^{-4}$; a cell that had
  settled in the previous sweep needs one pass, and the mean was 1.0--1.6 passes
  per cell (\dlog{H-011}).
\item \textbf{Products} (:5917--5952): $k_{\mathrm{Sol}}=\sum n_lk_l$;
  the absorption rate $\sum n_l\sum_uP_{lu}$; for comparison the
  chemistry's own single-line rate on the same field and columns
  (\rtn{h2SolChem}, 3D-PDR's \rtn{H2PDRATE}); diagnostics; and three
  heating terms per molecule: the atoms' kinetic energy from dissociation
  $\sum_l n_l H^{\mathrm{dis}}_l$, the net heating by collisions within $X$
  \begin{equation}
    \mathcal{H}_{\mathrm{coll}} = hc\sum_c\sum_{u>l} n_c k^c_{ul}
       \Bigl(n_u - n_l\frac{g_u}{g_l}e^{-hc(E_u-E_l)/kT}\Bigr)(E_u-E_l)
       \quad[\mathrm{erg\,s^{-1}}]
    \label{eq:hcoll}
  \end{equation}
  (\rtn{h2CollHeat}, \src{h2\_mod.f90:\allowbreak{}1425--1448}; Cloudy's
  \texttt{HeatDexc}), and the kinetic energy per molecule formed,
  $E_{\mathrm{kin}}$. $\mathcal{H}_{\mathrm{coll}}$ is the energy that collisions take out
  of the levels populated above their Boltzmann ratio at $T$ by pumping and
  formation; it is negative, i.e.\ \Htwo{} line cooling, where the gas is
  hotter than the level populations. A pump deposits in the gas only what
  $\mathcal{H}_{\mathrm{coll}}$ recovers from the levels it feeds, and a dissociation
  the kinetic energy of the atoms, so the three terms do not count any
  energy twice.
\end{enumerate}

\subsection{Sweeps, parallel strategy and convergence}
\label{sec:sweeps}

\rtn{h2LevelModel} reads the data once and, on every call:
\begin{itemize}
\item takes $n(\mathrm{H_2})=n_{\mathrm H}x(\mathrm{H_2})$ of every cell
  from the chemistry; on the first call it starts every cell from Boltzmann
  populations at $T_e$, or, on a warm start, from \src{output/\allowbreak{}h2Pop.bin}
  if that file describes the same grid and level set (\rtn{h2PopResume},
  \src{pdr\_mod.f90:\allowbreak{}5961--6012}, \dlog{H-020}); later calls start from the last
  call's populations. A cold start on a converged model overshoots the
  dissociations of its first sweep by 38\% (\dlog{H-016}); a resumed one is
  within 0.6\% of the checkpointed run (\dlog{H-020});
\item for each sweep, rebuilds the source-ward level columns of every cell,
  $N^\star_l=\sum n_l\,n(\mathrm{H_2})\,\ell$ along \rtn{rayCells} to each
  source, weighted over sources by $r_s^{-2}$ and excluding the cell itself
  (\rtn{starColumns}, \src{:\allowbreak{}5621--5659}); the other cells' populations in all
  columns are those of the previous sweep, so a sweep is a Jacobi iteration
  over the cells, whose result does not depend on the order in which the
  cells are solved;
\item distributes the cells dynamically: the ranks take cells one at a time
  from a shared counter in an MPI-3 one-sided window allocated on rank~0
  (\verb|MPI_Win_allocate|, \verb|MPI_Fetch_and_op|, \src{:\allowbreak{}5398--5426}); the
  first sweep takes them in a fixed pseudo-random order (a Fisher--Yates shuffle
  driven by a linear congruential generator), later sweeps in decreasing cost
  of the previous sweep (\rtn{orderCells}, \src{:\allowbreak{}5538--5559}), so that cores of
  different speed finish together; a cell's result does not depend on the rank
  that solves it;
\item reduces populations, rates and heating terms with \verb|MPI_Allreduce|
  (in single precision);
\item stops when the luminosities of 1--0\,S(1), 2--1\,S(1), 0--0\,S(0) and
  0--0\,S(1), summed over the grid in the optically thin limit, all change by
  less than 1\% between sweeps (\src{:\allowbreak{}5459--5495}), at most \rtn{nSweepMax}
  sweeps.
\end{itemize}
Inside the MC iteration \rtn{nSweepMax}$=1$: one sweep per iteration, the
level model converging together with the rest of the model. On the finished
model \rtn{nSweepMax}$=12$ (Section~\ref{sec:convergence}). On \texttt{pn150k}
($41^3$, 35\,384 cells, 14 ranks) a call took 3736\,s before and 1474\,s after the
optimisations of \dlog{H-011}; the overlap calculation was about half of the
timed work (\dlog{H-007}).

% ======================================================================
\section{Reaction rates handed to the chemistry}
\label{sec:photorates}
% ======================================================================

The chemistry's rate coefficients are computed per cell in \rtn{pdrStep}
(\src{pdr\_mod.f90:\allowbreak{}3180--4580}) by 3D-PDR's
\rtn{CALCULATE\_REACTION\_RATES} (\src{pdr/\allowbreak{}calc\_reac\_rates.F90}), called at
\src{pdr\_mod.f90:\allowbreak{}3692} with MOCASSIN's inputs, and then partly overwritten
with rates that MOCASSIN or the level model own. Table~\ref{tab:rateorigin}
summarises where each radiation-driven rate comes from; reaction numbers are
those of \src{dataPDR/\allowbreak{}rates\_reduced.d}.

\begin{longtable}{@{}L{0.28\linewidth}L{0.15\linewidth}L{0.37\linewidth}L{0.09\linewidth}@{}}
\caption{Origin of the radiation- and cosmic-ray-driven rates in the 3D build.}\label{tab:rateorigin}\\
\toprule
Reaction(s) & Number(s) & Rate used & Section \\
\midrule
\endfirsthead
\toprule
Reaction(s) & Number(s) & Rate used & Section \\
\midrule
\endhead
\Htwo\,+\,$h\nu\to$\,2H & 329 & the level model's $k_{\mathrm{Sol}}$ & \ref{sec:h2rate} \\
CO\,+\,$h\nu\to$\,C\,+\,O & 290 & $\alpha_{\mathrm{CO}}\sum_r\chi_rf_r\Theta$ \citep{vanDishoeckBlack1988} & \ref{sec:co} \\
24 other photoreactions & 266--268, 270--288, 291--292 & $\alpha\chi$ (Draine spectral shape) & \ref{sec:generic} \\
C, Mg\,+\,$h\nu\to$\,ion\,+\,$e^-$ & 269, 289 & set to zero & \ref{sec:linked} \\
ionisation of H, He, C, O, Mg & 293, 297, 298, 299, 312 & cosmic-ray (or, Mg, induced-photon) rate plus MOCASSIN's photoionisation rate & \ref{sec:linked} \\
\Htwo\,+\,CRP\,$\to$\,\Htwo$^+$\,+\,$e^-$ & 296 & cosmic-ray rate plus MOCASSIN's \Htwo{} photoionisation & \ref{sec:linked} \\
X$^+$\,+\,$e^-\to$\,X\,+\,$h\nu$ (X = H, He, C, O, Mg) & 242--246, 248, 249 & MOCASSIN's total recombination coefficient & \ref{sec:linked} \\
cosmic-ray ionisation & 293--300 & $\alpha\,\zeta/\zeta_0$ & \ref{sec:cr} \\
cosmic-ray-induced photoreactions & 301--315 & $\alpha(\zeta/\zeta_0)(T/300)^\beta\gamma/(1-\omega)$ & \ref{sec:cr} \\
charge transfer with H, \Hp & 156--157, 161--163, 167 & \rtn{pdrCtCoef} & \ref{sec:ct} \\
H\,+\,H\,(grain)\,$\to$\,\Htwo & 328 & \citet{CazauxTielens2002} & \ref{sec:formation} \\
\bottomrule
\end{longtable}

\subsection{The field as a set of rays}
\label{sec:rays}

3D-PDR evaluates each photoreaction as a sum over rays, each ray $r$ carrying
a field strength $\chi_r$ (in Draine units), a visual extinction $A_{V,r}$ and
shielding columns. In the 3D build the bridge builds
$n_{\mathrm{chem}} = 1 + N_{\mathrm{site}}$ (+1 fallback) rays
(\src{pdr\_mod.f90:\allowbreak{}3383--3455}, \dlog{D-348}, \dlog{D-350}), which
represent the local MC field split by the paths its photons took
(Sections~\ref{sec:split} and \ref{sec:sites}):

\begin{longtable}{@{}L{0.17\linewidth}L{0.20\linewidth}L{0.10\linewidth}L{0.17\linewidth}L{0.24\linewidth}@{}}
\caption{The rays handed to the chemistry in the 3D build.}\label{tab:rays}\\
\toprule
Ray & $\chi_r$ & $A_{V,r}$ & LW factor & Columns $N(\mathrm{H_2}),N(\mathrm{CO})$ \\
\midrule
\endfirsthead
\toprule
Ray & $\chi_r$ & $A_{V,r}$ & LW factor & Columns $N(\mathrm{H_2}),N(\mathrm{CO})$ \\
\midrule
\endhead
0 (direct) & $\chiD$ & 0 & $f_{\mathrm{dir}}$ & source-ward, Eq.~(\ref{eq:avradial}) \\
site $j$ & $\chiF\,w_j/\sum w$ & 0 & $f_{\mathrm{dif}}$ & two-leg, Eq.~(\ref{eq:twoleg}) \\
fallback & $\chiF$ & 0 & $f_{\mathrm{dif}}$ & source-ward (only when $\chiF>0$ and no site exists) \\
\bottomrule
\end{longtable}

$A_{V,r}=0$ on every ray (\src{:\allowbreak{}3439}) because the MC field has already been
attenuated by the dust; applying 3D-PDR's $e^{-\gamma A_V}$ again would count
the dust twice (\dlog{D-343}, \dlog{D-350}). The C\,\textsc{i} column handed
over is zero (\src{:\allowbreak{}3627}) because MOCASSIN's opacity already contains the
C\,\textsc{i} photoionisation continuum; the HD and S columns are zero
(neither species is in the network). If a source-ward column is missing
(negative), it falls back to $(\AV/(\AV/N_{\mathrm H}))\,x_X$
(\src{:\allowbreak{}3458}, \src{:\allowbreak{}3603--3605}); zero is a valid column and is kept
(\dlog{D-088}). The columns exclude the target cell (Section~\ref{sec:avradial});
only the level model averages the shielding through the cell's own gas
(Section~\ref{sec:ownavg}).

\subsection{Generic photoreactions}
\label{sec:generic}

For every reaction $X+\mathrm{PHOTON}\to\ldots$ without a dedicated treatment,
3D-PDR's rate is $k=\sum_r\alpha\,\chi_r\,e^{-\gamma A_{V,r}}$ [s$^{-1}$]
(\src{pdr/\allowbreak{}calc\_reac\_rates.F90:\allowbreak{}321--347}), with $\alpha$ the rate in the
unattenuated Draine field and $\gamma$ the dust-attenuation exponent of the UMIST
format; with $A_{V,r}=0$,
\begin{equation}
  k = \alpha\,(\chiD+\chiF) = \alpha\,\chi ,
  \label{eq:generic}
\end{equation}
the coefficient $\alpha$ of the rate file times the local MC field in the
6--13.6\,eV band. This applies to the photodissociation and photoionisation of
H$_2^+$, H$_3^+$, CH, CH$^+$, CH$_2$, CH$_2^+$, CH$_3$, CH$_3^+$, CH$_4$, OH,
OH$^+$, H$_2$O and O$_2$ (24 of the network's 28 photoreactions). The
coefficients $\alpha$ were computed for the Draine spectral shape, while the
MC field reddens with depth: its LW share falls from 0.30 to 0.008 through
\texttt{hiipdr}'s PDR (\dlog{D-349}). Reactions whose cross sections lie mainly
above 11.2\,eV are therefore overestimated at depth, and those driven by the
6--11\,eV part underestimated; none of these reactions is shielded by \Htwo{}
or CO lines. Rates computed from cross sections folded over the actual field
would remove the assumption; \citet{Heays2017} provide cross sections for 102
species for that purpose, and find that their rates for the interstellar field
agree with earlier compilations to about 30\% in general.

\subsection{\texorpdfstring{\Htwo}{H2} photodissociation: the level model's rate}
\label{sec:h2rate}

3D-PDR first computes its own rate on the rays (\rtn{H2PDRATE}, \rtn{H2SHIELD1};
\src{pdr/\allowbreak{}shield.F90:\allowbreak{}27--55, 231--281}; \src{pdr/\allowbreak{}calc\_reac\_rates.F90:\allowbreak{}259--268}),
\begin{equation}
  k^{\mathrm{net}}_{\mathrm{H_2}} = \sum_r \alpha_{\mathrm{H_2}}\,\chi_r f_r\,
     J_{\mathrm{FGK}}(N_{\mathrm{H_2},r})\,\rtn{SCATTER}(0,1000\,\text{\AA}),
  \qquad \alpha_{\mathrm{H_2}}=5.18\times10^{-11}\unit{s^{-1}},
\end{equation}
where $J_{\mathrm{FGK}}=J_D+J_R$ is the single-line self-shielding function of
\citet[their eqs.~A7--A9]{Federman1979} for one representative line at
1000\,\AA{} with oscillator strength 0.01, half the \Htwo{} in its lower level, a
Doppler width $v_{\mathrm{turb}}/\lambda$ and a radiative width of $8\times10^7$\,s$^{-1}$
(\src{pdr/\allowbreak{}shield.F90:\allowbreak{}46--51, 249--252}). The bridge then replaces it with
the level model's Solomon rate (\dlog{H-012}):
\begin{equation}
  k(\mathrm{H_2}+h\nu\to2\mathrm H) = k_{\mathrm{Sol}}
  \qquad(\src{pdr\_mod.f90:\allowbreak{}3736}).
\end{equation}
The same rate enters 3D-PDR's heating routine, so the chemistry and the
heating use one \Htwo{} photodissociation rate. The network's own rate is kept
as a diagnostic (\rtn{h2SolChem}, on the same field split and columns). On
\texttt{hiipdr} the level model's rate is 2.4--2.6 times the network's in the
ionised gas, 1.5--1.8 times in cold unshielded gas, and 0.6--0.7 and 0.4--0.6
times behind source-ward columns of $10^{15}$--$10^{17}$\,cm$^{-2}$ and more
(\dlog{H-007}). Three things separate them: the level model's rate in the
unshielded Draine field is 1.23 times $\alpha_{\mathrm{H_2}}$ and the stellar field is
harder; in warm gas the molecules sit in excited levels that dissociate
faster; and the network shields one line of fixed radiative width at the near
face of the cell, whereas the level model averages every line through the cell
and includes overlap.

\subsection{CO photodissociation}
\label{sec:co}

\rtn{COPDRATE} (\src{pdr/\allowbreak{}shield.F90:\allowbreak{}83--106}; called at
\src{pdr/\allowbreak{}calc\_reac\_rates.F90:\allowbreak{}283--292}):
\begin{equation}
  k_{\mathrm{CO}} = \alpha_{\mathrm{CO}}\sum_r \chi_r\,f_r\;
     \Theta\bigl(N_{\mathrm{CO},r},N_{\mathrm{H_2},r}\bigr)\;\rtn{SCATTER}(0,\bar\lambda),
  \qquad \alpha_{\mathrm{CO}} = 2.0\times10^{-10}\unit{s^{-1}},
\end{equation}
with $\alpha_{\mathrm{CO}}$ from the rate file (reaction 290). $\Theta$ is the combined
CO self-shielding and \Htwo{} screening factor for $^{12}$CO of
\citet[their Table~5]{vanDishoeckBlack1988}, interpolated by a two-dimensional
cubic spline in $\log_{10}\Theta$ over $(\log_{10}N_{\mathrm{CO}},\log_{10}N_{\mathrm{H_2}})$ and
clamped at the table edges (\rtn{COSHIELD}, \src{pdr/\allowbreak{}shield.F90:\allowbreak{}375--416}).
The table spans $N(\mathrm{CO})=10^{13}$--$10^{19}$ and
$N(\mathrm{H_2})=10^{19}$--$10^{23}$\,cm$^{-2}$ plus the unshielded entries, which 3D-PDR
places at $10^{12}$ and $10^{18}$\,cm$^{-2}$ (\src{pdr/\allowbreak{}modules.F90:\allowbreak{}353--356},
\src{pdr/\allowbreak{}allocations.F90:\allowbreak{}9--14}); smaller columns are treated as these.
$\bar\lambda$ is the mean wavelength of the 33 dissociating bands weighted by
their contributions to the total rate at the given columns, at which the dust
attenuation is evaluated (\rtn{LBAR}, their eq.~4,
\src{pdr/\allowbreak{}shield.F90:\allowbreak{}570--599}):
\begin{equation}
  \bar\lambda = \bigl[(5675-200.6\,w) - (571.6-24.09\,w)\,u + (18.22-0.7664\,w)\,u^2\bigr]\ \text{\AA},
  \label{eq:lbar}
\end{equation}
with $u=\log_{10}N_{\mathrm{CO}}$ and $w=\log_{10}N_{\mathrm{H_2}}$ [cm$^{-2}$],
clamped to 913.6--1076.1\,\AA. With $A_V=0$, $\rtn{SCATTER}=1$ and $\bar\lambda$ has
no effect. The window factor $f_r$ corrects for the spectral shape of the MC
field in the LW window (Section~\ref{sec:window}); CO is the one photoreaction
in this branch that uses it. The columns are those of the chemistry's own
$n(\mathrm{H_2})$ and $n(\mathrm{CO})$, on the source-ward or two-leg paths, and
exclude the target cell, so CO is under-shielded in the cells of its own
transition (Section~\ref{sec:limitations}).

The shielding functions of \citet{vanDishoeckBlack1988} were computed for
the unmodified Draine field, cool gas ($T\lesssim100$\,K, CO excitation
temperature 10\,K), Doppler parameters of 1\,km\,s$^{-1}$ for CO and 3\,km\,s$^{-1}$ for \Htwo{}
and a fixed H\,\textsc{i} column of $5\times10^{20}$\,cm$^{-2}$ (notes to their Table~5).
\citet{Visser2009} recomputed them from newer spectroscopic data and find an
unattenuated rate of $2.6\times10^{-10}$\,s$^{-1}$, 30\% above the 2.0$\times10^{-10}$ used
here, and depth-dependent rates that change by up to a factor of three when
the excitation temperature or the Doppler width is raised, for temperatures
above 100\,K. The warm CO of the illuminated PDR is therefore outside the
conditions of the table, and the \Htwo{} screening uses the chemistry's total
\Htwo{} column rather than the level model's line-by-line transmission
(Section~\ref{sec:limitations}).

\subsection{Photoionisation and recombination of the linked atoms}
\label{sec:linked}

Five elements are \emph{linked}: their neutral and first ion are species of
the network and stages of MOCASSIN (H, He, C, O, Mg;
\src{pdr\_mod.f90:\allowbreak{}553--564}). For them the network's own photoionisation is
replaced by MOCASSIN's, computed on the full MC field with the real cross
sections (\dlog{D-358}; \src{pdr\_mod.f90:\allowbreak{}3706--3722}; slots found by structure
at \src{:\allowbreak{}743--801}):
\begin{itemize}
\item the network reaction $X+\mathrm{PHOTON}\to X^++e^-$ is set to zero where it
  exists (C, reaction 269; Mg, 289; the network has none for H, He and O);
\item MOCASSIN's photoionisation rate from the ground state
  \begin{equation}
    \Gamma_X = \sum_k \frac{\Jste_k\,\sigma_X(\nu_k)}{h\nu_k}\quad[\mathrm{s^{-1}}]
    \label{eq:gammaX}
  \end{equation}
  (\src{update\_mod.f90:\allowbreak{}303--311}; outer-shell cross sections of
  \citealt{Verner1996} and \citealt{VernerYakovlev1995}, \src{ph\_mod.f90:\allowbreak{}473--476})
  is added to the reaction's first-order cosmic-ray slot $X+\mathrm{CRP}\to X^++e^-$
  (H, He, C, O: reactions 293, 297, 298, 299), on top of the cosmic rays; Mg has
  no such reaction and its cosmic-ray-induced photoionisation (312) takes the
  rate instead;
\item for hydrogen, MOCASSIN's collisional ionisation term
  $2.75\times10^{-16}\,T^{3/2}(157\,894/T+2)\,e^{-157\,894/T}$ is added as well
  (\src{update\_mod.f90:\allowbreak{}2620--2630}, stored at \src{:\allowbreak{}2859}). It has the form of a
  rate coefficient [cm$^3$\,s$^{-1}$]: at $10^4$\,K it is $6.8\times10^{-16}$, the
  magnitude of the H\,\textsc{i} collisional ionisation coefficient, but it is added
  to the photoionisation rate [s$^{-1}$] without the factor $n_e$, here and in
  MOCASSIN's own ionisation balance (\src{:\allowbreak{}2769--2770}), so collisional
  ionisation is underestimated by the factor $n_e$. With that factor it would
  reach 10\% of a photoionisation rate of $10^{-8}$\,s$^{-1}$ at about
  $1.5\times10^4$\,K for $n_e=10^4$\,cm$^{-3}$; the omission matters in hot, dense
  ionised gas and not in the PDR. The source comment
  attributes the formula to \citet{DrakeUlrich1980}; that paper takes its
  collisional ionisation rates from \citet{Johnson1972} and tabulates
  $4.88\times10^{-16}$\,cm$^3$\,s$^{-1}$ for the ground level at $10^4$\,K (their
  Table~3), so the formula's origin is not established
  (Section~\ref{sec:limitations});
\item the radiative recombination $X^++e^-\to X+\mathrm{PHOTON}$ takes
  MOCASSIN's total recombination coefficient (radiative, from the fits of
  \citealt{VernerFerland1996} and the others listed at
  \src{update\_mod.f90:\allowbreak{}4101--4115}, plus dielectronic) in the first of its
  temperature ranges, the other ranges set to zero;
\item \Htwo{} photoionisation, $\mathrm{H_2}+h\nu\to\mathrm{H_2^+}+e^-$, is
  added to the slot $\mathrm{H_2}+\mathrm{CRP}\to\mathrm{H_2^+}+e^-$ (296) with
  $\Gamma_{\mathrm{H_2}}=\sum_{h\nu_k\ge15.4\,\mathrm{eV}}\Jste_k\sigma_{\mathrm{H_2}}(\nu_k)/h\nu_k$
  and the fit of \citet{Yan1998}, as revised in their erratum \citep{Yan2001},
  to the cross section (\src{update\_mod.f90:\allowbreak{}2863--2877}; \rtn{yanH2xSec},
  \src{ionization\_mod.f90:\allowbreak{}385--411}).
\end{itemize}
These rates are stored by \rtn{updateCell} at every trial temperature of its
search, so the chemistry reads those of the temperature it settled on
(\src{update\_mod.f90:\allowbreak{}2838--2878}). In particular, \textbf{carbon is
photoionised at MOCASSIN's rate}: the MC field contains dust and C\,\textsc{i}
continuum absorption but no \Htwo{} line absorption, so the C\,\textsc{i} rate is
not reduced by the \Htwo{} lines that overlap the C\,\textsc{i} continuum
(Section~\ref{sec:limitations}). The C\,\textsc{i} treatment of
\citet{KampBertoldi2000} in 3D-PDR (\rtn{CIPDRATE},
\src{pdr/\allowbreak{}shield.F90:\allowbreak{}132--157}),
$k=\alpha\chi\exp[-\gamma\AV-1.1\times10^{-17}N_{\mathrm C}-0.9\,T^{0.27}(N_{\mathrm{H_2}}/1.59\times10^{21})^{0.45}]$,
is computed with $\AV=0$ and $N_{\mathrm C}=0$, i.e.\ with its \Htwo{} term only,
and then discarded. (Equation~8 of \citealt{KampBertoldi2000}, which the code
cites, normalises the \Htwo{} term at $10^{22}$\,cm$^{-2}$, not $1.59\times10^{21}$; the
vendored term is 2.3 times larger at a given column.) The ionisation balance of
the linked elements is moreover imposed from MOCASSIN before and after the
solve (Section~\ref{sec:impose}).

\subsection{Cosmic rays and cosmic-ray-induced photons}
\label{sec:cr}

The cosmic-ray ionisation rate is the input \kw{pdrZeta} (default
$10^{-16}$\,s$^{-1}$, \src{common\_mod.f90:\allowbreak{}902}), stored normalised to the
UMIST reference rate $\zeta_0=1.3\times10^{-17}$\,s$^{-1}$ (\src{pdr\_mod.f90:\allowbreak{}360}),
and is the same in every cell (3D-PDR's attenuation option \kw{CRATTENUATION}
is not compiled). Direct cosmic-ray reactions have $k=\alpha\,\zeta/\zeta_0$, with
$\alpha$ the rate at $\zeta_0$ (\src{pdr/\allowbreak{}calc\_reac\_rates.F90:\allowbreak{}353--370}). Reactions
driven by the ultraviolet photons that the secondary electrons of
cosmic-ray ionisation excite in \Htwo{} \citep{PrasadTarafdar1983} have
\begin{equation}
  k = \alpha\,\frac{\zeta}{\zeta_0}\left(\frac{T}{300\,\mathrm K}\right)^{\beta}\frac{\gamma}{1-\omega}
  \quad[\mathrm{s^{-1}}]
  \label{eq:crphot}
\end{equation}
(\src{pdr/\allowbreak{}calc\_reac\_rates.F90:\allowbreak{}381--401}), the UMIST form of the rate of
\citet[their eq.~7]{Gredel1989}, $\zeta\,p_M/(1-\omega)$ per molecule: $\alpha=\zeta_0$ for
all 15 such reactions of the network, $\gamma=p_M$ is the efficiency of the
photoreaction per ionisation, and $\omega$ the grain albedo, which enters because
the induced photons are removed by grain absorption, $\sigma_g(1-\omega)$, rather
than by extinction. That approximation holds when grains dominate the
absorption of the induced photons, which \citet{Gredel1989} note is the case
for all molecules except CO; the temperature factor is used only by CO
($\beta=1.17$, reaction 313). Here $\omega$ is MOCASSIN's grain albedo in the LW
band (Section~\ref{sec:dust}). The network has 8 \kw{CRP} and 15 \kw{CRPHOT}
reactions. \citet{Heays2017} find newly calculated cosmic-ray-induced rates
that differ from the standard values by up to a factor of five.

\subsection{Charge transfer with hydrogen}
\label{sec:ct}

The charge-transfer reactions of the linked elements with H and \Hp{}
(reactions 156--157, He$^+$\,+\,H; 161--163, O$^+$\,+\,H and \Hp\,+\,O; 167,
\Hp\,+\,Mg; the network has none for carbon) use one coefficient per reaction
in both the chemistry and MOCASSIN's ionisation balance (\rtn{pdrCtCoef},
\src{common\_mod.f90:\allowbreak{}1383--1436}; \src{pdr\_mod.f90:\allowbreak{}3742--3747}; \dlog{D-403}): the
network's $\alpha(T/300)^\beta e^{-\gamma/T}$ in its temperature ranges, replaced by the
fits of \citet{KingdonFerland1996}, $k=a\,t_4^{\,b}(1+c\,e^{d\,t_4})$ with $t_4=T/10^4$\,K,
inside their ranges for He$^+$\,+\,H (6000--$10^5$\,K) and \Hp\,+\,Mg
(5000--30\,000\,K) (\src{pdr\_mod.f90:\allowbreak{}848--855}). For oxygen the network's rates
are fits to \citet{Stancil1999} and are kept. The coefficients are evaluated
at the cell's own $T_e$, not at the clamped temperature of the other rates.

\subsection{\texorpdfstring{\Htwo}{H2} formation on grains}
\label{sec:formation}

The reaction $\mathrm H+\mathrm H\to\mathrm{H_2}$ on grains (reaction 328) takes the
rate of \citet[their eqs.~15--18]{CazauxTielens2002}, whose model
\citet{CazauxTielens2004} develop in full (\rtn{H2\_FORMATION\_RATE},
\src{pdr/\allowbreak{}h2\_form.F90:\allowbreak{}9--134}, flag \kw{H2FORM\_CT02}; called at
\src{pdr/\allowbreak{}calc\_reac\_rates.F90:\allowbreak{}116}):
\begin{gather}
  R_{\mathrm f} = \tfrac12\,v_{\mathrm{th}}\,
     \bigl(\sigma^{\mathrm{sil}}_{\mathrm H}\,\epsilon_{\mathrm{sil}}+\sigma^{\mathrm{gra}}_{\mathrm H}\,\epsilon_{\mathrm{gra}}\bigr)\,
     S_{\mathrm H}(T,T_d)\quad[\mathrm{cm^3\,s^{-1}}],
  \label{eq:rform}\\
  v_{\mathrm{th}} = 1.45\times10^{5}\sqrt{T/100\,\mathrm K}\unit{cm\,s^{-1}},\qquad
  S_{\mathrm H} = \bigl[1+0.04\sqrt{T+T_d}+0.002\,T+8\times10^{-6}\,T^2\bigr]^{-1},\\
  \epsilon = \frac{\xi}{1+F_1+F_2},\quad
  F_1 = \frac{\mu F}{2\nu_{\mathrm{H_2}}e^{-E_{\mathrm{H_2}}/T_d}},\quad
  F_2 = \frac{\bigl(1+\sqrt{(E_{\mathrm{HC}}-E_{\mathrm S})/(E_{\mathrm{HP}}-E_{\mathrm S})}\bigr)^2}{4}\,e^{-E_{\mathrm S}/T_d},\\
  \xi = \Bigl[1+\frac{\nu_{\mathrm{HC}}}{2F}\,e^{-1.5E_{\mathrm{HC}}/T_d}
        \bigl(1+\sqrt{(E_{\mathrm{HC}}-E_{\mathrm S})/(E_{\mathrm{HP}}-E_{\mathrm S})}\bigr)^2\Bigr]^{-1},
\end{gather}
with temperatures in K. $v_{\mathrm{th}}$ is the mean speed of H atoms, the factor
$\frac12$ makes one molecule of two atoms, and $\sigma_{\mathrm H}$ [cm$^2$] is the
geometric grain cross section per H nucleus. $S_{\mathrm H}$ is the sticking
coefficient of \citet[eq.~3.7]{HollenbachMcKee1979}, which they take from
calculations by Burke and Hollenbach; it is 0.58, 0.33 and 0.08 at $T=100$,
300 and 1000\,K (for $T_d=20$\,K). The recombination efficiency $\epsilon$ is
limited by three processes \citep{CazauxTielens2002}: $F_1$, newly formed
molecules that stay on a very cold surface; $F_2$, the evaporation of
physisorbed atoms before they meet a chemisorbed one, which controls the
decline above a few tens of K; and $\xi$, the evaporation of chemisorbed atoms at
several hundred K. The constants are $F=10^{-10}$ monolayers\,s$^{-1}$ (the
accretion rate), $\mu=0.005$, $\nu_{\mathrm{H_2}}=3\times10^{12}$\,s$^{-1}$,
$\nu_{\mathrm{HC}}=1.3\times10^{13}$\,s$^{-1}$, $E_{\mathrm{HC}}=3\times10^4$\,K, and
$(E_{\mathrm S},E_{\mathrm{H_2}},E_{\mathrm{HP}})=(110,320,450)$\,K for silicate and
$(260,520,800)$\,K for graphite. The functional form and $\mu$ are those of
\citet{CazauxTielens2002}; the per-material energies are those tabulated by
\citet[Table~C.1]{Rollig2013} from the work of Cazaux and Tielens, whereas
\citet[Table~1]{CazauxTielens2002} list $(E_{\mathrm S},E_{\mathrm{HP}},E_{\mathrm{HC}})=(200,600,10^4)$\,K
for their silicate surface, the values Cloudy uses; \citet{Rollig2013} also
set the $F_1$ term to zero. With the vendored parameters
(Equation~\ref{eq:rform} evaluated for this document) $\epsilon_{\mathrm{sil}}$ is
close to 1 between 6 and 15\,K, 0.90 at 20\,K, 0.25 at 50\,K and 0.10 at 100\,K, and
$\epsilon_{\mathrm{gra}}$ close to 1 between 10 and 40\,K, 0.43 at 100\,K and 0.12 at 300\,K;
both fall to zero below about 6 and 9\,K ($F_1$) and above about 800\,K ($\xi$).
The silicate energies are the main reason why the silicate coefficient is
0.50--0.91 of Cloudy's on \texttt{hiipdr} (\dlog{L-33}, \dlog{H-022}).

The grain areas per H $\sigma^{\mathrm{sil}}_{\mathrm H}$, $\sigma^{\mathrm{gra}}_{\mathrm H}$ are
MOCASSIN's (Section~\ref{sec:grainarea}), handed over per cell at
\src{pdr\_mod.f90:\allowbreak{}3295--3307}; $T_d$ is MOCASSIN's mean grain temperature
(Section~\ref{sec:Td}), floored at $T_{\mathrm{CMB}}$; the metallicity factor is 1
(\src{pdr\_mod.f90:\allowbreak{}347}). The formation rate per unit volume is
$R_{\mathrm f}\,n_{\mathrm H}\,n(\mathrm H)$ [cm$^{-3}$\,s$^{-1}$]. The same function, evaluated
with one material's area at a time, weights the two formation distributions of
the level model (Section~\ref{sec:formdist}).

% ======================================================================
\section{The chemical solve}
\label{sec:chemistry}
% ======================================================================

\subsection{Network and rate law}
\label{sec:network}

The network is 3D-PDR's REDUCED network (files
\src{dataPDR/\allowbreak{}species\_reduced.d} and \src{dataPDR/\allowbreak{}rates\_reduced.d}), in the
format of the UMIST Database for Astrochemistry. \citet{Bisbas2012}
introduced its 33 species with 320 reactions drawn from UMIST 2006
\citep{Woodall2007}; \citet{Bisbas2015} describe the 33-species network then in
use as a 330-reaction subset of UMIST 2012 \citep{McElroy2013}. The vendored
file, from upstream commit \texttt{fabcdc6}, has 331 reactions and records
neither its UMIST release nor the references of its coefficients, so the
release cannot be determined from the files. It has 33 species,
\begin{center}
\small H, H$_2$, H$^+$, H$_2^+$, H$_3^+$, He, He$^+$, C, C$^+$, CH, CH$^+$,
CH$_2$, CH$_2^+$, CH$_3$, CH$_3^+$, CH$_4$, CH$_4^+$, CH$_5^+$, O, O$^+$, OH,
OH$^+$, H$_2$O, H$_2$O$^+$, H$_3$O$^+$, O$_2$, O$_2^+$, CO, CO$^+$, HCO$^+$, Mg,
Mg$^+$, e$^-$,
\end{center}
and 332 reactions: 3D-PDR's 331 plus reaction 332,
$e^-+\mathrm{H_2^+}\to e^-+\mathrm H+\mathrm H^+$ ($2\times10^{-7}e^{-30720/T}$\,cm$^3$\,s$^{-1}$),
appended in this project because the network had no
electron-impact dissociation of \Htwo$^+$, Cloudy's main \Htwo$^+$ sink in ionised
gas (\src{pdr/\allowbreak{}UPSTREAM.md}, \dlog{D-395}); the Arrhenius form is the one of
Cloudy's \src{mole\_co\_umisthack.dat}, which attributes it to
\citet{Janev1987}. Of the 332, 28 are photoreactions, 8 cosmic-ray and 15
cosmic-ray-induced photoreactions (Section~\ref{sec:photorates}), and three
carry the grain marker \texttt{\#}: grain formation of \Htwo{} (328) and the
grain-assisted recombinations of He$^+$ and C$^+$ (330, 331). These two have no
term in the generated right-hand side, so their coefficients are computed but
unused and 330 reactions act in the solve. The other 278 are two-body reactions:
ion--molecule, neutral--neutral and charge-transfer reactions, dissociative and
radiative recombination, radiative association and collisional dissociation.
The ODE right-hand side is generated C code selected at compile time
(\src{pdr/\allowbreak{}odes\_reduced.c}, flag \kw{-DREDUCED}); the keyword \kw{pdrNetwork}
only validates that the requested network is the compiled one.

Two-body coefficients follow the modified Arrhenius form
\begin{equation}
  k(T) = \alpha\left(\frac{T}{300\,\mathrm K}\right)^{\beta} e^{-\gamma/T}\quad[\mathrm{cm^3\,s^{-1}}],
  \label{eq:arrhenius}
\end{equation}
with $\gamma$ in K (\src{pdr/\allowbreak{}calc\_reac\_rates.F90:\allowbreak{}172--252}). Where a reaction
has several entries for different temperature ranges, the first whose upper
limit is above $T$ (else the last) is used and the others are set to zero; a
fit with $\gamma<0$ is evaluated at the lower limit of its range below it, and one
with $\gamma<-200$\,K is set to zero there. The chemistry is evaluated at
MOCASSIN's $T_e$ clamped to [\kw{pdrTmin}, \kw{pdrTmax}] $=[1,\,2\times10^4]$\,K by
default (\src{pdr\_mod.f90:\allowbreak{}3346}), the range over which the fits are meant to be
used; MOCASSIN's own floor \kw{teFloor}, 2.725\,K by default, keeps $T_e$ above
$T_{\mathrm{CMB}}$. The vendored bound $\max(\kw{pdrTmin},T_{\mathrm{CMB}})$
(\src{:\allowbreak{}364}) is read only by 3D-PDR's own temperature iteration
(\src{pdr/\allowbreak{}changetemperature.F90}), which the bridge does not call. In terms of fractional abundances $x_i=n_i/n_{\mathrm H}$ the system is
\begin{equation}
  \frac{\dd x_i}{\dd t} = \sum_{j,k} k_{jk}\,x_jx_k\,n_{\mathrm H}
      + \sum_j k_j\,x_j - x_i\Bigl(\sum_j k_{ij}\,x_j\,n_{\mathrm H} + \sum k_i\Bigr),
  \label{eq:odes}
\end{equation}
with two-body coefficients $k_{jk}$, $k_{ij}$ [cm$^3$\,s$^{-1}$] and first-order rates
$k_j$, $k_i$ [s$^{-1}$] (photoreactions and cosmic rays, Section~\ref{sec:photorates}).
Grain formation enters as $R_{\mathrm f}\,x(\mathrm H)\,n_{\mathrm H}$ in $\dd x(\mathrm{H_2})/\dd t$
and twice that as a loss of H. Inside the right-hand side the electron
abundance is not a free variable: it is the sum of the network's cations plus
\rtn{xe\_extra}, the electrons of the ions the network has no species for
(higher stages, other elements), supplied by the bridge (\dlog{D-402};
\src{pdr/\allowbreak{}odes\_reduced.c:\allowbreak{}58}).

\subsection{Initial conditions and seeding}
\label{sec:seed}

\paragraph{First visit of a cell.} \rtn{seedAbunFromMocassin}
(\src{pdr\_mod.f90:\allowbreak{}3123--3157}, \dlog{D-208}): every molecule starts at zero;
each linked element's neutral and first ion take MOCASSIN's stage-1 and stage-2
fractions times the element abundance; $x_e=n_e/n_{\mathrm H}$. Nothing is read
from the network file's dark-cloud defaults.

\paragraph{First PDR iteration.} \rtn{seedMolecules}
(\src{pdr\_mod.f90:\allowbreak{}2855--3023}, \dlog{D-209}) places a first \Htwo{} and CO
content so that the first shielding columns are not zero. \Htwo{} comes from the
balance of grain formation and unshielded photodissociation, with the neutral
hydrogen of MOCASSIN, $x_{\mathrm n}n_{\mathrm H}$, shared between H and \Htwo:
\begin{gather}
  R_{\mathrm f}\,n_{\mathrm H}\,n(\mathrm H) = D\,n(\mathrm{H_2}),\qquad
  n(\mathrm H)+2n(\mathrm{H_2}) = x_{\mathrm n}n_{\mathrm H}
  \quad\Longrightarrow\quad
  x(\mathrm{H_2}) = \frac{R_{\mathrm f}\,x_{\mathrm n}\,n_{\mathrm H}}{D+2R_{\mathrm f}n_{\mathrm H}},\\
  D = \alpha_{\mathrm{H_2}}\,\chi\;\mathrm{SCATTER}(\AV,1000\,\text{\AA}),
\end{gather}
with $\alpha_{\mathrm{H_2}}$ the network's coefficient ($5.18\times10^{-11}$\,s$^{-1}$, read
from the rate file) and $\chi$ the local MC field of Equation~(\ref{eq:chi});
$x(\mathrm{H_2})$ is capped at $x_{\mathrm n}/2$. CO follows the ansatz
$x(\mathrm{CO}) = x(\mathrm C)\,[2x(\mathrm{H_2})/x_{\mathrm n}]$, capped by the oxygen,
with carbon and oxygen conserved. Steps 4--5 of Table~\ref{tab:steps} are then
repeated on the seeded abundances. The seed attenuates the already attenuated
local field $\chi$ once more with \rtn{SCATTER} and ignores self-shielding; it
only provides the starting point, which the iteration then forgets.

\subsection{Imposition of MOCASSIN's ionisation before the solve}
\label{sec:impose}

MOCASSIN owns the ionisation state; the chemistry owns what is locked in
molecules (\dlog{D-153}, \dlog{D-193}, \dlog{D-402};
\src{pdr\_mod.f90:\allowbreak{}3503--3584}). For each linked element $X$ with abundance
$A_X$ and molecular content $M_X=\sum_i \nu_{X,i}x_i$ (atoms of $X$ in each
species, from the species' formulae, \rtn{speciesAtoms}):
\begin{equation}
  x(X) = (A_X-M_X)\,\frac{f_1}{\sum_s f_s},\qquad
  x(X^+) = (A_X-M_X)\,\frac{f_2}{\sum_s f_s},
\end{equation}
where $f_s$ are MOCASSIN's stage fractions; what MOCASSIN holds in stages
above the second stays MOCASSIN's. The electron abundance is set to
$n_e/n_{\mathrm H}$ from MOCASSIN, and \rtn{xe\_extra} to $x_e$ minus the
linked first ions (\src{:\allowbreak{}3565--3581}). The imposition fixes the starting
point of the integration; inside it, the linked atoms are ionised and
recombined at MOCASSIN's own rates (Section~\ref{sec:linked}), so the solver's
neutral/ion ratio should stay close to MOCASSIN's, and the remaining
difference is carried by the network's other channels (charge transfer,
ion--molecule reactions).

\subsection{Integration to steady state}

\rtn{calculate\_abundances} (\src{pdr/\allowbreak{}calculate\_abundances.c}), called at
\src{pdr\_mod.f90:\allowbreak{}3749}, integrates Equation~(\ref{eq:odes}) with CVODE from
SUNDIALS 7.8.0 (\dlog{D-015}; \citealt{Hindmarsh2005,Gardner2022}). The line numbers
in this list are those of \src{pdr/\allowbreak{}calculate\_abundances.c} unless a file is
named:
\begin{itemize}
\item variable-order backward differentiation formulae (\verb|CV_BDF|, :196),
  the stiff integrator, with Newton iteration on a dense direct linear solver
  (\verb|SUNDenseMatrix|, \verb|SUNLinSol_Dense|, :229--238) and CVODE's
  difference-quotient Jacobian (the analytic Jacobian is not used);
\item relative and absolute tolerances $10^{-8}$ and $10^{-30}$ on the fractional
  abundances (set at \src{pdr\_mod.f90:\allowbreak{}369--370}, applied at :212); at most
  $10^7$ internal steps per output interval (:115);
\item output times from $10^{-4}$\,yr, multiplied by 10 until
  \kw{pdrEndTime} (default $10^7$\,yr) (:246, :364);
\item \kw{CHEMRETRY} (:248--384): a failed integration is restarted from the
  initial state up to four more times with a forced first step of
  $10^{-2}$\,s, divided by 100 at each retry;
\item \kw{CHEMSTEADY} (:387--437): the integration then continues in blocks that
  each double the elapsed time, at most seven ($128\times$\kw{pdrEndTime},
  $1.3\times10^9$\,yr by default), until no species above the absolute tolerance
  changes by more than 1\% over a block;
\item abundances below the absolute tolerance are set to zero and the
  electron entry is restored (:440--452).
\end{itemize}
The result is meant to be the steady state of the network at fixed
temperature, field, columns and imposed ionisation; there is no time
dependence carried between MC iterations. A species that relaxes
exponentially on a time scale $\tau$ from an offset of order unity passes the 1\%
test over a doubling once the elapsed time exceeds about $4.6\,\tau$, so species
that relax more slowly than about $10^8$\,yr cannot pass it by the last block;
if the seventh block does not pass, the abundances of that time are returned
without a flag. An
unchanged abundance vector is counted (\rtn{pdrNstuck},
\src{pdr\_mod.f90:\allowbreak{}3860}) but not treated as a failure, since it is also the
signature of a converged cell. The slowest chemical process in molecular gas
is usually \Htwo{} formation, whose time scale $1/(R_{\mathrm f}n_{\mathrm H})$ is
$\sim10^7$\,yr at $n_{\mathrm H}=100$\,cm$^{-3}$; where the age of the object is shorter
the steady state overstates the molecular content
(Section~\ref{sec:assumptions}).

\subsection{After the solve}

\begin{enumerate}
\item The solver's own neutral/ion split of each linked element is
  overwritten with MOCASSIN's stage-1:stage-2 ratio, keeping the solver's
  total $x(X)+x(X^+)$ (\src{pdr\_mod.f90:\allowbreak{}3770--3784}, \dlog{D-162}); the
  electron abundance is reset to MOCASSIN's (\src{:\allowbreak{}3801--3804}, \dlog{D-231}).
  The solver's own ratio is kept for diagnostics (\dlog{D-402}, \dlog{D-404}).
\item The abundances are stored (\src{:\allowbreak{}3922}) and $n(\mathrm{H_2})=x(\mathrm{H_2})n_{\mathrm H}$
  is written for MOCASSIN's opacity (\src{:\allowbreak{}4573}).
\item What the network does to each linked element's $X$ and $X^+$ beyond the
  channels \rtn{ionBalance} computes itself (photoionisation, cosmic rays,
  radiative recombination, charge transfer with H) is classified reaction by
  reaction at the solver's state, as net conversions $X\leftrightarrow X^+$ and
  as sources and sinks to and from molecules, and stored for MOCASSIN's next
  ionisation solve (\rtn{pdrIonRate}, \src{:\allowbreak{}3924--4054}, \dlog{D-402}). The
  classified totals reproduce the right-hand side's own formation and
  destruction totals to $4\times10^{-16}$ (\dlog{D-402}).
\end{enumerate}

% ======================================================================
\section{Feedback from the molecules on MOCASSIN}
\label{sec:feedback}
% ======================================================================

The molecular content computed in Section~\ref{sec:chemistry} acts back on
the radiation field, the ionisation and the temperature, which closes the
loop of Figure~\ref{fig:chain}. It does so through three channels only: the
fraction of each element held in molecules, the \Htwo{} photoionisation
continuum, and the heating and cooling terms. No molecular line enters the MC
transfer.

\subsection{Elements locked in molecules}

The whole handback for an element is the fraction of it \emph{not} in
molecules (\rtn{setElemAvail}, \src{pdr\_mod.f90:\allowbreak{}5120--5144}, called by
\rtn{pdrToMocassin}, \src{:\allowbreak{}5038--5104}):
\begin{equation}
  \rtn{elemAvail}_X = \min\!\Bigl[\max\Bigl(1-\frac{M_X}{A_X},\,0\Bigr),\,1\Bigr],
  \qquad M_X = \sum_i \nu_{X,i}\,x_i ,
\end{equation}
for the five linked elements H, He, C, O and Mg, with $A_X$ MOCASSIN's own
abundance ($A_{\mathrm H}=1$) and $\nu_{X,i}$ the number of atoms of $X$ in
species $i$, read from the species' formula for every species other than $X$
and $X^+$ themselves (\src{:\allowbreak{}585--603}, \dlog{D-402}). MOCASSIN's
\rtn{ionBalance} constrains the sum of the element's ion stages to it
(\src{update\_mod.f90:\allowbreak{}3130--3134}): carbon in CO therefore appears in no
atomic stage, emits no atomic line and gives no free electron. The chemistry
writes no ionisation state and no temperature (\dlog{D-193}). After each
iteration \rtn{rescaleIonToAvail} (\src{pdr\_mod.f90:\allowbreak{}2764--2803}) rescales the
stages of every element to \rtn{elemAvail}, keeping MOCASSIN's ratios between
them, so that the hydrogen budget
$x(\mathrm H^0)+x(\mathrm H^+)+2x(\mathrm{H_2})=1$ holds in every cell (up to
the hydrogen in other molecules) before the next transfer builds its opacity
(\dlog{D-220}). It prints the size of the pre-rescale defect (\texttt{[D220]}),
which is the informative test of the budget: since the kept second solve
(below) no cell of either fixture is off by more than 1\% in the written model
(\dlog{L-29}).

The network channels classified in Section~\ref{sec:chemistry} (sources and
sinks of $X$ and $X^+$ to and from molecules, and net conversions between
them) enter \rtn{ionBalance}'s rate matrix as sinks on the diagonal and
sources on the right-hand side, the sum still constrained to \rtn{elemAvail}
(\src{update\_mod.f90:\allowbreak{}3255--3300}, \dlog{D-402}). They are those of the
previous chemistry, i.e.\ one iteration behind (\dlog{L-47}).
Cosmic-ray ionisation, which MOCASSIN otherwise lacks, is added as MOCASSIN's
own physics with the network's coefficients: $\alpha\,\zeta/\zeta_0$ for direct
ionisation and
$\alpha\,(\zeta/\zeta_0)(T/300\,\mathrm K)^\beta\gamma/(1-\omega)$ for the
cosmic-ray-induced photons, $\zeta_0=1.3\times10^{-17}$\,s$^{-1}$, for every
linked neutral (\src{update\_mod.f90:\allowbreak{}2916--2930}; \dlog{D-193}, \dlog{D-280}).

\subsection{\texorpdfstring{\Htwo}{H2} in the gas opacity}

$n(\mathrm{H_2})$ (\rtn{H2den}) adds the \Htwo{} photoionisation continuum
above 15.4\,eV to the gas opacity of the next transfer
(\rtn{addOpacity}, \src{ionization\_mod.f90:\allowbreak{}435--472}, \dlog{D-107}),
\begin{equation}
  \kappa_\nu \mathrel{+}= n(\mathrm{H_2})\,\sigma_{\mathrm{H_2}}(\nu),\qquad h\nu\ge15.4\unit{eV},
\end{equation}
with the cross section fitted by \citet{Yan1998}, as revised in their erratum
\citep{Yan2001} (\rtn{yanH2xSec},
\src{ionization\_mod.f90:\allowbreak{}385--411}; zero below 15.4\,eV, as in Cloudy, and
clipped at zero just above threshold where the fit goes negative). The
comment above the opacity loop (\src{:\allowbreak{}444--453}) still describes ``a
single power law anchored at threshold''; the function actually used is the
piecewise fit of \citet{Yan1998}. This continuum matters only for
hydrogen-ionising photons that reach molecular gas. No \Htwo{}, CO or
H\,\textsc{i} Lyman-series \emph{line} opacity is added to the MC transfer:
line absorption is represented only through the shielding columns.

\subsection{Heating and cooling}
\label{sec:thermal}

MOCASSIN solves the thermal balance of every cell, heating equal to cooling
per H nucleus, with \rtn{thermBalance} (\src{update\_mod.f90:\allowbreak{}1734--2456})
inside its temperature search \rtn{iterateT} (\src{:\allowbreak{}609--1191}). Its own
terms are the photoionisation heating by every ion on the MC field; the grain
photoelectric heating (\rtn{setPhotoelHeatCool},
\src{update\_mod.f90:\allowbreak{}1460--1556}), summed over every grain species and
size from the MC field above the photoelectric threshold (work function plus
grain potential), with the yield of \citet{WeingartnerDraine2001}
(\dlog{D-289}) and the grain charge of MOCASSIN's own charge balance, whose
electron sticking was set to the same paper's value (\dlog{D-326}); the
recombination, free--free and collisionally excited line cooling of every ion;
and the gas--grain collisional exchange (\rtn{setDustGasCollHeatCool},
\src{:\allowbreak{}1562--1731}), which omits collisions with \Htwo{} (\dlog{L-9}).
The molecules enter it in three ways.

\begin{enumerate}
\item \textbf{Coolants.} The coolants named by \kw{pdrCoolants} (in the
  fixtures \Cp, C, O and $^{12}$CO, LAMDA-format files, \citealt{Schoier2005})
  are treated in two ways.
  \begin{itemize}
  \item \Cp, C and O are cooled by MOCASSIN's own model atoms (its electron
    collision data and all its levels), to which the coolant files' rates for
    collisions with H, \Hp, He and \Htwo{} and the chemistry's escape
    probability of each transition are added (\src{emission\_mod.f90:\allowbreak{}1861--1887};
    \src{common\_mod.f90:\allowbreak{}1489--1635}; \dlog{D-363}). They are solved at every
    trial temperature, and the same atoms give the reported line fluxes, so
    the thermal balance and the lines share one atom per ion.
  \item CO, which MOCASSIN has no atom for, is solved by 3D-PDR's
    escape-probability level-population solver in \rtn{pdrStep} at 16
    temperatures log-spaced over [\kw{pdrTmin}, \kw{pdrTmax}], a factor 1.94
    apart for the defaults (\rtn{initPDRcoolTab},
    \src{pdr\_mod.f90:\allowbreak{}2713--2732}, \dlog{D-215}). The temperature search reads
    this table by log--log interpolation, clamped at its ends (\rtn{pdrCoolAt},
    \src{common\_mod.f90:\allowbreak{}1445--1486}), and adds it to MOCASSIN's cooling
    (\src{update\_mod.f90:\allowbreak{}2232--2249}). For CO the table misses an exact solve
    at the cell's temperature by a mean $|\Delta\ln\Lambda|$ of 0.02
    (\texttt{hiipdr}) and 0.004 (\texttt{pn150k}) per cell (\dlog{L-36}).
  \end{itemize}
  In both the level populations are in statistical equilibrium with the
  colliders H, \Hp, $e^-$, He and \Htwo{} at the chemistry's abundances (the
  linked ions and $n_e$ being MOCASSIN's). The escape probability of a
  transition $u\to l$ is the mean over the twelve line rays of
  Section~\ref{sec:lineRays} of $\beta=(1-e^{-\tau})/\tau$, with
  \begin{equation}
    \tau = \frac{A_{ul}\,c^3}{8\pi\nu_{ul}^3\,b}\sum_{\mathrm{cells}}
           \Bigl(n_l\frac{g_u}{g_l}-n_u\Bigr)\ell ,
  \end{equation}
  each crossed cell contributing its own level populations as last stored
  (\src{pdr/\allowbreak{}escape\_probability.F90:\allowbreak{}121--137, 196--202}; \dlog{D-377}). With
  \kw{-DSOBOLEV} the rate equations use the equivalent net-rate form, the
  spontaneous rate $A_{ul}\beta$ and the radiative rates driven by
  $\beta J_{\mathrm{bg}}$ (\src{:\allowbreak{}219--233}), where the background
  $J_{\mathrm{bg}}$ is MOCASSIN's own $J_\nu$ at the line plus the CMB
  (\dlog{D-377}). The cooling is the net escaping power,
  $\Lambda_{ul}=A_{ul}h\nu_{ul}n_u\beta\,(S_{ul}-J_{\mathrm{bg}})/S_{ul}$ with
  $S_{ul}$ the line source function (\src{:\allowbreak{}209--211}).
\item \textbf{Extra heating.} 3D-PDR's heating routine \rtn{CALC\_HEATING}
  (\src{pdr/\allowbreak{}heatingfunctions.F90}) is called at \src{pdr\_mod.f90:\allowbreak{}4499}
  with the local field $\chi$ at MOCASSIN's $T_e$. The terms MOCASSIN computes
  itself are excluded: grain and PAH photoelectric heating (3D-PDR's
  \citealt{BakesTielens1994} expression with the modifications of
  \citealt{Wolfire2003}, and its alternative \citealt{WeingartnerDraine2001}
  fit), carbon photoionisation heating and gas--grain collisions; these are
  the only terms that read $\chi$. The rest, \rtn{pdrHeatXtra}
  (\src{pdr\_mod.f90:\allowbreak{}4545--4546}), is added to MOCASSIN's heating
  (\src{update\_mod.f90:\allowbreak{}2391}) and held fixed during the temperature search:
  \begin{itemize}
  \item \emph{\Htwo{} terms from the level model} (\src{pdr\_mod.f90:\allowbreak{}4504--4528},
    \dlog{H-013}): formation, $R_{\mathrm f}\,n_{\mathrm H}\,n(\mathrm H)\,E_{\mathrm{kin}}$,
    with $E_{\mathrm{kin}}$ the kinetic energy of a newly formed molecule from the
    formation distribution; dissociation, $n(\mathrm{H_2})\sum_l n_l H^{\mathrm{dis}}_l$,
    the kinetic energy of the atoms; and the net collisional term
    $n(\mathrm{H_2})\,\mathcal{H}_{\mathrm{coll}}$, which is the heating by the
    de-excitation of pumped and newly formed molecules and, where the gas is
    warmer than the level populations, \Htwo's line cooling
    (Section~\ref{sec:solvecell}). They replace 3D-PDR's fixed energies
    (1.5\,eV per formation, which the vendored code attributes to
    \citealt{HollenbachTielens1999}; 0.4\,eV per dissociation;
    $9\times2.2$\,eV per dissociation for pumping, with the critical-density
    factor it attributes to \citealt{HollenbachMcKee1979}), which are kept
    beside them for comparison (\src{h2Cells.out});
  \item cosmic-ray heating, $\Gamma_{\mathrm{CR}}=9.4\,\mathrm{eV}\times\zeta\,n(\mathrm{H_2})$,
    i.e.\ $1.5\times10^{-11}\zeta\,n(\mathrm{H_2})$\,erg\,cm$^{-3}$\,s$^{-1}$, the
    expression of \citet[their eq.~A23]{TielensHollenbach1985}: about 8\,eV of
    heat per primary ionisation in molecular gas, helium included, with $\zeta$
    the rate per \Htwo{}. Ionisations in atomic gas deposit no heat in this
    term, which \citet{TielensHollenbach1985} justify by cosmic-ray heating
    being unimportant there;
  \item turbulent dissipation, $\Gamma_{\mathrm{turb}}=3.5\times10^{-28}\,(v_{\mathrm{turb}}/\mathrm{km\,s^{-1}})^3\,
    (L_{\mathrm{turb}}/\mathrm{pc})^{-1}\,n_{\mathrm H}$\,erg\,cm$^{-3}$\,s$^{-1}$ with
    $L_{\mathrm{turb}}=5$\,pc fixed and $v_{\mathrm{turb}}$ = \kw{pdrVturb}, which the
    vendored code attributes to \citet{Black1987}: $7\times10^{-29}$\,erg\,s$^{-1}$
    per H nucleus at the default 1\,km\,s$^{-1}$;
  \item the energy released by exothermic reactions (for example the
    dissociative recombination of HCO$^+$, H$_3^+$ and H$_3$O$^+$, and He$^+$ with
    \Htwo{} and CO),
    the reactions being found automatically in the network, after
    \citet{Clavel1978}, who found that such reactions could contribute
    significantly to the heating if their energy is released entirely as
    kinetic energy, as is assumed here.
  \end{itemize}
\item \textbf{Photoelectric heating} is MOCASSIN's own, from its grains, as
  described above (\dlog{D-289}, \dlog{D-326}); 3D-PDR's PAH term is not used
  (\dlog{D-134}, \dlog{D-352}). The model grains include no PAHs (\dlog{L-7}),
  so the heating lacks the contribution of PAHs and very small grains that
  dominates in the \citet{BakesTielens1994} prescription of most PDR codes
  and of the code comparison \citep{Rollig2007}; the Cloudy reference models
  omit them too. Against Cloudy the photoelectric efficiency per H per unit
  field is 0.66--1.10 with no trend with depth on \texttt{pn150k}
  (\dlog{D-304}); on \texttt{hiipdr} the gas temperature at $\AV=0.73$ and 0.97
  is 1.07 and 1.01 times Cloudy's since the sticking fix, the deeper gas
  remaining colder (\dlog{L-22}).
\end{enumerate}

The gas temperature found this way is the one the next chemistry runs at,
and, one iteration later, the level model. It acts on the molecules through
the rate coefficients, through the \Htwo{} formation coefficient (the thermal
speed of the H atoms and their sticking probability,
Section~\ref{sec:formation}), through the collisional excitation of \Htwo{}
and of the coolants, and through the CO cooling it controls. The dust
temperature acts through the formation efficiency (Section~\ref{sec:Td}).
These are channels by which the field sets the molecular content even where
no FUV photon reaches the molecules.

\subsection{The second ionisation solve}

After the chemistry, \rtn{ionBalance} is solved a second time in every cell,
at the temperature of the first solve, on the channels and the molecular
content this iteration's chemistry has just left; it is counted against the
first and the model is put back (\rtn{ionSecondSolve},
\src{iteration\_mod.f90:\allowbreak{}3361--3434}, called at \src{:\allowbreak{}2168}; \dlog{D-404}). A
control pass on the old terms, before the chemistry (\src{:\allowbreak{}1274}), moves
nothing. On the finished model the second solve's stages replace the first's
(\rtn{ionKeepSecond}, \src{:\allowbreak{}3455--3521}), so the written ions are consistent
with the last chemistry, and the level model's final call runs on them
(\dlog{H-029}). The second solve removes the one-iteration delay of
\dlog{L-47}: the neutral cells whose solver $X^+/X$ is more than 10\% from
MOCASSIN's fall from 28.4\% to 0.0\% for hydrogen on \texttt{hiipdr}, and to
0.0\% for every linked element except oxygen there (9.7\%, \dlog{L-50}). Neither
the temperature nor the chemistry is solved again on the second solve's ions,
and in the iterations before the last the transfer, the thermal balance and
the chemistry still run on the first solve's.

% ======================================================================
\section{Iteration control and convergence}
\label{sec:convergence}
% ======================================================================

\subsection{The fixed point}

The model is converged when the field, the ionisation, the temperatures, the
molecular abundances and the \Htwo{} level populations reproduce one another
from one MC iteration to the next. Because the columns and the level
populations lag by one iteration (Section~\ref{sec:lags}) and the level model
makes one sweep per iteration, the shielding and the populations converge
together with the rest of the model rather than inside each iteration. This
is a Picard (successive-substitution) iteration on the coupled system; there
is no acceleration of the outer iteration. Damping across iterations
(\kw{iterDamp}, \src{iteration\_mod.f90:\allowbreak{}2424--2467}) is off by default
(\dlog{D-406}). Because the MC field is redrawn at every iteration, the fixed
point is reached only to within the MC noise of
Section~\ref{sec:mcnoise}.

\subsection{Tests}

\begin{itemize}
\item \textbf{Integrated quantities} (\src{common\_mod.f90:\allowbreak{}798--818},
  test at \src{iteration\_mod.f90:\allowbreak{}3243}; \dlog{D-391}, \dlog{D-392}): the
  emission measure $\int n_e n_i\,\dd V$ of every ion stage holding at least
  1\% of its element, the H$^0$, \Htwo{} and CO contents $\int n\,\dd V$, and the
  gas temperature weighted by $n_e\,n(\mathrm H^+)$, by $n(\mathrm H^0)$ and by
  $n(\mathrm{H_2})$. The run stops when the largest
  $|q_n/q_{n-3}-1|\le$ \kw{convTotals} (default 0.02), and not before the first
  iteration after \kw{pdrStartIter}. The test does not see the packet count:
  on \texttt{hiipdr} it stops at iteration 14 or 15 from $6\times10^6$ to
  $3.84\times10^8$ packets while the carbon turn moves by 0.5\,mag
  (\dlog{L-52}, \dlog{D-408}); and the temperature weighted by $n(\mathrm H^0)$ is
  carried by the few hundred hot cells of the ionisation front, whose
  iteration-to-iteration scatter alone puts the 2\% limit at only 1.4--1.5
  standard deviations (\dlog{L-52}).
\item \textbf{Per-cell tests} (MOCASSIN's own, on $x(\mathrm H^0)$ and $T_e$),
  extended to $x(\mathrm{H_2})$: the abundance after the chemistry against the
  copy taken at the start of the iteration (\src{iteration\_mod.f90:\allowbreak{}1173--1188,
  2484--2512}, \dlog{D-179}, \dlog{D-391}). They are reported; the stop is the
  integrated test, since the converged-cell fraction plateaus at about 72\% in
  the PDR fixtures from the per-cell MC noise (\src{iteration\_mod.f90:\allowbreak{}3236--3241}).
\item \textbf{Level model}: inside its final call, the luminosities of the
  four \Htwo{} lines 1--0\,S(1), 2--1\,S(1), 0--0\,S(0) and 0--0\,S(1) to 1\% between
  sweeps (Section~\ref{sec:sweeps}). Each call also reports the median and
  90th percentile of $|k_{\mathrm{Sol}}/k_{\mathrm{Sol}}^{\mathrm{last\ call}}-1|$ over cells with
  $x(\mathrm{H_2})\ge10^{-6}$ (\src{pdr\_mod.f90:\allowbreak{}5499--5520}).
\end{itemize}

\subsection{Monte Carlo noise}
\label{sec:mcnoise}

Since every photorate reads its own cell's MC field (\dlog{D-350}), the
chemistry of a cell inherits the cell's sampling noise, redrawn at every
iteration because $J$ is rebuilt from zero. On \texttt{hiipdr} at
$6\times10^6$ packets the seed-to-seed scatter of the FUV field per cell is
2--4\% at $\AV\sim1$, 7\% at 1.6--2.1, 14\% at 2.1--2.6 and 22\% (90th percentile
38\%) at 2.6--3.05 (\dlog{L-26}). It follows $N^{-1/2}$ in the number of packet
tracks per cell: on a $61^3$ grid with $5\times10^8$ packets, 21 times the tracks
per cell, the iteration-to-iteration noise is 0.8, 1.5, 2.8 and 4.8\% at
$\AV=1$--1.5, 1.5--2, 2--2.5 and 2.5--3.5 (\dlog{D-365}). For the level model the
noise is that of a single 2\%-wide frequency bin per line rather than of the
band: its per-cell Solomon rate scatters about 1.5 times as much as the
chemistry's band-integrated rate on the same field, and moved by 50--60\%
(median) between iterations against 3.4\% on a fixed field (\dlog{H-012}).

The noise has three consequences beyond scatter. Through the non-linear
dependence of the photorates on the field it biases cell means, by an
estimated $+5\%$ at 22\% scatter (\dlog{L-26}), and the bias of the means is one
way: fewer packets give a shallower carbon turn and colder deep gas on
\texttt{hiipdr} (\dlog{D-408}). At the carbon turn the noise is not scatter about
a mean but two states: 85--97\% of the cells in each 0.1\,mag bin between
$\AV=1.9$ and 2.6 are either ionised in carbon or in the state maintained by
cosmic rays alone, depending on whether a packet of the 11.26--13.6\,eV band
crossed them in that iteration, so ``C$^+=$C'' marks where half the cells are
missed (\dlog{D-405}). And the noise is redrawn in full at every iteration: the
spread of a cell's mean over seven to nine iterations among
symmetry-equivalent cells is 0.33--0.41 of one iteration's (\dlog{L-26}). An
estimator accumulated over iterations is registered for future development
(\dlog{D-411}).

What converges at the default $6\times10^6$ packets was measured by
quadrupling the count three times (\dlog{D-408}, \dlog{H-032}). The \Htwo{} lines,
the ionised gas's lines, [O\,\textsc{i}], H$\beta$ and the H/\Htwo{} crossing are
converged at the default ($\AV=1.395$ at $6\times10^6$ and 1.394 at
$3.84\times10^8$ packets on \texttt{hiipdr}; 1--0\,S(1) changes by 0.8\%). The
carbon turn and the [C\,\textsc{i}] lines need 16--64 times the default, and CO
still moves at 64 times (by 5--10\% per factor of four on \texttt{hiipdr}).
Per-cell values of the deep chemistry are not converged at any of these
counts.

\subsection{End of the run}

When the stop test passes, or at the maximum iteration count, the code keeps
the second ionisation solve (\rtn{ionKeepSecond}), calls
\rtn{h2LevelModel(grid,12)} (\src{iteration\_mod.f90:\allowbreak{}3254, 3283}) and writes
the outputs; the chemistry is not solved again on the final level populations
(Section~\ref{sec:limitations}). The level model's outputs are written by
\src{h2Lines\_mod.f90}: \src{output/\allowbreak{}lineFluxH2.out} (every \Htwo{} $X$-state
line above $10^{-6}$ of the strongest, integrated over the grid, optically
thin, no extinction), \src{output/\allowbreak{}h2Cells.out} (per cell: $n(\mathrm{H_2})$,
$T_e$, $k_{\mathrm{Sol}}$ and the chemistry's own rate on the same field, pumping
rate, ortho/para, fraction above 4100\,cm$^{-1}$, the \Htwo{} heating terms beside
3D-PDR's and MOCASSIN's heating and cooling) and \src{output/\allowbreak{}h2Pop.bin} (all
populations, read back by a warm start).

% ======================================================================
\section{Physical assumptions}
\label{sec:assumptions}
% ======================================================================

The assumptions below are made, explicitly or implicitly, on the path from
the field to the molecular content. Each names the section where it enters,
says why it is made and where it fails, and gives the direction of the
resulting bias and its measured size where the project has measured it.
Section~\ref{sec:limitations} lists the defects and open issues.

\subsection{Time dependence and equilibrium}
\begin{enumerate}
\item \textbf{Chemical steady state.} The network is integrated until no
  species changes by more than 1\% over a doubling of the integration time, at
  fixed temperature, field, columns and imposed ionisation
  (Section~\ref{sec:chemistry}). No chemical age, advection or time-dependent
  field is represented. The slowest process is \Htwo{} formation, with a time
  scale $1/(R_{\mathrm f}n_{\mathrm H})\approx10^7\,(n_{\mathrm H}/100\,\mathrm{cm^{-3}})^{-1}$\,yr
  for $R_{\mathrm f}\approx3\times10^{-17}$\,cm$^3$\,s$^{-1}$, close to the model's
  coefficient (\dlog{L-33}): $10^7$\,yr at \texttt{hiipdr}'s density and $10^6$\,yr
  at \texttt{pn150k}'s. This is comparable to or longer than the age of an
  H\,\textsc{ii} region and much longer than that of a planetary nebula, so where molecules
  have to form from atomic gas the steady state overstates the molecular
  content, and it describes neither gas just uncovered by an advancing front
  nor molecules inherited from an earlier phase.
\item \textbf{Statistical equilibrium of the \Htwo{} levels}, with the
  excited electronic states treated as instantaneous (radiative lifetimes
  0.5--9\,ns in the data; collisions compete only above
  $\sim10^{14}$\,cm$^{-3}$) (Sections~\ref{sec:balance}, \ref{sec:pump}). The slowest radiative
  relaxation within $X$, the quadrupole decay of the lowest rotational levels
  ($A=2.9\times10^{-11}$\,s$^{-1}$ for $J=2\to0$, about 1100\,yr, in
  \src{dataPDR/\allowbreak{}h2/\allowbreak{}transprob\_X.dat}), is short compared with the dynamical
  times of the modelled objects, and collisions shorten it further.
\item \textbf{Formation closure in the level model}: every molecule removed by
  the Solomon process is replaced by one formed on a grain, and the level
  populations are normalised to one. The formation flux into the levels
  therefore equals the Solomon flux; destruction by other channels and the
  corresponding formation are not in the level balance
  (Section~\ref{sec:balance}). Where the Solomon rate has been shielded away
  and cosmic rays and ion--molecule reactions dominate the destruction,
  formation pumping, and the internal energy it brings, are underestimated
  (Section~\ref{sec:limitations}). With the network's \Htwo\,+\,CRP channels
  alone the other destruction is $0.99\,\zeta$, $10^{-16}$\,s$^{-1}$ at the
  default \kw{pdrZeta}, so this happens where the Solomon rate is shielded
  below a few $10^{-16}$\,s$^{-1}$ (Section~\ref{sec:balance}).
\item \textbf{The coupling is a fixed-point iteration}: shielding columns, level
  populations of other cells, the chemistry's channels in \rtn{ionBalance} and
  the heating terms held in the temperature search are one or two iterations
  behind (Section~\ref{sec:lags}). The lags vanish at a fixed point; because
  the MC field is redrawn at every iteration, the fixed point is statistical,
  and the noise biases cell means through the non-linear rates
  (Section~\ref{sec:mcnoise}).
\item \textbf{Homogeneous cells.} Every quantity is uniform within a cell, and
  the MC field is the cell average. A front thinner than a cell is not
  resolved: on \texttt{hiipdr}'s $31^3$ grid a cell is 0.65\,pc and 0.104\,mag of
  $\AV$ (\dlog{H-026}), so the H/\Htwo{} transition is located only to a cell,
  and an ionisation-front cell spreads over its whole volume the warm, partly
  ionised gas of a layer about 0.18\,pc thick, which then dominates
  [O\,\textsc{i}] 63\,\textmu m (\dlog{L-37}). The shielding columns of the chemistry run to the near face of
  the cell (Section~\ref{sec:avradial}), which under-shields CO in the cells of
  its transition by up to half a cell's column.
\end{enumerate}

\subsection{Radiation field}
\begin{enumerate}\setcounter{enumi}{5}
\item \textbf{Point sources only.} The field is that of the sources of the
  input and their reprocessing; no interstellar radiation field illuminates the
  outer boundary and the cosmic microwave background is not transported
  (Sections~\ref{sec:field}, \ref{sec:Td}). Where the attenuated field of the
  sources falls below the interstellar one, the model's outer gas is more
  molecular and colder, and its grains colder, than the same cloud immersed in
  the interstellar field would be. On \texttt{hiipdr} the source's field falls
  below $\chi=1$ beyond about 0.6 of the outer radius and to a median
  $\chi=0.08$ in the outermost shell; on \texttt{pn150k} it stays above 1.6
  everywhere (\src{output/\allowbreak{}pdrFUV.out} of the fixtures' runs). The grain
  temperature is floored at $T_{\mathrm{CMB}}$ for the chemistry only
  (\dlog{L-30}); the line background of the coolants does include the CMB
  (\dlog{D-377}).
\item \textbf{The MC field is a continuum field.} It contains dust absorption
  and scattering and the gas continuum opacities (including C\,\textsc{i} and
  \Htwo{} photoionisation), but no \Htwo, CO or H\,\textsc{i} Lyman-series
  \emph{line} opacity (Section~\ref{sec:eta}). The photons that the \Htwo{} lines
  absorb remain in the transported field, where they heat dust, eject
  photoelectrons and photoionise C\,\textsc{i}. Behind the H/\Htwo{} transition
  this overstates dust heating and photoelectric emission slightly, since the
  saturated line cores hold a small part of the band's energy (the whole window
  holds 15\% of a Draine-shaped field's), and C\,\textsc{i} photoionisation, which
  happens only in the window, by more: a few tens of per cent at $\AV\sim3$ by
  the project's estimate (Section~\ref{sec:limitations}). Unmeasured.
\item \textbf{Separability of continuum and line attenuation.} The rate seen by
  a molecule is the MC continuum at the line times a line transmission computed
  from column densities along a characteristic path: the straight line to each
  source for the direct beam, and the two-leg path through a scattering site
  for diffuse light. Dust is never applied twice: $\AV=0$ on every ray
  (Section~\ref{sec:rays}). Exact for the direct beam; for the diffuse light it
  assumes that a photon's history is represented by the one site it is
  assigned to.
\item \textbf{Direct beam geometry.} The direct field arrives along the straight
  line from each point source. With several sources the columns and the
  in-cell chord are averages weighted by $r_s^{-2}$ alone, without the sources'
  luminosities or attenuation, applied to the summed direct field
  (Section~\ref{sec:avradial}). Exact for one source, as in all the fixtures;
  for several, the error is unmeasured and has no definite sign.
\item \textbf{Diffuse field geometry.} Its total is the MC's $\chiF$; its
  distribution over twelve fixed directions and over sites along them is the
  formal solution with an isotropic emissivity $\sum\kappa^{\mathrm{sca}}J$ and a
  grey FUV extinction $\RFUV\kappa_V$ (Section~\ref{sec:sites}). The same window
  factor is applied at every site. Measured: the formal solution recovers
  52--130\% of $\chiF$, falling with depth (\dlog{L-4}); the forward-peaked phase
  function ($g=0.29$--0.61 on the fixtures) puts more of the real diffuse field
  on the source side than the weights do.
\item \textbf{Angular quadrature}: twelve HEALPix $N_{\mathrm{side}}=1$
  directions fixed in the grid frame, for the sites and for line escape
  (Sections~\ref{sec:healpix}, \ref{sec:lineRays}). Biased along the ring axis:
  cells nearest the plane $z=0$ hold 5--11\% more \Htwo{} before and across the
  transition than their symmetry-equivalent cells; 48 directions remove this,
  moving the totals by under 1\% (\dlog{L-53}, \dlog{D-410}).
\item \textbf{Spectral shape.} The generic photoreactions assume the Draine
  spectral shape across 6--13.6\,eV (Section~\ref{sec:generic}), while the MC
  field's LW share falls from 0.30 to 0.008 through \texttt{hiipdr}'s PDR
  (\dlog{D-349}); CO takes the ratio of the MC field's LW share to Draine's,
  with the shape inside the window assumed Draine's (Section~\ref{sec:window});
  \Htwo{} uses the MC specific intensity at each line, constant across the
  2\%-wide frequency bin that contains it (Section~\ref{sec:eta}).
\item \textbf{Symmetry.} For a \kw{symmetricXYZ} octant the lower faces are
  mirrors for every walk to the edge of the model; the walks towards the
  sources do not reflect (Section~\ref{sec:raycells}). Valid for a model
  symmetric in the three planes with its sources on them.
\end{enumerate}

\subsection{\texorpdfstring{\Htwo}{H2} microphysics}
\begin{enumerate}\setcounter{enumi}{13}
\item \textbf{Line profile}: Voigt, with Doppler parameter
  $b^2=2kT_e/m_{\mathrm{H_2}}+v_{\mathrm{turb}}^2$ and damping from the
  upper level's total radiative decay rate (to the bound levels and the
  continuum of $X$; the data contain no predissociated levels)
  (Section~\ref{sec:shield}). Damping each line by its own $A$, as Cloudy does,
  would put \texttt{hiipdr}'s H/\Htwo{} crossing 0.21\,mag and \texttt{pn150k}'s
  0.10\,mag shallower (\dlog{H-026}).
\item \textbf{Self-shielding} of each line by the column of its lower level,
  averaged through the cell's own column with the curve of growth; the cell's
  own column is taken over its mean chord along the source direction (direct)
  or over all directions (diffuse) (Section~\ref{sec:ownavg}).
\item \textbf{Line overlap} by the other \Htwo{} lines, as a six-point mean of
  their transmission over the line's absorption; photons absorbed in another
  line's wing are lost, not counted as pumping it; H\,\textsc{i} Lyman lines do
  not shade (Section~\ref{sec:overlap}).
\item \textbf{Trapping}: complete redistribution; a trapped photon re-excites
  the same upper level (on-the-spot); escape probability per line, without
  overlap, averaged over twelve directions (Section~\ref{sec:escape}).
\item \textbf{Colliders}: H (including reactive collisions), He, ortho- and
  para-\Htwo{} and \Hp; no electrons, so the excitation of \Htwo{} in ionised
  gas is underestimated. Missing same-spin pairs (86--99.9\% of them, depending
  on the collider) take g-bar estimates; missing spin-changing pairs get no
  collisional rate; outside the tabulated temperatures (for H, 100--5000\,K) the
  rates are held at the table's end values (Section~\ref{sec:coll}).
\item \textbf{Formation distribution} of \citet{TakahashiUehara2001},
  silicate and carbonaceous surfaces mixed in the proportion of the formation
  rates (Section~\ref{sec:formdist}). Formation through H$^-$ and \Htwo$^+$,
  cosmic-ray excitation and destruction other than Solomon are not in the
  level model.
\end{enumerate}

\subsection{Grains}
\begin{enumerate}\setcounter{enumi}{19}
\item \textbf{\Htwo{} formation} follows
  \citet{CazauxTielens2002,CazauxTielens2004} with the
  \citet{HollenbachMcKee1979} sticking coefficient and 3D-PDR's surface
  energies per material, on the geometric grain area per H nucleus of the
  model's own grains, split into silicate and carbonaceous by the species'
  labels (Sections~\ref{sec:grainarea}, \ref{sec:formation}). On
  \texttt{hiipdr}'s silicates this gives 0.50--0.91 of Cloudy's coefficient,
  mainly through the surface energies and partly through the colder gas; on
  \texttt{pn150k}'s graphite 0.86--1.00 (\dlog{L-33}).
\item \textbf{One grain temperature per cell}: the number-weighted mean over
  species and sizes, the size weights normalised over the whole size grid,
  enters an efficiency that is steep in $T_d$; a size grid wider than a
  species' sizes pulls the mean towards zero and can switch formation off
  (\dlog{H-035}; Section~\ref{sec:Td}).
\item \textbf{Grain temperature} from MOCASSIN's radiative equilibrium with no
  stochastic heating by default (\dlog{L-8}), floored at $T_{\mathrm{CMB}}$ because
  the CMB is not transported (Section~\ref{sec:Td}).
\item \textbf{Extinction ratio} $\RFUV$ weighted by source~1's emitted
  spectrum; it overstates the band extinction of the reddened local field at
  depth, which matters only for the diffuse-site weights (Section~\ref{sec:dust}).
\item \textbf{Photoelectric heating} from the model's grains alone: no PAHs or
  very small grains (\dlog{L-7}), so the heating is lower than PDR models with
  the \citet{BakesTielens1994} prescription would have at the same field, and
  the gas colder (Section~\ref{sec:thermal}).
\end{enumerate}

\subsection{Ownership and coupling}
\begin{enumerate}\setcounter{enumi}{24}
\item \textbf{MOCASSIN owns the ionisation} of H, He, C, O and Mg and the
  electron density: imposed on the chemistry before and after the solve, with
  the photoionisation and recombination rates of these atoms taken from
  MOCASSIN inside the solve (Sections~\ref{sec:linked}, \ref{sec:impose}).
  In particular the \Cp/C ratio, which controls the carbon available to CO, is
  MOCASSIN's, and its C\,\textsc{i} photoionisation has no \Htwo-line shielding.
\item \textbf{MOCASSIN owns the temperatures}: the chemistry and the level
  model run at $T_e$ (clamped to the fitted range of the rate data, while 3D-PDR's
  heating terms are evaluated at the unclamped $T_e$, \dlog{L-30}) and $T_d$
  (Section~\ref{sec:chemistry}). The thermal balance is solved at fixed
  chemistry, with the extra heating terms held fixed during the temperature
  search (Section~\ref{sec:thermal}).
\item \textbf{The chemistry owns the molecular content}, which reaches MOCASSIN
  only as the fraction of each element locked in molecules, as $n(\mathrm{H_2})$
  in the opacity, and through the coolants and heating terms
  (Section~\ref{sec:feedback}).
\item \textbf{Cosmic rays}: a uniform ionisation rate (\kw{pdrZeta}), with no
  attenuation with column (Section~\ref{sec:cr}); their heating, 9.4\,eV per
  \Htwo{} ionisation, is counted in molecular gas only
  \citep{TielensHollenbach1985}.
\end{enumerate}

\subsection{Relation to the standard PDR benchmark}

The PDR code comparison \citep{Rollig2007} fixed, for its benchmark models,
the field in Draine units with $\chi=1.71\,G_0$, $\AV=6.289\times10^{-22}N_{\mathrm H}$,
an FUV attenuation $e^{-3.02\AV}$ wherever a code did not treat it as a
function of wavelength, an \Htwo{} formation coefficient
$3\times10^{-18}T^{1/2}$\,cm$^3$\,s$^{-1}$, a cosmic-ray ionisation rate of
$5\times10^{-17}$\,s$^{-1}$, and photoelectric heating after
\citet{BakesTielens1994}, the only heating process it standardised; the
\Htwo{} shielding was left to each code (their Sect.~4.2.3 and Table~5). The
present model departs from all of these choices by construction: its field is
the transported MC field, scattering included (Section~\ref{sec:split}); its
$\AV/N_{\mathrm H}$, $\RFUV$ and grain area are its own grains'
(Table~\ref{tab:dustfix}); its formation coefficient follows the grains'
temperature (Section~\ref{sec:Td}); $\zeta$ defaults to $10^{-16}$\,s$^{-1}$; its
photoelectric heating comes from its own grains without PAHs; and its \Htwo{}
shielding is the level model's. A benchmark comparison therefore tests the
coupling only after these inputs have been matched, as was done against the
vendored 3D-PDR in \dlog{D-257} (in $\zeta$, composition, $\AV$ conversion and
dust); unmatched, it measures differences of physics rather than of numerics.

% ======================================================================
\section{Known limitations and open issues}
\label{sec:limitations}
% ======================================================================

This section lists the limitations that bear on the path from the field to
the molecules, with the direction of the bias each introduces and its
measured size where one exists. The first group is recorded in the project's
own logs. The second group was found while reading the code for this document
and is not, as far as could be found, recorded in \src{DEVLOG.md},
\src{DEVLOG-H2.md} or \src{doc/\allowbreak{}pdr-status.md}; its sizes have not been
measured.

\subsection{Recorded in the project's logs}

\begin{enumerate}
\item \textbf{Per-cell Monte Carlo noise in the photorates} (\dlog{L-26},
  \dlog{D-405}, \dlog{H-012}). Every photorate reads its own cell's field, whose
  per-cell scatter reaches 14--22\% beyond $\AV\sim2$ at the default packet
  count; the level model's per-line field is noisier than the band integral.
  Besides scatter, the noise biases cell means through the non-linear rates
  (estimated $+5\%$ at 22\% scatter) and, at the carbon turn, decides which of
  two states a cell is in (\dlog{D-405}). An estimator accumulated over
  iterations is registered for future development (\dlog{D-411}).
  Section~\ref{sec:mcnoise} gives the numbers.
\item \textbf{Packet convergence of carbon and CO} (\dlog{D-408}, \dlog{H-032}).
  At the default $6\times10^6$ packets the \Htwo{} lines and the H/\Htwo{}
  crossing are converged, but the carbon turn is placed too shallow (C$^+=$C
  moves by 0.5\,mag on \texttt{hiipdr} up to $3.84\times10^8$ packets, towards
  Cloudy's), the deep gas is too cold, [C\,\textsc{i}] needs 16--64 times the
  packets and CO is still moving at 64 times. The stop test does not see this.
\item \textbf{The twelve-direction quadrature} (\dlog{L-53}, \dlog{D-410}):
  cells that the grid's symmetry makes equal differ in $n(\mathrm{H_2})$ by up
  to about 11\% following the ring axis of the twelve directions, independently
  of the packet count; 48 directions remove the pattern at 1.48 times the run
  time, moving totals by under 1\%. The twelve were kept by decision, since no
  line map carries the pattern (CO's abundance and the \Htwo{} density do).
\item \textbf{Damping width of the \Htwo{} lines} (\dlog{H-008}, \dlog{H-026}).
  The code damps each line with the upper level's total radiative decay rate
  (to the bound levels and the continuum of $X$), which is the physical natural
  width; damping each line
  with its own $A$, as Cloudy does, narrows the wings, strengthens the
  shielding and moves the H/\Htwo{} crossing 0.21\,mag shallower on \texttt{hiipdr}
  (from $\AV=1.396$ to 1.188, Cloudy's being 1.280) and 0.10\,mag shallower on
  \texttt{pn150k}. Between $\AV=1.05$ and 1.2 on \texttt{hiipdr} the model's
  photodissociation rate per molecule is 5--7 times Cloudy's. On \texttt{hiipdr}
  the width accounts for $+0.21$\,mag of the difference between the two codes'
  crossings and everything else for $-0.09$\,mag (\dlog{H-026}); it is the
  largest single identified term and should accompany any comparison with
  Cloudy, which uses the total decay rate for every species but \Htwo{}
  (\dlog{H-008}).
\item \textbf{No H\,\textsc{i} Lyman-series shading} of the \Htwo{} band
  (\src{DEVLOG-H2.md}, \dlog{H-006}, ``What is not yet known''). Neither the MC transfer
  nor the level model contains it, so pumping and dissociation are
  overestimated where a large H\,\textsc{i} column lies in front of the molecular
  gas, and the transition is placed too deep. The hydrogen column between the
  ionisation front and the H/\Htwo{} crossing is about $1.8\times10^{21}$\,cm$^{-2}$ on
  \texttt{hiipdr} ($\AV$ 0.48 to 1.40 at $1.93\times10^{21}$\,cm$^{-2}$\,mag$^{-1}$,
  Table~\ref{tab:dustfix}; \dlog{D-260}, \dlog{H-027}) and at most
  $1.1\times10^{21}$\,cm$^{-2}$ on \texttt{pn150k} (crossing at $\AV=0.39$,
  $2.79\times10^{21}$\,cm$^{-2}$\,mag$^{-1}$), part of it molecular. The damping wings
  of Ly$\beta$, Ly$\gamma$ and Ly$\delta$ alone cover about 4, 1.4 and 0.7\,\AA{} of the
  912--1110\,\AA{} band at $N(\mathrm{H\,I})=10^{21}$\,cm$^{-2}$, growing as $N^{1/2}$
  (an estimate, Section~\ref{sec:overlap}). Unmeasured.
\item \textbf{No cross-pumping and no overlap in trapping} (\dlog{H-006}): a
  photon absorbed in another line's wing is lost rather than pumping that
  line, and the escape of each line ignores the others.
\item \textbf{Electrons are not colliders} of the \Htwo{} levels, and the level
  model has no destruction other than Solomon (\src{h2\_mod.f90} header); this
  matters in the partly ionised gas, where \Htwo{} is already 14--230 times
  Cloudy's for chemical reasons (\dlog{L-43}).
\item \textbf{C\,\textsc{i} photoionisation without \Htwo-line shielding}
  (\src{DEVLOG.md}, item 8 of the audit list near line 28742): MOCASSIN's rate
  replaces \rtn{CIPDRATE}, and MOCASSIN does not transport LW lines, so the
  rate is overestimated behind the H/\Htwo{} transition; estimated at a few
  tens of per cent at $\AV\sim3$, with the opposite sign to the deep \Cp{}
  deficit (\dlog{L-3}).
\item \textbf{Chlorine is absent from the network}; in Cloudy it carries 25--50\%
  of the \Htwo{} sink from $\AV=0.26$ on in the PN model, and Cloudy without
  chlorine turns molecular 0.042\,mag earlier (\dlog{L-46}, \dlog{H-026}). Its
  absence moves the model's crossing shallower, by about 0.04\,mag on
  \texttt{pn150k}, partly offsetting there the 0.10\,mag of item~4 (\dlog{L-46}).
\item \textbf{Silicate formation efficiency}: the per-material Cazaux \& Tielens
  surface energies of 3D-PDR give 0.50--0.91 of Cloudy's coefficient on
  \texttt{hiipdr} (\dlog{L-33}, \dlog{H-022}); reported, not tuned.
\item \textbf{The mean grain temperature can switch formation off}
  (\dlog{H-035}): the number-weighted mean over the whole size grid counts the
  bins outside a species' size range at zero temperature; on a grid wider
  than the grains present this put $T_d$ at $10^{-6}$\,K and no \Htwo{} formed.
  Fixed after \texttt{8fc1512}: \dlog{D-414} takes the mean over the bins that hold
  grains (Section~\ref{sec:scope}).
\item \textbf{Diffuse-site weights} recover 52--130\% of the MC's diffuse field
  (\dlog{L-4}); the magnitude is the MC's, the directional distribution, and so
  the columns assigned to the diffuse light, is approximate
  (Section~\ref{sec:sites}).
\item \textbf{Unresolved fronts} (\dlog{L-37}, \dlog{D-088}): on uniform grids
  the ionisation front and the H/\Htwo{} transition are thinner than a cell; the
  front cells, which carry the warm partly ionised gas through their whole
  volume, make 92\% of \texttt{hiipdr}'s [O\,\textsc{i}] 63\,\textmu m, 9.3--9.5 times
  Cloudy's at $31^3$ and 5.5 times at $61^3$. Front-sensitive lines need a
  resolution check; a sub-cell estimate of [O\,\textsc{i}] 6300\,\AA{} is available
  offline (\kw{frontLines}).
\item \textbf{The gas re-emits from the hydrogen and helium spectrum whatever
  absorbed} (\dlog{L-51}). A packet absorbed by neutral carbon in the PDR is
  re-emitted, some of the time, as a hydrogen- or helium-ionising photon:
  1.9\% of \texttt{hiipdr}'s cells per iteration at $\AV=2.15$--2.43 are
  over-ionised, and their chemistry, solved to steady state on that iteration's
  ionisation, can lose more than half its \Htwo{} (202 such drops in six
  iterations at the default packet count, 23 at four times, none at sixteen
  times).
\item \textbf{No PAHs or very small grains} (\dlog{L-7}): the photoelectric
  heating lacks the component that dominates in real PDRs, so the gas is colder
  than in PDR models with that component at the same field; common-mode with
  the Cloudy references.
\item \textbf{Gas--grain collisional coupling omits \Htwo} (\dlog{L-9}): 68
  times weaker than Cloudy's dust cooling at $\AV=0.65$ on \texttt{pn150k},
  negligible at $n_{\mathrm H}=10^3$\,cm$^{-3}$ and decisive at high density.
\item \textbf{The \Htwo{} line emission is computed optically thin and without
  dust extinction} in \src{lineFluxH2.out} (\src{h2Lines\_mod.f90} header).
\end{enumerate}

\subsection{Found while preparing this document}

\begin{enumerate}\setcounter{enumi}{17}
\item \textbf{Formation flux in the level model is tied to the Solomon flux.}
  \rtn{solveCell} calls \rtn{h2Solve} without a sink or a source
  (\src{pdr\_mod.f90:\allowbreak{}5904}), so the rate at which new molecules enter the
  levels in the formation distribution $\phi$ equals the Solomon dissociation
  rate. The chemistry's actual formation rate, $R_{\mathrm f}n_{\mathrm H}n(\mathrm H)$,
  balances all destruction (cosmic rays, ion--molecule reactions, Solomon) in
  steady state. Where the Solomon rate has been shielded away, formation
  pumping is therefore underestimated. With the network's \Htwo\,+\,CRP channels
  alone the other destruction is $0.99\,\zeta$, $10^{-16}$\,s$^{-1}$ at the default
  \kw{pdrZeta}, so this happens where the Solomon rate is shielded below a few
  $10^{-16}$\,s$^{-1}$ (Section~\ref{sec:balance}). The energy budget is split the same
  way: the internal energy that newly formed molecules bring to the levels
  (and that the collisional term returns to the gas) follows the Solomon flux,
  while the kinetic part of the formation energy (\texttt{heat(5)}) follows the
  chemistry's formation rate. Both numbers needed for a consistent closure are
  already available in \rtn{solveCell} ($w_{\mathrm{sil}}+w_{\mathrm{gra}}$,
  $n(\mathrm H)$, $n(\mathrm{H_2})$): passing the non-Solomon destruction as
  \rtn{sink}, or the chemistry's formation flux per molecule as \rtn{src}, would
  make the two agree. A per-cell print of formation flux over Solomon flux
  would size the effect first.
\item \textbf{The final output mixes two Solomon rates.} The last
  \rtn{h2LevelModel(grid,12)} (\src{iteration\_mod.f90:\allowbreak{}3254, 3283}) runs after
  the last chemistry, which is not re-solved. The written level populations and
  $k_{\mathrm{Sol}}$ are those of up to twelve sweeps; the written
  $n(\mathrm{H_2})$ was obtained with the previous one-sweep rate. The size of
  the difference is printed by the call itself (``Solomon rate against the
  last call's'').
\item \textbf{The \Htwo{} collisional heating term is frozen and two
  iterations old in the temperature search.} $\mathcal{H}_{\mathrm{coll}}$ changes
  sign and depends steeply on $T$ (it is \Htwo's line cooling where the gas is
  warmer than the levels), but it enters \rtn{pdrHeatXtra}, which MOCASSIN holds
  constant in \rtn{iterateT}, and the value used to find $T_e^{(n)}$ was
  computed by the level model at $T_e^{(n-2)}$ (Section~\ref{sec:lags}).
  \dlog{H-030} associates this term, together with the damping, with the two
  derailed iterations at the start of the chemistry. Tabulating it against
  $T$, as the CO cooling is (\dlog{D-215}), would remove the freeze.
\item \textbf{Generic photoreactions keep the Draine spectral shape}
  (Section~\ref{sec:generic}). \dlog{L-24} corrected the LW-window reactions;
  the other 24 photoreactions still multiply the 6--13.6\,eV band strength by
  coefficients computed for the Draine shape, although the MC field's shape
  changes strongly with depth (its LW share falls from 0.30 to 0.008 on
  \texttt{hiipdr}, \dlog{D-349}). Reactions whose thresholds lie in the window
  are overestimated at depth, those driven by the 6--11\,eV part underestimated.
  Folding each reaction's cross section, or a few sub-bands, over the MC field
  would remove the assumption.
\item \textbf{The level model's \Htwo{} transmission is not shared.} The
  per-line \Htwo{} shielding computed in \rtn{solveCell} is exactly what is
  missing from MOCASSIN's C\,\textsc{i} rate (item~8) and what CO's
  van Dishoeck \& Black factor approximates for \Htwo{} screening; neither uses
  it.
\item \textbf{CO shielding and the chemistry's own \Htwo{} diagnostic} exclude
  the target cell's own column (\dlog{D-088}); only the level model averages
  through it. On \texttt{hiipdr}'s grid half a cell can hold $10^{16}$\,cm$^{-2}$ of
  CO, so CO is under-shielded in the cells of its own transition and the
  transition is placed up to about half a cell too deep
  (Section~\ref{sec:avradial}).
\item \textbf{Several sources are combined by geometry alone.} The effective
  $\AV$, the source-ward columns, the in-cell chord (\rtn{computeAVradial},
  \src{pdr\_mod.f90:\allowbreak{}1986}) and the level model's source-ward level columns
  (\rtn{starColumns}, \src{:\allowbreak{}5651}) weight the sources by $r_s^{-2}$, without
  their luminosities or their attenuation, and apply the mean column to the
  summed direct field through non-linear shielding functions. Exact for one
  source, which is every fixture; untested for several. Weighting each source
  by its own attenuated direct field, $L_{\mathrm{FUV},s}r_s^{-2}e^{-\RFUV\tau_{V,s}}$,
  and evaluating the shielding per source would remove it.
\item \textbf{$\RFUV$ is weighted by the emitted spectrum of source~1}
  (\src{pdr\_mod.f90:\allowbreak{}1381--1390}). It overstates the band extinction of the
  local, reddened field at depth, and so the optical depth between the
  scattering sites and the target; this suppresses distant sites increasingly
  with depth and is a candidate for the fall with depth of the formal
  solution's recovery in \dlog{L-4}.
\item \textbf{The grain temperature is a number-weighted mean applied to an
  area-weighted rate} (Section~\ref{sec:Td}). Besides the zero-bin trap of
  item~11, the mean favours the smallest grains and one efficiency
  $\epsilon(\langle T_d\rangle)$ replaces $\langle\epsilon(T_d)\rangle$ over the sizes. On
  silicate at 20--50\,K, where $\epsilon$ falls by a factor of 3.6
  (Section~\ref{sec:Td}), the formation rate depends on how the temperatures
  of the sizes are spread. Evaluating the rate per size with $\pi a^2$ weights
  would remove both; a temperature per grain material is registered for
  future development as \dlog{D-417} (\dlog{L-54}), after the commit described
  here.
\item \textbf{No interstellar field at the outer boundary}
  (Section~\ref{sec:assumptions}, item~6). On \texttt{hiipdr} the source's field
  is below $\chi=1$ in the outer 40\% of the radius, so the outer PDR is more
  molecular and colder than the same cloud in the interstellar field.
\item \textbf{The coolant line escape depends weakly on the sweep order.} The
  optical depths of the coolant lines read the other cells' level populations
  as stored on the rank (\src{pdr\_mod.f90:\allowbreak{}4315--4316}), which cells solved
  earlier on the same rank have already updated
  (\src{:\allowbreak{}4440--4442}); the shielding columns are snapshots and do not have
  this dependence.
\item \textbf{The seed double-attenuates.} \rtn{seedMolecules} multiplies the
  local, already attenuated $\chi$ by \rtn{SCATTER}$(\AV)$
  (\src{pdr\_mod.f90:\allowbreak{}2937--2938}; Section~\ref{sec:seed}). It affects only the
  starting point of the first PDR iteration.
\item \textbf{The plane-parallel branch includes the near half of the cell in
  its columns} (\src{pdr\_mod.f90:\allowbreak{}2492--2507}), against \dlog{D-088}'s rule in
  the point-source branch. The level model does not run in that geometry.
\item \textbf{Collisional ionisation of hydrogen lacks the electron density.}
  MOCASSIN adds $2.75\times10^{-16}T^{3/2}(157\,894/T+2)e^{-157\,894/T}$, a rate
  coefficient in cm$^3$\,s$^{-1}$, to the photoionisation rate in s$^{-1}$ without the
  factor $n_e$, in its ionisation balance (\src{update\_mod.f90:\allowbreak{}2769--2770})
  and in the rate handed to the chemistry (\src{:\allowbreak{}2859}). Collisional
  ionisation is therefore misscaled by a factor $n_e$ (underestimated wherever
  $n_e>1$\,cm$^{-3}$); with the factor, the term would reach 10\% of a
  photoionisation rate of $10^{-8}$\,s$^{-1}$ at about $1.5\times10^4$\,K for
  $n_e=10^4$\,cm$^{-3}$, so hot dense ionised gas is affected and the PDR is not.
  Inherited from MOCASSIN (Section~\ref{sec:linked}).
\item \textbf{Two reactions of the rate file are inactive.} The grain-assisted
  recombinations of He$^+$ and C$^+$ (reactions 330 and 331) have no term in
  \src{pdr/\allowbreak{}odes\_reduced.c}; their coefficients are computed but unused
  (Section~\ref{sec:network}).
\item \textbf{H\,+\,\Htwo{} rates below 100\,K} are those of 100\,K, the first
  temperature of the \citet{Lique2015} table, which overestimates the reactive
  ortho--para conversion by H in cold gas, where \Hp{} exchange usually
  dominates it (Section~\ref{sec:coll}).
\item \textbf{The \Htwo{} line profile is normalised to its own centre}
  ($\tau_0H(a,x)/H(a,0)$ with Doppler-core $\tau_0$), overstating the Voigt optical
  depth by $1/H(a,0)$, at most 3\% for these data (Section~\ref{sec:shield}).
\item \textbf{Stale or inaccurate comments.} \src{constants\_mod.f90:\allowbreak{}18--20} labels
  the FUV lower edge 2400\,\AA{} (the value is 6.0\,eV, 2068\,\AA); the comment in
  \rtn{addOpacity} (\src{ionization\_mod.f90:\allowbreak{}444--453}) describes a power-law
  \Htwo{} cross section, whereas \rtn{yanH2xSec} is the piecewise fit of
  \citet{Yan1998}; and the comment of \rtn{draineWindowShare}
  (\src{pdr\_mod.f90:\allowbreak{}2285--2286}), \dlog{D-349} and \dlog{L-24} give the Draine
  field's LW share of the continuous band as 0.1516 (0.152), whereas
  Equation~(\ref{eq:draine}) gives 0.148. The code uses the binned value
  (0.151--0.164 on the fixtures), consistently with the MC sums, so no result
  is affected. The comment of \rtn{h2FormDist} (\src{h2\_mod.f90:\allowbreak{}578--579})
  gives the vibrational distribution as ``their eqs 4-5'' about
  $E_{vm}=X_{\mathrm{vib}}D$ and the rotational one as ``their eqs 6 and 7'', whereas
  in \citet{TakahashiUehara2001} the vibrational distribution is eqs.~5--7 (with
  $E_{vm}=X_{\mathrm{vib}}D+E_{\mathrm{off}}$) and the rotational models are eqs.~12
  and 13 (the code is right); and \src{update\_mod.f90:\allowbreak{}2621} attributes the
  collisional-ionisation formula to \citet{DrakeUlrich1980}, which does not
  contain it (Section~\ref{sec:linked}).
\end{enumerate}


\appendix
% ======================================================================
\section{Index of code locations}
\label{app:index}
% ======================================================================

Files are under \src{source/\allowbreak{}} of the \texttt{h2-model} worktree at commit
\texttt{8fc1512}; \src{pdr/\allowbreak{}} is the vendored 3D-PDR. A range is the whole
routine unless the process column names a part of it.

\begin{longtable}{@{}L{0.39\linewidth}L{0.25\linewidth}L{0.30\linewidth}@{}}
\caption{Index of code locations.}\label{tab:index}\\
\toprule
Process & Routine & File:lines \\
\midrule
\endfirsthead
\toprule
Process & Routine & File:lines \\
\midrule
\endhead
\multicolumn{3}{@{}l}{\emph{Iteration and transfer}}\\
MC iteration, PDR block & \rtn{iterateMC} & \src{iteration\_mod.f90:\allowbreak{}56--3345} \\
packet transport & \rtn{energyPacketDriver}, \rtn{pathSegment} & \src{photon\_mod.f90:\allowbreak{}37--3575, 1249--3344} \\
path-length estimator ($\Jste$, $\Jdif$) & \rtn{pathSegment} & \src{photon\_mod.f90:\allowbreak{}1900--1924, 2236--2250} \\
scattering, re-flag of the packet as diffuse & \rtn{pathSegment} & \src{photon\_mod.f90:\allowbreak{}2025--2035, 2154--2165} \\
absorption by dust or gas & \rtn{pathSegment} & \src{photon\_mod.f90:\allowbreak{}2050--2103} \\
$\Jdif$ always accumulated in PDR runs & \rtn{readInput} & \src{set\_input\_mod.f90:\allowbreak{}775--778} \\
estimator units, octant & \rtn{iterateMC} & \src{iteration\_mod.f90:\allowbreak{}924, 940} \\
FUV band, LW window, Draine unit, $V$ band & constants & \src{constants\_mod.f90:\allowbreak{}16--35} \\
FUV and $V$ bin indices & \rtn{setContinuum} & \src{continuum\_mod.f90:\allowbreak{}92--95} \\
\multicolumn{3}{@{}l}{\emph{Dust}}\\
grain area per H by material & \rtn{measureH2grainArea} & \src{pdr\_mod.f90:\allowbreak{}1266--1338} \\
material from the grain label & \rtn{isCarbonGrain}, \rtn{isSilicateGrain} & \src{pdr\_mod.f90:\allowbreak{}1228--1247, 1341--1360} \\
$\RFUV$, $\AV/N_{\mathrm H}$, area per cell & \rtn{dustPerCell} & \src{pdr\_mod.f90:\allowbreak{}1363--1436} \\
model-wide fallbacks, own extinction curve & \rtn{calibratePDRdust} & \src{pdr\_mod.f90:\allowbreak{}1439--1662} \\
grain LW albedo & \rtn{grainAlbedoLW} & \src{update\_mod.f90:\allowbreak{}4658--4673} \\
grain temperatures; their mean & \rtn{getDustT} & \src{update\_mod.f90:\allowbreak{}4282--4562; 4546--4555} \\
size weights of the distribution & \rtn{makeDustXsec} & \src{ph\_mod.f90:\allowbreak{}999--1012} \\
grain attenuation factor, extinction curve & \rtn{SCATTER}, \rtn{XLAMBDA} & \src{pdr/\allowbreak{}shield.F90:\allowbreak{}449--497, 525--549} \\
\multicolumn{3}{@{}l}{\emph{Geometry}}\\
exact ray--cell traversal & \rtn{rayCells} & \src{pdr\_mod.f90:\allowbreak{}1710--1833} \\
source-ward $\AV$ and columns & \rtn{computeAVradial} & \src{pdr\_mod.f90:\allowbreak{}1859--2015} \\
dispatch, plane-parallel case & \rtn{computeAVdust} & \src{pdr\_mod.f90:\allowbreak{}2410--2523} \\
diffuse snapshot & \rtn{snapDiffuseSources} & \src{pdr\_mod.f90:\allowbreak{}2082--2121} \\
scattering sites, two-leg columns & \rtn{traceDiffuseSites} & \src{pdr\_mod.f90:\allowbreak{}2136--2276} \\
Draine LW share & \rtn{draineWindowShare} & \src{pdr\_mod.f90:\allowbreak{}2288--2314} \\
line-escape rays & \rtn{traceLineRays} & \src{pdr\_mod.f90:\allowbreak{}2327--2376} \\
$\chi$, $\chi_0$ & \rtn{mocassinToPDR} & \src{pdr\_mod.f90:\allowbreak{}2538--2677} \\
\multicolumn{3}{@{}l}{\emph{\Htwo{} level model}}\\
driver, sweeps, MPI scheduling & \rtn{h2LevelModel} & \src{pdr\_mod.f90:\allowbreak{}5201--6014} \\
source-ward level columns & \rtn{starColumns} & \src{pdr\_mod.f90:\allowbreak{}5621--5659} \\
per-cell solution & \rtn{solveCell} & \src{pdr\_mod.f90:\allowbreak{}5665--5954} \\
warm-start populations & \rtn{h2PopResume} & \src{pdr\_mod.f90:\allowbreak{}5961--6012} \\
1D-build stub & \rtn{h2LevelModel} & \src{pdr\_mod.f90:\allowbreak{}6015--6027} \\
$X$ levels, $A$ values, collision tables & \rtn{h2ReadX} & \src{h2\_mod.f90:\allowbreak{}206--334} \\
excited states, lines, dissociation & \rtn{h2ReadElec} & \src{h2\_mod.f90:\allowbreak{}366--516} \\
UV pumping, Solomon rate, trapping & \rtn{h2PumpRates} & \src{h2\_mod.f90:\allowbreak{}540--572} \\
formation distribution & \rtn{h2FormDist} & \src{h2\_mod.f90:\allowbreak{}589--639} \\
collision-table interpolation & \rtn{tabRate} & \src{h2\_mod.f90:\allowbreak{}642--660} \\
Voigt function & \rtn{h2Voigt} & \src{h2\_mod.f90:\allowbreak{}664--686} \\
$S$, $W$, absorption quantiles & \rtn{shieldQuad}, \rtn{shieldTable} & \src{h2\_mod.f90:\allowbreak{}692--760} \\
shielding look-up, cell average & \rtn{h2Shield}, \rtn{tabShieldJ}, \rtn{h2ShieldJ} & \src{h2\_mod.f90:\allowbreak{}768--866} \\
direct-path shielding (prepare/evaluate) & \rtn{h2LineShieldPrep}, \rtn{h2LineShieldEval} & \src{h2\_mod.f90:\allowbreak{}876--907} \\
Cloudy's FGK79 fit (comparison only) & \rtn{h2ShieldFederman} & \src{h2\_mod.f90:\allowbreak{}915--934} \\
line overlap & \rtn{lineOverlap} & \src{h2\_mod.f90:\allowbreak{}1018--1289} \\
line parameters $\sigma$, $a$ & \rtn{h2LineParams} & \src{h2\_mod.f90:\allowbreak{}1328--1336} \\
escape probabilities & \rtn{h2LineEscape} & \src{h2\_mod.f90:\allowbreak{}1344--1361} \\
collisional and radiative rates & \rtn{h2CollRad} & \src{h2\_mod.f90:\allowbreak{}1369--1415} \\
collisional heating & \rtn{h2CollHeat} & \src{h2\_mod.f90:\allowbreak{}1425--1448} \\
statistical equilibrium & \rtn{balance}, \rtn{h2Solve}, \rtn{gaussSolve} & \src{h2\_mod.f90:\allowbreak{}1456--1568} \\
outputs & \rtn{h2Lines\_mod} & \src{h2Lines\_mod.f90} \\
\multicolumn{3}{@{}l}{\emph{Chemistry}}\\
network set-up, linked elements, replaced channels & \rtn{initPDRnetwork} & \src{pdr\_mod.f90:\allowbreak{}243--948} \\
atoms of each element per species & \rtn{initPDRnetwork} & \src{pdr\_mod.f90:\allowbreak{}585--603} \\
first-visit seed & \rtn{seedAbunFromMocassin} & \src{pdr\_mod.f90:\allowbreak{}3123--3157} \\
first-iteration molecular seed & \rtn{seedMolecules} & \src{pdr\_mod.f90:\allowbreak{}2855--3023} \\
one cell of chemistry & \rtn{pdrStep} & \src{pdr\_mod.f90:\allowbreak{}3180--4580} \\
rays: $\chiD$, $\chiF$, sites, window factors & \rtn{pdrStep} & \src{pdr\_mod.f90:\allowbreak{}3383--3455} \\
imposition before the solve & \rtn{pdrStep} & \src{pdr\_mod.f90:\allowbreak{}3503--3584} \\
columns handed over & \rtn{pdrStep} & \src{pdr\_mod.f90:\allowbreak{}3593--3628} \\
rate coefficients & \rtn{CALCULATE\_REACTION\_RATES} & \src{pdr/\allowbreak{}calc\_reac\_rates.F90}; call \src{pdr\_mod.f90:\allowbreak{}3692} \\
MOCASSIN's atomic rates into the network & \rtn{pdrStep} & \src{pdr\_mod.f90:\allowbreak{}3706--3722} \\
\Htwo{} photodissociation $=k_{\mathrm{Sol}}$ & \rtn{pdrStep} & \src{pdr\_mod.f90:\allowbreak{}3736} \\
charge transfer coefficients & \rtn{pdrCtCoef} & \src{common\_mod.f90:\allowbreak{}1383--1436} \\
\Htwo{} formation rate & \rtn{H2\_FORMATION\_RATE} & \src{pdr/\allowbreak{}h2\_form.F90:\allowbreak{}9--134} \\
\Htwo, CO, C\,\textsc{i} photorates and shielding & \rtn{H2PDRATE}, \rtn{COPDRATE}, \rtn{CIPDRATE}, \rtn{COSHIELD}, \rtn{SCATTER}, \rtn{LBAR} & \src{pdr/\allowbreak{}shield.F90:\allowbreak{}27--55, 83--106, 132--157, 231--281, 375--416, 449--497, 570--599} \\
steady-state integration & \rtn{calculate\_abundances} & \src{pdr/\allowbreak{}calculate\_abundances.c}; call \src{pdr\_mod.f90:\allowbreak{}3749} \\
ODE right-hand side & generated & \src{pdr/\allowbreak{}odes\_reduced.c} \\
post-solve imposition & \rtn{pdrStep} & \src{pdr\_mod.f90:\allowbreak{}3770--3784, 3801--3804} \\
channels for \rtn{ionBalance} & \rtn{pdrStep} & \src{pdr\_mod.f90:\allowbreak{}3924--4054} \\
coolant level populations, CO table & \rtn{pdrStep} & \src{pdr\_mod.f90:\allowbreak{}4144--4490} \\
escape probabilities of the coolant lines & \rtn{escape\_probability} & \src{pdr/\allowbreak{}escape\_probability.F90:\allowbreak{}121--137, 196--233} \\
heating, \Htwo{} terms from the level model & \rtn{pdrStep} & \src{pdr\_mod.f90:\allowbreak{}4499--4546} \\
3D-PDR's heating terms & \rtn{CALC\_HEATING} & \src{pdr/\allowbreak{}heatingfunctions.F90:\allowbreak{}7--375} \\
\multicolumn{3}{@{}l}{\emph{Feedback}}\\
handback & \rtn{pdrToMocassin}, \rtn{setElemAvail} & \src{pdr\_mod.f90:\allowbreak{}5038--5104, 5120--5144} \\
hydrogen budget & \rtn{rescaleIonToAvail} & \src{pdr\_mod.f90:\allowbreak{}2764--2803} \\
molecular locking in the ion solver & \rtn{ionBalance} & \src{update\_mod.f90:\allowbreak{}3130--3134, 3255--3300} \\
cosmic-ray ionisation in the ion solver & \rtn{ionBalance} & \src{update\_mod.f90:\allowbreak{}2916--2930} \\
atomic photoionisation for the chemistry & \rtn{updateCell} & \src{update\_mod.f90:\allowbreak{}303--311, 2838--2878} \\
\Htwo{} continuum opacity & \rtn{addOpacity}, \rtn{yanH2xSec} & \src{ionization\_mod.f90:\allowbreak{}385--472} \\
thermal balance, temperature search & \rtn{thermBalance}, \rtn{iterateT} & \src{update\_mod.f90:\allowbreak{}1734--2456, 609--1191} \\
grain photoelectric heating & \rtn{setPhotoelHeatCool} & \src{update\_mod.f90:\allowbreak{}1460--1556} \\
gas--grain collisional exchange & \rtn{setDustGasColl\allowbreak{}HeatCool} & \src{update\_mod.f90:\allowbreak{}1562--1731} \\
CO table and look-up & \rtn{initPDRcoolTab}, \rtn{pdrCoolAt} & \src{pdr\_mod.f90:\allowbreak{}2713--2732}; \src{common\_mod.f90:\allowbreak{}1445--1486} \\
CO cooling added in the balance & \rtn{thermBalance} & \src{update\_mod.f90:\allowbreak{}2232--2249} \\
C\,\textsc{i}, C\,\textsc{ii}, O\,\textsc{i} atoms with the chemistry's colliders & atom solve; \rtn{pdrNeutralRates} & \src{emission\_mod.f90:\allowbreak{}1861--1887}; \src{common\_mod.f90:\allowbreak{}1489--1635} \\
extra heating into the thermal balance & \rtn{thermBalance} & \src{update\_mod.f90:\allowbreak{}2391} \\
second ionisation solve, kept at the end & \rtn{ionSecondSolve}, \rtn{ionKeepSecond} & \src{iteration\_mod.f90:\allowbreak{}3361--3434, 3455--3521} \\
\multicolumn{3}{@{}l}{\emph{Convergence}}\\
integrated stop test & \rtn{cvTotalsUpdate} & \src{iteration\_mod.f90:\allowbreak{}3579--3765}; test at \src{:\allowbreak{}3243} \\
per-cell $x(\mathrm{H_2})$ test & \rtn{iterateMC} & \src{iteration\_mod.f90:\allowbreak{}2484--2512} \\
\bottomrule
\end{longtable}

% ======================================================================
\section{Input keywords and defaults on this path}
\label{app:keywords}
% ======================================================================

Keywords are read in \rtn{readInput} (\src{set\_input\_mod.f90}); defaults are
set there or in \src{common\_mod.f90}. Retired keywords that once selected
ownership or prescriptions are of two kinds (\src{set\_input\_mod.f90:\allowbreak{}731--757};
the reason for each is printed by \rtn{retiredKeyword}, \src{:\allowbreak{}1101--1172}).
Those whose choice is settled or whose value the model now measures
(\kw{pdrAVfac}, \kw{pdrUVfac}, \kw{pdrAllCells}, \kw{pdrCoolFeedback},
\kw{pdrHealpix}, \kw{pdrCoupleIter}, \kw{pdrChemIter}, \ldots) are still parsed
and ignored with a message, so old inputs run. Those that selected alternative
physics or stood in for what MOCASSIN knows (\kw{pdrZoneLimit},
\kw{pdrFUVmode}, \kw{pdrDustModel}, \kw{pdrOwnDust}, \kw{pdrPEfromPDR},
\kw{pdrMetallicity}, \kw{pdrGrainRadius}, \ldots) are refused and the run stops.

\begin{longtable}{@{}L{0.16\linewidth}L{0.15\linewidth}L{0.40\linewidth}L{0.20\linewidth}@{}}
\caption{Input keywords on the path from the field to the molecules.}\label{tab:keywords}\\
\toprule
Keyword & Default & Effect & Read / default \\
\midrule
\endfirsthead
\toprule
Keyword & Default & Effect & Read / default \\
\midrule
\endhead
\kw{pdr} & off & switches the coupling on & \src{set\_input:\allowbreak{}638} \\
\kw{pdrNetwork} & \texttt{reduced} & validated against the compiled network; cannot switch it & \src{set\_input:\allowbreak{}645}; \src{common:\allowbreak{}889} \\
\kw{pdrDataDir} & install prefix & directory of \src{dataPDR} (network, coolants, \src{h2/\allowbreak{}}) & \src{set\_input:\allowbreak{}648} \\
\kw{pdrStartIter} & 1 & first MC iteration with the PDR coupling & \src{set\_input:\allowbreak{}654}; \src{common:\allowbreak{}893} \\
\kw{pdrZeta} & $10^{-16}$\,s$^{-1}$ & cosmic-ray ionisation rate per \Htwo{} (uniform), in the network, in MOCASSIN's ion balance and in the cosmic-ray heating & \src{set\_input:\allowbreak{}651}; \src{common:\allowbreak{}902} \\
\kw{pdrEndTime} & $10^{7}$\,yr & first integration span; steady state extends it up to 128$\times$ & \src{set\_input:\allowbreak{}670}; \src{common:\allowbreak{}903} \\
\kw{pdrVturb} & 1\,km\,s$^{-1}$ & microturbulence in $b$ (\Htwo{} lines, coolant lines); velocity of the turbulent heating & \src{set\_input:\allowbreak{}673}; \src{common:\allowbreak{}907} \\
\kw{pdrTmin}, \kw{pdrTmax} & 1\,K, $2\times10^4$\,K & clamp of the chemistry's gas temperature (\src{pdr\_mod.f90:\allowbreak{}3346}) and range of the CO cooling table (\src{:\allowbreak{}2720}); the vendored routines' own lower bound is raised to $T_{\mathrm{CMB}}$ (\src{:\allowbreak{}364}) & \src{set\_input:\allowbreak{}676, 679}; \src{common:\allowbreak{}908--909} \\
\kw{pdrCoolants} & none & coolant data files (fixtures: \src{12co.dat 12c+.dat 12c.dat 16o.dat}); CO is tabulated, the atoms' collision rates and escape probabilities go to MOCASSIN's own atoms & \src{set\_input:\allowbreak{}757} \\
\kw{iterDamp} & 0 & damping of abundances, ion stages and temperature across iterations & \src{set\_input:\allowbreak{}700}; \src{common:\allowbreak{}1193} \\
\kw{teFloor} & 2.725\,K & floor of MOCASSIN's $T_e$ & \src{set\_input:\allowbreak{}691}; \src{common:\allowbreak{}1278} \\
\kw{convTotals} & 0.02 & integrated stop test & \src{set\_input:\allowbreak{}468}; \src{common:\allowbreak{}816} \\
\kw{maxIterateMC} & 30 & maximum number of MC iterations (its second number is reported, not used to stop) & \src{set\_input:\allowbreak{}413}; \src{:\allowbreak{}75} \\
\kw{nPhotons} & input & packets per iteration; sets the MC noise of every photorate (Section~\ref{sec:mcnoise}) & \src{set\_input:\allowbreak{}422} \\
\kw{randomSeed} & $-1$ (clock) & a non-negative seed makes a run repeatable at fixed rank count & \src{set\_input:\allowbreak{}403}; \src{common:\allowbreak{}734} \\
\kw{pathStretch} & off & composite path-length biasing \citep{Baes2016} & \src{set\_input:\allowbreak{}719}; \src{common:\allowbreak{}1102} \\
\kw{isotropic\allowbreak{}Scattering} & off & isotropic instead of Henyey--Greenstein dust scattering & \src{set\_input:\allowbreak{}173}; \src{:\allowbreak{}40} \\
\kw{quantumHeat\allowbreak{}Grain} & off & stochastic heating of small grains (emitted spectrum only) & \src{set\_input:\allowbreak{}200} \\
\kw{symmetricXYZ} & off & octant symmetry; mirrors in the walks & \src{set\_input:\allowbreak{}359} \\
\kw{nbins} & 600 & frequency bins of the MC field & \src{set\_input:\allowbreak{}77, 350} \\
\bottomrule
\end{longtable}

There is no keyword for the \Htwo{} level model: in the 3D build it always
runs, and its Solomon rate always replaces the network's (\dlog{H-012}).

% ======================================================================
\section{Data files on this path}
\label{app:data}
% ======================================================================

The table lists the data files that enter the path from the field to the
molecules, with the source of their content. For the \Htwo{} files the source
is the one the file's own header gives, completed where the header is not the
primary source (Section~\ref{sec:h2data}); \src{dataPDR/\allowbreak{}h2/\allowbreak{}} is a copy of
Cloudy c25.00's \src{data/\allowbreak{}h2/\allowbreak{}}, unmodified.

\begin{longtable}{@{}L{0.34\linewidth}L{0.62\linewidth}@{}}
\caption{Data files on the path from the field to the molecules.}\label{tab:datafiles}\\
\toprule
File & Content and provenance \\
\midrule
\endfirsthead
\toprule
File & Content and provenance \\
\midrule
\endhead
\src{dataPDR/\allowbreak{}species\_reduced.d} & the 33 species, initial abundances (unused: seeding is MOCASSIN's) and masses; 3D-PDR \citep{Bisbas2012} \\
\src{dataPDR/\allowbreak{}rates\_reduced.d} & 332 reactions in UMIST format: number, reactants, products, $\alpha$, $\beta$, $\gamma$, source code, $T_{\min}$, $T_{\max}$, accuracy; 331 from 3D-PDR (a subset of UMIST, \citealt{Woodall2007} or \citealt{McElroy2013}, release not recorded; Section~\ref{sec:network}); reaction 332 added in this project with the coefficient of Cloudy's \src{mole\_co\_umisthack.dat}, attributed there to \citet{Janev1987} \\
\src{dataPDR/\allowbreak{}12c+.dat}, \src{12c.dat}, \src{16o.dat}, \src{12co.dat} & coolant level data in LAMDA format \citep{Schoier2005}; for \src{16o.dat} and \src{12c.dat} the $^1D$/$^1S\to{}^3P_J$ collision rates are split by $g_J/9$ (\src{pdr/\allowbreak{}UPSTREAM.md}, \dlog{D-285}) \\
\src{dataPDR/\allowbreak{}h2/\allowbreak{}energy\_X.dat} & $X$ level energies (302 levels), \citet[their Table~1]{Komasa2011}, compiled by K.~Kaplan \\
\src{dataPDR/\allowbreak{}h2/\allowbreak{}transprob\_X.dat} & $X$ electric-quadrupole $A$ values, \citet{Wolniewicz1998}; magnetic-dipole components of 30 $v=1\to0$, $\Delta J=0$ transitions, \citet{PachuckiKomasa2011}, listed separately and summed by the reader \\
\src{dataPDR/\allowbreak{}h2/\allowbreak{}coll\_rates\_H\_15.dat} & H\,+\,\Htwo{} de-excitation rate coefficients, reactive channels included, 100--5000\,K, \citet{Lique2015} \\
\src{dataPDR/\allowbreak{}h2/\allowbreak{}coll\_rates\_He\_ORNL.dat} & He\,+\,\Htwo{}, 1--$10^4$\,K, the ORNL data set; the header cites ``Lee, T.\ G.\ et al.\ 2007, ApJ, in preparation'', and no published source for the table was identified \\
\src{dataPDR/\allowbreak{}h2/\allowbreak{}coll\_rates\_H2ortho\_ORNL.dat}, \src{coll\_rates\_H2para\_ORNL.dat} & $o$-\Htwo{} and $p$-\Htwo\,+\,\Htwo{}, $v=0$, 1--$10^4$\,K, \citet{Wan2018} \\
\src{dataPDR/\allowbreak{}h2/\allowbreak{}coll\_rates\_Hp.dat} & \Hp\,+\,\Htwo{} including proton exchange, $v=0$, $J\le9$, 3--6000\,K, \citet{Gerlich1990} \\
\src{dataPDR/\allowbreak{}h2/\allowbreak{}energy\_\{state\}.dat}, \src{transprob\_\{state\}.dat}, \src{dissprob\_\{state\}.dat} & excited electronic states $B$, $C^\pm$, $B'$, $D^\pm$: levels, lines to $X$, spontaneous radiative dissociation rates and kinetic energies, \citet{Abgrall2000}, whose line transition probabilities rest on \citet{Abgrall1993} and \citet{Abgrall1994}; the $J=0$ energies of the $\Pi$ components that the files carry as comments are extrapolations by Cloudy's authors and are not read \\
\src{dataPDR/\allowbreak{}h2/\allowbreak{}energy\_dissoc.dat} & dissociation energies of the electronic states, \citet{Sharp1970} per the header (which dates it 1971), with values revised in Cloudy in 2002; only the value for $X$, 36\,118.11\,cm$^{-1}$, is used \\
\src{dataPDR/\allowbreak{}h2/\allowbreak{}README.md} & provenance and Cloudy's licence notice (zlib) \\
MOCASSIN atomic data (\src{data/\allowbreak{}ph1.dat}, \src{ph2.dat}, \src{radrec.dat}) & photoionisation cross-section fits \citep{Verner1996,VernerYakovlev1995} and radiative recombination fits \citep[and the sources listed at \src{update\_mod.f90:\allowbreak{}4101--4115}]{VernerFerland1996}, used for the linked atoms' rates; the \Htwo{} photoionisation cross section is a fit in the code \citep{Yan1998,Yan2001} \\
\bottomrule
\end{longtable}


\clearpage
\bibliographystyle{plainnat}
\bibliography{references}

\end{document}
